
\chapter{三角}


%\section{正弦、余弦、正切、余切}


\section{锐角的正弦、余弦、正切、余切}
\zsd 锐角的正弦、余弦、正切、余切

A组：
\begin{enumerate}
\item 填空\\
\begin{minipage}{\textwidth}
\centering
\includegraphics[scale=0.9]{./bodyxb/src/main42xbb2ch06ruijiandesanjiaobi01t1.png}
\end{minipage}
\begin{enumerate}[(1)]
\item $\sin{A}=$\blkda{$\frac{a}{c}$}, $\cos{A}=$\blkda{$\frac{b}{c}$}, $\tan{A}=$\blkda{$\frac{a}{b}$}, $\cot{A}=$\blkda{$\frac{b}{a}$}. 
\item $\sin^2A+\cos^2A=$\blkda{1}. 
\item $\sin\left(\ang{90}-A\right)=$\blkda{$cosA$}, $\cos\left(\ang{90}-A\right)=$\blkda{$sinA$}, \\
$\tan\left(\ang{90}-A\right)=$\blkda{$cotA$}, $\cot\left(\ang{90}-A\right)=$\blkda{$tanA$}. 
\item 
\begin{minipage}[t]{\textwidth}
\begin{tblr}{hlines,vlines,column{1}={2cm,c},column{2,3,4,5}={2.5cm,c},row{1}={blue!10}}
角度$\alpha$ & $\sin\alpha$ & $\cos\alpha$ & $\tan\alpha$ & $\cot\alpha$\\
$\ang{30}$ & \blkda{$\frac{1}{2}$} & \blkda{$\frac{\sqrt{3}}{2}$} &\blkda{$\frac{\sqrt{3}}{3}$} & \blkda{$\sqrt{3}$} \\
$\ang{45}$ & \blkda{$\frac{\sqrt{2}}{2}$} & \blkda{$\frac{\sqrt{2}}{2}$} &\blkda{$1$} & \blkda{$1$} \\
$\ang{60}$ & \blkda{$\frac{\sqrt{3}}{2}$} & \blkda{$\frac{1}{2}$} & \blkda{$\sqrt{3}$} & \blkda{$\frac{\sqrt{3}}{3}$} \\
\end{tblr}
\end{minipage}
\end{enumerate}

\item  如图, 在 $x$ 轴的上方, 直角 $\angle B O A$ 绕原点 $O$ 按顺时针方向旋转. 若 $\angle \mathrm{BOA}$ 的两边分别与函数 $y=-\frac{1}{x}, y=\frac{2}{x}$ 的图象交于$B 、 A$ 两点, 则 $\angle O A B$ 大小的变化趋势为\dotfill{(\kh{D})}
\xx{逐渐变小}{逐渐变大}{时大时小}{保持不变}
\begin{minipage}{\textwidth}
\centering
\includegraphics[scale=0.9]{./bodyxb/src/31247891329857981275.jpg}
\end{minipage}
\begin{jieda}
    过$A, B$做$x$轴垂线，设垂足为$A', B'$，根据三角形相似可知$|OB'| \cdot |AA'| = \sqrt{2}$，进一步可以推导出$\dfrac{|OA|}{|OB|} = \sqrt{2}$，即$\angle \mathrm{BOA}$的正切值固定，因此选D.
\end{jieda}


\item 试判断下列命题是否正确:
\begin{enumerate}
    \item 小于 $90^{\circ}$ 的角是锐角;\dotfill{( )}
    \item 第一象限角小于第二象限角；\dotfill{( )}
    \item 终边相同的角一定相等；\dotfill{( )}
    \item 相等的角终边一定相同；\dotfill{( )}
    \item 若 $90^{\circ}<a<180^{\circ}$, 则 $a$ 是第二象限角.\dotfill{( )}
\end{enumerate}
\begin{jieda}
\begin{enumerate}[(1)]
\item $-\ang{45}$角不是锐角, 故假. 
\item $\ang{30}>\ang{120}-\ang{360}=-\ang{240}$, 故假. 
\item $\alpha=\ang{30}$, $\beta=\ang{30}+\ang{360}$, 故假. 
\item 真. 如果有破绽的话, 破绽在于始边是否重合. 
\item 真. 
\end{enumerate}
\end{jieda}



\item $-2021^{\circ}$ 是第\blkda{II}象限角.
\begin{jieda}
$-2021=-5\cdot360-221=-6\cdot 360+139$, 故II. 
\end{jieda}


\item 角 $a$ 是第四象限角, 则角 $180^{\circ}-a$ 是第\blkda{三}象限角.


\item 分别写出第一、二、三、四象限角的集合. 
\begin{jieda}
一: $\{\theta\mid\theta=k\cdot\ang{360}+\alpha,\ang{0}\alpha<\ang{90},k\zz\}$; \\
二: $\{\theta\mid\theta=k\cdot\ang{360}+\alpha,\ang{90}\alpha<\ang{180},k\zz\}$; \\
三: $\{\theta\mid\theta=k\cdot\ang{360}+\alpha,\ang{180}\alpha<\ang{270},k\zz\}$; \\
四: $\{\theta\mid\theta=k\cdot\ang{360}+\alpha,\ang{270}\alpha<\ang{360},k\zz\}$. 
\end{jieda}


\end{enumerate}

B组：
\begin{enumerate}
\item $\dfrac{\pi}{5}$=\blkda{36}度; $-72^{\circ}$=\blkda{$-\dfrac{2\pi}{5}$}弧度.


\item 若时钟的分针转了2小时15分, 则所转的角度是\blkda{$-\ang{810}$}. 
\begin{jieda}
$s=vt=360\cdot 2.25=810$, 就是这个东西. 
\end{jieda}


\item 角的顶点在坐标系的原点, 始边与$x$轴的正半轴重合, 那么终边在下列位置的角的集合分别是: 
\begin{enumerate}[(1)]
\item $x$轴正半轴: \blkda[5]{$\{\theta\mid \theta=k\cdot\ang{360},k\zz\}$}; 
\item $x$轴: \blkda[5]{$\{\theta\mid \theta=k\cdot\ang{180},k\zz\}$}; 
\item $y$轴负半轴: \blkda[5]{$\{\theta\mid \theta=k\cdot\ang{360}-\ang{90},k\zz\}$}; 
\item $y$轴: \blkda[5]{$\{\theta\mid \theta=k\cdot\ang{180}+\ang{90},k\zz\}$}; 
% \item 坐标轴: \blkda{}; 
\item 直线$y=-x$: \blkda[5]{$\{\theta\mid \theta=k\cdot\ang{180}+\ang{135},k\zz\}$}; 
\item 直线$y=\pm{x}$: \blkda[5]{$\{\theta\mid \theta=k\cdot\ang{90}+\ang{45},k\zz\}$}. 
\end{enumerate}
\begin{jieda}
以(6)为例: 每逆时针旋转$\ang{90}$, 恰好到达一条射线. 于是从角的角度看, 遍历了所有的角. 
相当于有了一个起始值$\ang{45}$, 每次加上或减去$\ang{90}$, 就有了所有角. 
\end{jieda}



\item 已知$\alpha$与$\beta$都是锐角, $\alpha+\beta$的终边与$-280^{\circ}$的终边相同; 
$\alpha-\beta$的终边与$-670^{\circ}$的终边相同, 则$\alpha$=\blkda{$\ang{65}$}, $\beta$=\blkda{$\ang{15}$}. 
\begin{jieda}
$\begin{cases}
\alpha+\beta=m\cdot\ang{360}-\ang{280}\\
\alpha-\beta=n\cdot\ang{360}-\ang{670}
\end{cases}$, 
故$\alpha=\left(m+n\right)\cdot\ang{180}-\ang{475}$, $\beta=\left(m-n\right)\cdot\ang{180}+\ang{195}$. \\
又$\alpha,\beta$是锐角, 故$\alpha=\ang{65}$, $\beta=\ang{15}$. 
\end{jieda}


\item 已知圆周上有一点$P$(初始位置在$x$轴正半轴上)依逆时针方向作匀速圆周运动. 已知点$P$每分钟转过$\theta$角($\ang{0}<\theta\leq\ang{180})$, 
2分钟到达第三象限, 15分钟回到原来的位置, 则角$\theta$等于\blkda[2]{$\ang{96}$或$\ang{120}$}. 
\begin{jieda}
$\begin{cases}
\ang{0}<\theta<\ang{180}\\
\ang{180}<2\theta<\ang{270}\\
15\theta=k\cdot\ang{360}
\end{cases}$, 
将第三式代入前两式得到$k$的范围, 结合$k\zz$得$k=4,5$, 故$\theta=\ang{96}$或$\ang{120}$. 
\end{jieda}


\item 若$\alpha$与$\beta$是终边相同的角, 则$\alpha-\beta$的终边一定在\dotfill{(\kh{A})}
\xx{$x$轴的非负半轴上}
{$x$轴的非正半轴上}
{$y$轴的非正半轴上}
{$y$轴的非负半轴上}



\item 设$k\zz$, 终边相同的一组角是\dotfill{(\kh{C})}
\xx{$k\cdot 90^{\circ}$与$k\cdot 180^{\circ}+90^{\circ}$}
{$k\cdot 180^{\circ}+60^{\circ}$与$k\cdot 60^{\circ}$}
{$\left(2k-1\right)\cdot 180^{\circ}$与$\left(4k\pm 1\right)\cdot 180^{\circ}$}
{$k\cdot 180^{\circ}+30^{\circ}$与$k\cdot 360^{\circ}-150^{\circ}$}



\item 在与2021弧度角的终边相同的角中, 绝对值最小的角是\blkda[3]{$2021-644\pi$}. 
\begin{jieda}
$\abs{2021+2k\pi}_{\min}=\abs{2021-322\cdot2\pi}$. 
\end{jieda}



\item 若角$\alpha$和角$\beta$的和是$1$弧度, 差为$1^{\circ}$, 则$\alpha$和$\beta$的弧度数分别为\blkda{$\frac{1}{2}\pm\frac{\pi}{360}$}. 

\item 若中心角为$1^{\circ}$的扇形, 它的弧长为$2\pi$, 则该扇形所在圆的半径为\blkda{360}. 
\begin{jieda}
$r=\dfrac{l}{\alpha}=\dfrac{2\pi}{\frac{\pi}{180}}=360$. 
\end{jieda}


\item 若扇形的圆心角是$120^{\circ}$, 则该扇形面积与其内切圆面积的比值是\blkda{$\dfrac{7+4\sqrt{3}}{9}$}. 
\begin{jieda}
相切找切点, 于是结合简单的几何知识易得$R=r+\dfrac{2r}{\sqrt{3}}$. 
\end{jieda}


\item 若一条弦长等于半径的$\sqrt{2}$倍, 则此弦所对圆心角等于\blkda{$\dfrac{\pi}{2}$}弧度.
\begin{jieda}
$\sin\frac{\theta}{2}=\frac{\sqrt{2}}{2}$. 
\end{jieda}

\item 如果$1$弧度的圆心角所对的弦长为$2$, 那么这个圆心角所对的弧长是\blkda{$\dfrac{1}{\sin\frac{1}{2}}$}.
\begin{jieda}
``弧''与``弦''; 弦取中点. 
\end{jieda}



\item 若扇形的周长是$16$, 圆心角是$2$弧度, 则扇形面积是\blkda{$16$}.


\item
\begin{enumerate}[(1)]
\item 若角$\alpha$的终边过点$\left(2^x,\log_{\frac{1}{2}}2\right)$, $x\zr$, 则$\dfrac{\alpha}{2}$终边在第\blkda{II、IV}象限. 

\item 设$\alpha$是第二象限角, 则$\dfrac{\alpha}{3}$的终边不在\blkda{第三象限}.

\end{enumerate}

\begin{jieda}

先看图: \\
\begin{center}
    \includegraphics[scale=0.8]{./bodyxb/src/main42chap05renyijiaodesanjiaobi01.pdf}
\end{center}

该图的信息有点多: 

\begin{itemize}
\item 由$\alpha$所在的象限, 可得$\dfrac{\alpha}{2}$所在象限, 且只能是象限中的一部分. 
\item 由$\alpha$所在的象限, 可得$2\alpha$所在象限. 
\item 可推广到$\dfrac{\alpha}{n}\left(n\zn\right)$的情形. 
\end{itemize}

\begin{enumerate}[(1)]
\item $2^x>0, \log_{\frac{1}{2}}2=-1$, 故$\alpha$是第四象限角, 于是$\dfrac{\alpha}{2}$是第II,IV象限角. 

\item 每个象限等分为3份, 然后逆时针依次标上1234(周期). 

\end{enumerate}

\end{jieda}



\item 一扇形周长是\SI{32}{cm}, 扇形的圆心角为多少弧度时, 这个扇形的面积最大? 最大面积是多少? 
\begin{jieda}
$l+2r=32$, $S=\frac{1}{2}lr=\frac{1}{4}\cdot l\cdot2r\leq\frac{1}{4}\left(\frac{l+2r}{2}\right)^2=64$. \\
$\alpha=\frac{l}{r}=\SI{2}{rad}$, $S_{\min}=\SI{64}{cm^2}$. 
\end{jieda}


\item 已知扇形的圆心角为 $\frac{\pi}{3}$, 半径为 $R$ , 圆 $C$ 与扇形的两条半径及扇形的弧都相切, 求圆 $C $中圆心角为 $\frac{\pi}{3}$ 的扇形与原扇形的面积之比.
\begin{jieda}
    根据几何关系，半径之比为$\dfrac{1}{3}$，因此面积之比$\dfrac{1}{9}$
\end{jieda}

\end{enumerate}

C组：
\begin{enumerate}
    \item 求$\dfrac{1-2log_65 \cdot log_{10}3 \cdot log_{15}2}{log_65 \cdot log_{10}3 + log_{10}3 \cdot log_{15}2 + log_{15}2 \cdot log_65}$的值.
    \begin{jieda}
        考虑换底，$log_65 = \dfrac{log_25}{log_23+1}, log_{10}3 = \dfrac{log_23}{log_25+1}, log_{15}2 = \dfrac{1}{log_23 + log_25}$\\
        故设$a = log_23, b = log_25$\\
        原式等于$\dfrac{1 - 2 \cdot \dfrac{b}{a+1} \cdot \dfrac{a}{b+1} \cdot \dfrac{1}{a+b}}{\dfrac{b}{a+1}\cdot \dfrac{a}{b+1} + \dfrac{a}{b+1} \cdot \dfrac{1}{a+b} + \dfrac{1}{a+b} \cdot \dfrac{b}{a+1}} = \dfrac{\left(a+1\right)\left(b+1\right)\left(a+b\right) - 2ab}{ab\left(a+b\right) + a\left(a+1\right) + b\left(b+1\right)} = 1$
    \end{jieda}
    \item 已知定义在$R^{+}$上的函数$f\left(x\right) = \left \{\begin{array}{ll}
        |log_3x - 1|, & 0 < x \leq 9 \\
        4-\sqrt{x}, & x > 9
    \end{array} \right.$，设$a, b, c$是三个互不相同的实数，满足$f\left(a\right) = f\left(b\right) = f\left(c\right)$，求$abc$的取值范围.
    \begin{jieda}
        不妨设$a < b < c$，易知$0 < a < 3 < b < 9< c$，因此有$1 - log_3a = log_3b - 1$，即$ab = 9$. \\
        故$f\left(a\right) \in \left(0, 1\right)$，因此$c \in \left(9, 16\right)$，因此有$abc \in \left(81, 144\right)$
    \end{jieda}
    \item 若关于$x$的方程$\dfrac{|x|}{x+6} = kx^2$有四个不同的实数解，求实数$k$的取值范围.
    \begin{jieda}
        易知必有一根为0 \\
        参变分离有$\dfrac{1}{k} = \dfrac{\left(x+6\right)x^2}{|x|} = \left \{ \begin{array}{ll}
            \dfrac{\left(x+6\right)x^2}{x}, & x > 0 \\
            -\dfrac{\left(x+6\right)x^2}{x},  & x < 0
        \end{array}\right.$
        由于原方程必有一根为0，故可以转化为方程$\dfrac{1}{k} = \left \{ \begin{array}{ll}
            \left(x+6\right)x, & x > 0 \\
            -\left(x+6\right)x,  & x < 0
        \end{array}\right.$有三个不等实根.因此$\dfrac{1}{k} \in \left(0, 9\right)$，即$k \in \left(\dfrac{1}{9}, +\infty\right)$
    \end{jieda}
\end{enumerate}

\section{任意角的正弦、余弦、正切、余切(1)}
\zsd 任意角的正弦、余弦、正切、余切及其符号. 同角三角比关系(求值)

A组：
\begin{enumerate}
\item 已知角$\alpha$的终边经过下列各点, 求角$\alpha$的正弦、余弦、正切、余切. 
\begin{enumerate}[(1)]
\item $A\left(-8,-6\right)$,     则$\sin\alpha$=\blkda{$-\frac{3}{5}$}, $\alpha$=\blkda{$-\frac{4}{5}$}, $\tan\alpha$=\blkda{$\frac{3}{4}$}, $\cot\alpha$=\blkda{$\frac{4}{3}$}; 
\item $B\left(\sqrt3,-1\right)$, 则$\sin\alpha$=\blkda{$-\frac{1}{2}$}, $\cos\alpha$=\blkda{$\frac{\sqrt{3}}{2}$}, $\tan\alpha$=\blkda{$-\frac{\sqrt{3}}{3}$}, $\cot\alpha$=\blkda{$-\sqrt{3}$}. 
\end{enumerate}


\item 下列各式为正号的是\dotfill{(\kh{C})}
\xx{$\cos2-\sin2$}
{$\cos2\cdot\sin2$}
{$\tan2\cdot\sec2$}
{$\sin2\cdot\tan2$}
\begin{jieda}
$2\in\left(\frac{\pi}{2},\pi\right)$, 故$\sin2>0$, $\cos2<0$, $\tan2<0$, $\cot2<0$, $\sec2<0$, $\csc2>0$. 
\end{jieda}



\item 填表:

\begin{tabular}{|c|c|c|c|}\hline
$\alpha$ & $\sin\alpha$ & $\cos\alpha$ & $\tan\alpha$\\\hline
$\frac{2\pi}{3}$ & \blkda{$\frac{\sqrt{3}}{2}$} & \blkda{$-\frac{1}{2}$} & \blkda{$-\sqrt{3}$} \\\hline
$\frac{5\pi}{4}$ & \blkda{$-\frac{\sqrt{2}}{2}$} & \blkda{$-\frac{\sqrt{2}}{2}$} & \blkda{1} \\\hline
$\frac{5\pi}{6}$ & \blkda{$\frac{1}{2}$} & \blkda{$-\frac{\sqrt{3}}{2}$} & \blkda{$-\frac{\sqrt{3}}{3}$}\\\hline
\end{tabular}


\item 填表:

\begin{tabular}{|c|c|c|c|c|c|}\hline
$\alpha$ & 点$P$的坐标 & ~~$\sin\alpha$~~ & ~~$\cos\alpha$~~ & ~~$\tan\alpha$~~ & ~~$\cot\alpha$~~ \\\hline
$\alpha=2k\pi\left(k\zz\right)$ & $\left(1,0\right)$ & 0 & 1 & 0 & 不存在\\\hline
\blkda{$\alpha=2k\pi+\frac{\pi}{2}\left(k\zz\right)$} & $\left(0,1\right)$ & \blkda{1} & \blkda{0} & \blkda{不存在} &  \blkda{0} \\\hline
\blkda{$\alpha=2k\pi+\pi\left(k\zz\right)$} & $\left(-1,0\right)$  & \blkda{0} & \blkda{-1} & \blkda{0} & \blkda{不存在} \\\hline
\blkda{$\alpha=2k\pi-\frac{\pi}{2}\left(k\zz\right)$}& $\left(0,-1\right)$ & \blkda{-1} & \blkda{0} & \blkda{不存在}  & \blkda{0}\\\hline
\end{tabular}



\item 已知$\sin\alpha=\dfrac{4}{5}$, $\alpha\in\left(0,\pi\right)$, 则$\tan\alpha$等于\blkda{$\pm\dfrac{4}{3}$}. 
\begin{jieda}
三角比的值分为两部分: 符号与绝对值. 结合定义, 直接赋值: $y=4,r=5$, 于是$x=\pm3$. 
\end{jieda}


\item 已知角 $\theta$ 满足 $\sin \theta<0$ 且 $\cos \theta>0$, 则角 $\theta$ 是第\blkda{四}象限的角.

\item 已知角 $a$ 的终边过点 $P\left(2 a, 3 a\right)\left(a<0\right)$, 则角 $a$ 的余弦值为\blkda{$-\frac{2\sqrt{13}}{13}$}.


\item 已知角 $\theta$ 满足 $\frac{\sin \theta}{\cos \theta}>0$, 则角 $\theta$ 是第\blkda{一或三}象限的角.


\item 设 $x \in \mathrm{R}$, 且 $x \neq k \pi+\frac{\pi}{2}\left(k \in Z\right)$, 则 $\frac{\sin x}{|\sin x|}+\frac{|\cos x|}{\cos x}+\frac{\tan x}{|\tan x|}$ 的取值集合为\blkda{$\{-1, 3\}$}


\end{enumerate}

B组：
\begin{enumerate}
\item 计算: $\cos\dfrac{\pi}{3}-\tan\dfrac{\pi}{4}+\dfrac{3}{4}\tan^2\dfrac{7\pi}{6}-\sin\dfrac{5\pi}{6}+\cos^2\dfrac{\pi}{6}+\sin\dfrac{3\pi}{2}$=\blkda{$-1$}. 

\item 若$\sin\theta=\dfrac{m-3}{m+5}$, $\cos\theta=\dfrac{4-2m}{m+5}$都有意义, 则$m$=\blkda{0或8}. 
\begin{jieda}
$\sin^2\theta+\cos^2\theta=1$. 
\end{jieda}


\item 已知 $\alpha$ 为第二象限的角, 其终边上有一点 $P\left(x, \sqrt{5}\right)$, 且 $\cos \alpha=\frac{\sqrt{2}}{4} x$, 求 $\tan \alpha =$ \blkda{$-\frac{\sqrt{15}}{3}$}.
\begin{jieda}
    $\dfrac{1}{r} = \dfrac{\sqrt{2}}{4}$，因此$sin\alpha = \dfrac{\sqrt{5}}{r} = \sqrt{5} \cdot \dfrac{\sqrt{2}}{4} = \dfrac{\sqrt{10}}{4}$，故$cos\alpha = -\sqrt{1 - \dfrac{10}{16}} = -\dfrac{\sqrt{6}}{4}$，因此$tan\alpha = \dfrac{sin\alpha}{cos\alpha} = -\dfrac{\sqrt{10}}{\sqrt{6}} = -\dfrac{\sqrt{15}}{3}$
\end{jieda}


\item 已知角$\theta$的终边上有一点$P\left(4,y\right)$, 且$\sin\theta=-\dfrac{2\sqrt{5}}{5}$, 则$y$=\blkda{$-8$}. 

\item 已知角$\alpha$的终边与射线$5x+12y=0\left(x\leq 0\right)$的图像重合, $\cos\alpha+\cot\alpha-\tan\alpha$的值为\blkda{$-\dfrac{2267}{780}$}. 
\begin{jieda}
$5x+12y=0\left(x\leq 0\right)$上取点$\left(-12,5\right)$, 
则$E=\dfrac{-12}{13}+\dfrac{-12}{5}-\dfrac{5}{-12}=-\dfrac{2267}{780}$. 
\end{jieda}


\item 已知集合$A=\left\{y\mid y=\dfrac{\abs{\sin{x}}}{\sin{x}}+\dfrac{\cos{x}}{\abs{\cos{x}}}+\dfrac{\abs{\tan{x}}}{\tan{x}}\right\}$, 
用列举法表示集合$A$是\blkda{$\{-1,3\}$}. 
\begin{jieda}
四个象限讨论一下即可. 
\end{jieda}


\item 若角$\alpha$是第三象限角, 则
\ding{172} $\sin\alpha+\cos\alpha<0$; 
\ding{173} $\tan\alpha-\sin\alpha>0$; 
\ding{174} $\cot\alpha\cdot\csc\alpha<0$; 
\ding{175} $\sin\alpha\cdot\sec\alpha>0$. 
其中正确的是\blkda{\ding{172}\ding{173}\ding{174}\ding{175}}.  

\item 已知角$\alpha$的终边过点$\left(3a-9,a+2\right)$, 且$\cos\alpha\leq 0$, $\sin\alpha>0$, 则$a$的取值范围是\blkda{$\left(-2,3\right]$}. 

\item 在$\triangle{ABC}$中, 有$\cos{A}\cdot\tan{B}\cdot\cot{C}<0$, 则这三角形是\blkda{钝角}三角形.
\begin{jieda}
首先没有直角, 否则要么积为零, 要么积不存在. \\
其次$\cos,\tan,\cot$在$\left(0,\frac{\pi}{2}\right)$上都是正的, 在$\left(\frac{\pi}{2},\pi\right)$上都是负的, \\
故结合条件可得$\cos{A},\tan{B},\cot{C}$中一负两正, 于是钝角三角形. 
\end{jieda}



\item 设$y=\sin\alpha\cdot\cos\alpha\cdot\cot\alpha$, 且$\alpha$是象限角, 则$y$的符号为\dotfill{(\kh{A})}
\xx{恒正}{恒负}{可能为0}{不确定}



\item 若角$\alpha$的终边落在直线$x+y=0$上, 则$\dfrac{\sin\alpha}{\sqrt{1-\sin^2\alpha}}+\dfrac{\sqrt{1-\cos^2\alpha}}{\cos\alpha}$的值等于\blkda{$0$}.
\begin{jieda}
考虑一下符号基本上就可以把答案选出来了. 
\end{jieda}


\item 已知实数$\alpha,\beta$满足$\abs{\cos\alpha-\cos\beta}=\abs{\cos\alpha}+\abs{\cos\beta}$, 且$\alpha\in\left(\frac{\pi}{2},\pi\right)$, 
则化简$\sqrt{\left(\cos\alpha-\cos\beta\right)^2}$的结果是\dotfill{(\kh{C})}
\xx{$\cos\alpha-\cos\beta$}
{$\abs{\cos\alpha}-\abs{\cos\beta}$}
{$\cos\beta-\cos\alpha$}
{$\abs{\cos\alpha}+\abs{\cos\beta}$}
\begin{jieda}
$\abs{\cos\alpha-\cos\beta}=\abs{\cos\alpha}+\abs{\cos\beta}\iff\cos\alpha\cos\beta\leq 0$. \\
又$\cos\alpha<0$, 故$\cos\beta\geq 0$. \\
故$E=\abs{\cos\alpha-\cos\beta}=\cos\beta-\cos\alpha$. 
\end{jieda}


\item 已知角$\alpha$的终边上一点$P$与点$A\left(-3,2\right)$关于$y$轴对称, 角$\beta$的终边上一点$Q$与点$A$关于原点对称, 求$2\sin\alpha+3\sin\beta$的值. 
\begin{jieda}
$A\left(-3,2\right)\xleftrightarrow{y\text{轴}}P\left(3,2\right)$, $A\left(-3,2\right)\xleftrightarrow{\text{原点}}Q\left(3,-2\right)$. \\
故$2\sin\alpha+3\sin\beta=2\cdot\dfrac{2}{\sqrt{13}}+3\cdot\dfrac{-2}{\sqrt{13}}=-\dfrac{2\sqrt{13}}{13}$. 
\end{jieda}


\item 已知角$\theta$终边上一点$P\left(-\sqrt{3},m\right)$, 且$\sin\theta=\dfrac{\sqrt{2}m}{4}$, 求$\cos\theta$和$\tan\theta$的值. 
\begin{jieda}
$\dfrac{m}{\sqrt{3+m^2}}=\dfrac{\sqrt{2}m}{4}$, 则$m=0,\pm\sqrt{5}$. 
\begin{itemize}
\item $m=0$时, $\cos\theta=-1$, $\tan\theta=0$; 
\item $m=\sqrt{5}$时, $\cos\theta=-\dfrac{\sqrt{6}}{4}$, $\tan\theta=-\dfrac{\sqrt{15}}{3}$; 
\item $m=-\sqrt{5}$时, $\cos\theta=-\dfrac{\sqrt{6}}{4}$, $\tan\theta=\dfrac{\sqrt{15}}{3}$. 
\end{itemize}
\end{jieda}


\item 已知 $\sin \alpha+\cos \alpha=\frac{1}{5}\left(0<\alpha<\pi\right)$, 求: \\
\begin{enumerate*}
    \item  $\sin \alpha-\cos \alpha$;
    \item  $\sin ^3 \alpha+\cos ^3 \alpha$;
    \item $\sin ^4 \alpha+\cos ^4 \alpha$.
\end{enumerate*}
\begin{jieda}
    \begin{enumerate}
        \item $\left(sin\alpha+cos\alpha\right)^2 = 1 + 2sin\alpha cos\alpha = \dfrac{1}{25}$，因此$sin\alpha cos\alpha = -\dfrac{12}{25}$ \\
        因此有$\left(sin\alpha - cos\alpha\right)^2 = 1 - 2sin\alpha cos\alpha = \dfrac{49}{25}$ \\
        故$sin\alpha - cos\alpha = \pm \dfrac{7}{5}$
        \\
        根据$\alpha \in \left(0, \pi\right)$，故$sin \alpha - cos\alpha > -1$，因此$sin\alpha - cos\alpha = \dfrac{7}{5}$ \\
        \ps 可以根据$sin\alpha + cos\alpha$在不同象限的值域直接判断出来$\alpha$是第二象限角.
        \item $sin^3 \alpha + cos^3 \alpha = \left(sin\alpha + cos\alpha\right)\left(sin^2 \alpha + cos^2 \alpha - sin\alpha cos\alpha\right) = \left(\dfrac{1}{5}\right)\left(1 + \dfrac{12}{25}\right) = \dfrac{37}{125}$ \\
        \ps 也可以在第一问的结果上构建方程组$\left \{ \begin{array}{l}
            sin\alpha + cos\alpha = \frac{1}{5} \\
            sin\alpha - cos\alpha = \frac{7}{5} 
        \end{array}\right.$解得$\left \{ \begin{array}{l}
            sin\alpha = \frac{4}{5} \\
            cos\alpha = -\frac{3}{5} 
        \end{array}\right.$后带入计算(下同)
        \item $\sin ^4 \alpha+\cos ^4 \alpha = \left(\sin ^2 \alpha+\cos ^2 \alpha\right)^2 - 2\left(sin \alpha cos\alpha\right)^2 = 1 - 2 \times \left(-\dfrac{12}{25}\right)^2 = \dfrac{337}{625}$
    \end{enumerate}
\end{jieda}
 
\end{enumerate}

C组：
\begin{enumerate}
    \item 求函数$f\left(x\right) = \sqrt{x^4 + 3x^2 -6x + 10} - \sqrt{x^4 - 3x^2 + 2x + 5}$的最大值.
    \begin{jieda}
        $f\left(x\right) = \sqrt{\left(x^2+1\right)^2 + \left(x-3\right)^2} - \sqrt{\left(x^2-2\right)^2 + \left(x+1\right)^2}$ \\
        其几何意义是点$P\left(x, x^2\right)$到$A\left(3, -1\right)$和$B\left(-1, 2\right)$的距离之差$|PA| - |PB|$.由三角形两边之和大于第三边可知$|PA| - |PB| \leq |AB|$，而等号可以取到.\\
        因此$f\left(x\right) \leq 5$，$x = \dfrac{\sqrt{89}-3}{8}$时取最大值.
    \end{jieda}
    \item 求函数$y = 2x + \sqrt{4x^2-8x+3}$的最小值.
    \begin{jieda}
        定义域先行，$4x^2-8x+3 \geq 0$，因此定义域为$\left(-\infty, \dfrac{1}{2}] \cup [\dfrac{3}{2}, +\infty\right)$ ，下面分析单调性
        \begin{enumerate}
            \item 当$x \geq \dfrac{3}{2}$时，容易判断函数单调递增
            \item 当$x \leq \dfrac{1}{2}$时，$f\left(x\right) = \left(2x-2\right) + \sqrt{\left(2x-2\right)^2-1}+2 = 2 + \dfrac{1}{\left(2x-2\right) - \sqrt{\left(2x-2\right)^2-1}}$，可以判断函数单调递减.
        \end{enumerate}
        因此最值应在端点处取得，$f\left(\dfrac{1}{2}\right) = 1, f\left(\dfrac{3}{2}\right) = 3$，因此函数最小值为$1$
    \end{jieda}
    \item 已知实数$a > b > 0$，求函数$f\left(x\right) = \dfrac{x}{\sqrt{a-x^2}-\sqrt{b-x^2}}$的最大值.
    \begin{jieda}
        $f\left(x\right) = \dfrac{x}{\sqrt{a-x^2}-\sqrt{b-x^2}} = \dfrac{x\left(\sqrt{a-x^2}+\sqrt{b-x^2}\right)}{a-b} \\
        = \dfrac{1}{a-b}\sqrt{x^2\left(a+b\right)-2x^4+2x^2\sqrt{\left(a-x^2\right)\left(b-x^2\right)}} \\
        = \dfrac{1}{a-b}\sqrt{ab - [x^2 - \sqrt{\left(a-x^2\right)\left(b-x^2\right)}]^2} \\
        \leq \dfrac{\sqrt{ab}}{a-b}$
        当且仅当$x = \sqrt{\dfrac{ab}{a+b}}$时取等号.
    \end{jieda}
\end{enumerate}



\section{任意角的正弦、余弦、正切、余切(2)}


\begin{enumerate}
\item 已知$\tan\alpha=\dfrac{2ab}{a^2-b^2}$, 其中$a>b>0$, $\alpha\in\left(0,\frac{\pi}{2}\right)$, 则$\sin\alpha$=\blkda{$\dfrac{2ab}{a^2+b^2}$}.
\begin{jieda}
$\sin\alpha=\dfrac{2ab}{\sqrt{\left(2ab\right)^2+\left(a^2-b^2\right)^2}}$. 
\end{jieda}


\item 已知$\dfrac{\cos\alpha}{1+\sin\alpha}=-\dfrac{1}{2}$, 则$\dfrac{\cos\alpha}{\sin\alpha-1}$=\blkda{2}. 
\begin{jieda}
结合$\cos^2\alpha+\sin^2\alpha=1$解得$\left(\cos\alpha,\sin\alpha\right)=\left(-\frac{4}{5},\frac{3}{5}\right)$(另一组增根已舍去), 
代入所求即得. 
\end{jieda}



\item 若$\cot\left(\sin\theta\right)\cdot\tan\left(\cos\theta\right)>0$, 则$\theta$是第\blkda{I、III}象限角. 
\begin{jieda}
四个象限讨论一下即可. $\cos\theta,\sin\theta\in[-1,1]\subseteq\left(-\frac{\pi}{2},\frac{\pi}{2}\right)$. 
\end{jieda}



\item 若$\theta$是锐角, $\sin\theta-\cos\theta=\dfrac{1}{2}$, 则$\sin^3\theta-\cos^3\theta$=\blkda{$\dfrac{11}{16}$}. 
\begin{jieda}
围绕$\left(\sin\theta\pm\cos\theta\right)^2=1\pm2\sin\theta\cos\theta$展开. 
\begin{align*}
\sin^3\theta-\cos^3\theta & =\left(\sin\theta-\cos\theta\right)\left(\sin^2\theta+\sin\theta\cos\theta+\cos^2\theta\right)\\
& = \left(\sin\theta-\cos\theta\right)\left(1+\dfrac{1-\left(\sin\theta-\cos\theta\right)^2}{2}\right)\\
& = \dfrac{1}{2}\cdot\dfrac{3-\frac{1}{4}}{2}\\
& = \dfrac{11}{16}.
\end{align*}
\end{jieda}



\item 设$\cot{x}=2$, 则$\dfrac{2\cos{x}-4\sin{x}}{5\cos{x}+3\sin{x}}$=\blkda{0}, 
$3\sin^2x-4\cos^2x$=\blkda{$-\dfrac{13}{5}$}. 
\begin{jieda}
$\dfrac{2\cos{x}-4\sin{x}}{5\cos{x}+3\sin{x}}=\dfrac{2\cot{x}-4}{5\cot{x}+3}=0$. \\
$3\sin^2x-4\cos^2x=\dfrac{3\sin^2x-4\cos^2x}{\cos^2x+\sin^2x}=\dfrac{3-4\cot^2x}{\cot^2x+1}=-\dfrac{13}{5}$. 
\end{jieda}

\item 用列举法写出集合$A=\left\{y\mid y=\dfrac{1}{\cos\alpha\sqrt{1+\tan^2\alpha}}
+\dfrac{2\tan\alpha}{\sqrt{\sec^2\alpha-1}}\right\}$=\blkda{$\{\pm1,\pm3\}$}.
\begin{jieda}
$y=\dfrac{\abs{\cos\alpha}}{\cos\alpha}+\dfrac{2\tan\alpha}{\abs{\tan\alpha}}$. 
\end{jieda}


\item 使函数$y=\sqrt{\sin{x}}+\sqrt{16-x^2}$有意义的$x$的取值范围是\blkda[3]{$[-4,-\pi]\cup[0,\pi]$}; 
使函数$y=\lg\sin\left(\cos{x}\right)$有意义的$x$的取值范围是\blkda[5]{$\left(-\dfrac{\pi}{2}+2k\pi,\dfrac{\pi}{2}+2k\pi\right)\left(k\zz\right)$}. 
\begin{jieda}
$\sin{x}\geq 0\land 16-x^2\geq 0$, 故$x\in[-4,-\pi]\cup[0,\pi]$. \\
$\sin\left(\cos{x}\right)>0\land\cos{x}\in[-1,1]\iff\cos{x}\in\left(0,1\right]\iff x\in\left(-\dfrac{\pi}{2}+2k\pi,\dfrac{\pi}{2}+2k\pi\right)\left(k\zz\right)$. 
\end{jieda}


\item 若$\tan\alpha=\sqrt3$, $\pi<\alpha<\dfrac{3\pi}{2}$, 那么$\cos\alpha-\sin\alpha$的值是\blkda{$\dfrac{-1+\sqrt3}{2}$}.
\begin{jieda}
由条件可取终边上一点$\left(-1,-\sqrt{3}\right)$. 
\end{jieda}


\item 已知$\sin\alpha=m\left(\abs{m}<1\right)$, $\dfrac{\pi}{2}<\alpha<\dfrac{3\pi}{2}$, 那么$\tan\alpha=$\blkda{$-\dfrac{m}{\sqrt{1-m^2}}$}.
\begin{jieda}
因为$\dfrac{\pi}{2}<\alpha<\dfrac{3\pi}{2}$, 故终边上点的横坐标为负. \\
$m=\dfrac{m}{1}$, 故终边上取点$\left(-\sqrt{1-m^2},m\right)$, 于是立得. 
\end{jieda}


\item 若角$\alpha$的终边落在直线$x+y=0$上, 则$\dfrac{\sin\alpha}{\sqrt{1-\sin^2\alpha}}+\dfrac{\sqrt{1-\cos^2\alpha}}{\cos\alpha}$的值等于\blkda{0}.
\begin{jieda}
考虑一下符号基本上就可以把答案选出来了. 
\end{jieda}


\item 已知$\theta\in\left(0,\dfrac{\pi}{2}\right)$, 且$\log_{\tan\theta+\cot\theta}\sin\theta=-\dfrac{3}{4}$, 
则$\log_{\tan\theta}\cos\theta=$\blkda{$\dfrac{1}{2}$}.
\begin{jieda}
$\left(\tan\theta+\cot\theta\right)^{-\frac{3}{4}}=\sin\theta$, 
即$1=\sin^4\theta\left(\tan\theta+\cot\theta\right)^3=\sin^4\theta\cdot\left(\dfrac{\sin\theta}{\cos\theta}+\dfrac{\cos\theta}{\sin\theta}\right)^3$. \\
于是$1=\dfrac{\sin^4\theta}{\sin^3\theta\cos^3\theta}=\dfrac{\sin\theta}{\cos^3\theta}=\dfrac{\tan\theta}{\cos^2\theta}$. \\
故$\cos^2\theta=\tan\theta$, 所以$\log_{\tan\theta}\cos\theta=\dfrac{1}{2}$. 
\end{jieda}



\item 证明下列恒等式: 
\begin{enumerate}[(1)]
\item $\dfrac{1+2\sin{x}\cos{x}}{\cos^2x-\sin^2x}=\dfrac{1+\tan{x}}{1-\tan{x}}$; 
\item $\dfrac{1-\sin^6x-\cos^6x}{1-\sin^4x-\cos^4x}=\dfrac{3}{2}$;
\item $\dfrac{\tan\alpha\cdot\sin\alpha}{\tan\alpha-\sin\alpha}=\dfrac{\tan\alpha+\sin\alpha}{\tan\alpha\cdot\sin\alpha}$; 
\item $\dfrac{\cos\alpha}{1+\sin\alpha}-\dfrac{\sin\alpha}{1+\cos\alpha}=\dfrac{2\left(\cos\alpha-\sin\alpha\right)}{1+\sin\alpha+\cos\alpha}$. 
\end{enumerate}
\begin{jieda}
\begin{enumerate}[(1)]
\item $\LHS=\dfrac{\left(\sin{x}+\cos{x}\right)^2}{\left(\cos{x}-\sin{x}\right)\left(\cos{x}+\sin{x}\right)}=\dfrac{\cos{x}+\sin{x}}{\cos{x}-\sin{x}}$. 

\item ``$\sin^2\theta+\cos^2\theta=1$''. 
\begin{align*}
\LHS & = \dfrac{1-\left(\sin^2x+\cos^2x\right)\left(\sin^4x-\sin^2x\cos^2x+\cos^4x\right)}{1-\left(\sin^2x+\cos^2x\right)^2+2\sin^2x\cos^2x}\\
& = \dfrac{1-\left(\sin^2x+\cos^2x\right)^2+3\sin^2x\cos^2x}{2\sin^2x\cos^2x}\\
& = \dfrac{3\sin^2x\cos^2x}{2\sin^2x\cos^2x}
\end{align*}

\item $\iff\tan^2\alpha\sin^2\alpha=\tan^2\alpha-\sin^2\alpha\iff\sin^2\alpha=\tan^2\alpha-\tan^2\alpha\sin^2\alpha$. 

\item 
\begin{align*}
\LHS & = \dfrac{1-sin\alpha}{cos\alpha} - \dfrac{1-cos\alpha}{sin\alpha} \\
& = \dfrac{sin\alpha - sin^2\alpha - \left(cos\alpha - cos^2\alpha\right)}{sin\alpha \cdot cos\alpha} \\
& = \dfrac{\left(cos\alpha - sin\alpha\right) \cdot \left(sin\alpha + cos\alpha - 1\right)}{sin\alpha \cdot cos\alpha} \\
& = \dfrac{cos\alpha - sin\alpha}{sin\alpha \cdot cos\alpha} \cdot \dfrac{\left(sin\alpha + cos\alpha\right)^2 - 1}{sin\alpha + cos\alpha + 1} \\
& = \dfrac{cos\alpha - sin\alpha}{sin\alpha \cdot cos\alpha} \cdot \dfrac{2sin\alpha cos\alpha}{sin\alpha + cos\alpha + 1} \\
& = \RHS
\end{align*}
\end{enumerate}
\end{jieda}


\item 已知$\cot\alpha=\dfrac{1}{3}$, 求$\sin\alpha$、$\cos\alpha$及$\tan\alpha$. 
\begin{jieda}
$\alpha$为第一象限角时, $\sin\alpha=\dfrac{3}{\sqrt{10}}$, $\cos\alpha=\dfrac{1}{\sqrt{10}}$, $\tan\alpha=3$; \\
$\alpha$为第三象限角时, $\sin\alpha=-\dfrac{3}{\sqrt{10}}$, $\cos\alpha=-\dfrac{1}{\sqrt{10}}$, $\tan\alpha=3$. 
\end{jieda}

\item 已知$\sin\alpha=\dfrac{m^2-1}{m^2+1}\left(m>0\right)$, 求$\cos\alpha$与$\tan\alpha$的值. 
\begin{jieda}
$\alpha$为第一或第四象限角时, $\cos\alpha=\dfrac{2m}{m^2+1}$, $\tan\alpha=\dfrac{m^2-1}{2m}$; \\
$\alpha$为第二或第三象限角时, $\cos\alpha=-\dfrac{2m}{m^2+1}$, $\tan\alpha=-\dfrac{m^2-1}{2m}$. 
\end{jieda}



\item 已知$\tan\alpha=\sqrt{2}$, 求下列各式的值:
\begin{enumerate}[(1)]
\item $\sin\alpha+2\cos\alpha$; 
\item $\dfrac{\cos\alpha-5\sin\alpha}{3\cos\alpha+\sin\alpha}$; 
\item $\dfrac{\sin^2\alpha-\sin\alpha\cos\alpha-3\cos^2\alpha}{5\sin\alpha\cos\alpha+\sin^2\alpha+1}$; 
\item $2\sin^2\alpha-\sin\alpha\cos\alpha+\cos^2\alpha$. 
\end{enumerate}
\begin{jieda}
齐次化. \\
\begin{enumerate*}[\left(1\right)]
\item $\pm\dfrac{\sqrt{6}+2\sqrt{3}}{3}$; 
\item $\dfrac{13-16\sqrt2}{7}$; 
\item $-\dfrac{1}{5}$; 
\item $\dfrac{5-\sqrt{2}}{3}$. 
\end{enumerate*}
\end{jieda}



\item
\begin{enumerate}
\item 已知$3\sin^2\alpha+2\sin^2\beta=5\sin\alpha$, 求$\sin^2\alpha+\sin^2\beta$的范围; 
\item 已知$6\sin3\alpha-\cos^22\beta=6$, 求角$\alpha$的值. 
\end{enumerate}
\begin{jieda}
\begin{enumerate}
\item $2\sin^2\beta=5\sin\alpha-3\sin^2\alpha\in[0,2]\land\sin\alpha\in[-1,1]$, 则$\sin\alpha\in[0,\frac{2}{3}]\cup\{1\}$. \\
故$E=\sin^2\alpha+\frac{5}{2}\sin\alpha-\frac{3}{2}\sin^2\alpha=-\frac{1}{2}\sin^2\alpha+\frac{5}{2}\sin\alpha\in[0,\frac{13}{9}]\cup\{2\}$. 

\item $6\sin3\alpha-6=\cos^22\beta$. \\
$\LHS\in[-12,0],\RHS\in[0,1]$, 故$6\sin3\alpha-6=\cos^22\beta=0$, \\
于是$3\alpha=2k\pi+\dfrac{\pi}{2}$, 即$\alpha=\dfrac{2k\pi}{3}+\dfrac{\pi}{6}\left(k\zz\right)$. 
\end{enumerate}
\end{jieda}

\end{enumerate}

C组：
\begin{enumerate}
    \item 设实数$x, y, z, w$满足$x+y+z+w = 1$，求$M = xw + 2yw + 3xy + 4xz + 5yz$的最大值.
    \begin{jieda}
        $w = 1 - \left(x+y+z\right)$，消去$w$有$M = x\cdot\left(1-x-y-z\right) + 2y \cdot \left(1-x-y-z\right) + 3xy  + 4xz + 5yz \\ = -x^2 + x - 2y^2 + 2y - 3z^2 + 3z = -\left(x-\dfrac{1}{2}\right)^2 - 2\left(y - \dfrac{1}{2}\right)^2 - 3\left(z - \dfrac{1}{2}\right)^2 + \dfrac{3}{2}$ \\
        故当$x = y = z = \dfrac{1}{2}, w = -\dfrac{1}{2}, M$取得最大值$\dfrac{3}{2}$
    \end{jieda}
    \item 已知实数$m$满足：当关于$x$的实系数一元二次方程$ax^2+bx+c=0$有实数根时，$\left(a-b\right)^2+\left(b-c\right)^2+\left(c-a\right)^2 \geq ma^2$恒成立，求$m$的最大值.
    \begin{jieda}
        不等式变形为$\dfrac{\left(a-b\right)^2}{a^2}+\dfrac{\left(b-c\right)^2}{a^2}+\dfrac{\left(c-a\right)^2}{a^2} \geq m$，即$\left(1-\dfrac{b}{a}\right)^2 + \left(\dfrac{b}{a} - \dfrac{c}{a}\right)^2 + \left(\dfrac{c}{a} - 1\right)^2 \geq m$。设方程的两根为$x_1, x_2$，故根据韦达定理有$\left(1+x_1+x_2\right)^2 + \left(x_1+x_2+x_1x_2\right)^2 + \left(x_1x_2-1\right)^2 \geq m$。\\
        展开得$1+x_1^2+x_2^2+2x_1 + 2x_2+2x_1x_2 + x_1^2 + x_2^2 + x_1^2x_2^2 + 2x_1x_2 + 2x_1^2x_2 + 2x_1x_2^2 + x_1^2x_2^2 -2x_1x_2 + 1 \geq m$
        整理得$2\left(x_1^2x_2^2 + x_1^2x_2 + x_1x_2^2 + x_1^2 + x_2^2 + x_1x_2 + x_1 + x_2 + 1\right) \geq m$ \\
        观察式子可知$x_1, x_2$对称，共9项，故可尝试因式分解。得到$2\left(x_1^2 + x_1 + 1\right)\left(x_2^2 + x_2 + 1\right) \geq m$。可以看成以$x_1, x_2$为变量的两个二次函数相乘，该二次函数当自变量取$-\dfrac{1}{2}$时取最小值$\dfrac{3}{4}$，因此$m \leq 2 \times \dfrac{3}{4} \times \dfrac{3}{4} = \dfrac{9}{8}$
    \end{jieda}
    \item 设正数$x, y$满足$x^2+y^2+\dfrac{1}{x} + \dfrac{1}{y} = \dfrac{27}{4}$，求$P = \dfrac{15}{x} - \dfrac{3}{4y}$的最小值.
    \begin{jieda}
        $x^2+y^2+\dfrac{16}{x} + \dfrac{1}{4y} = \dfrac{27}{4} + P$ \\
        $\dfrac{27}{4} + P \geq \sqrt[3]{x^2 \cdot \dfrac{8}{x} \cdot \dfrac{8}{x}} + \sqrt[3]{y^2 \cdot \dfrac{1}{8y} \cdot \dfrac{1}{8y}} = 4 + \dfrac{1}{4} = \dfrac{25}{4}$
    \end{jieda}
    \item 记$F\left(x, y\right) = \left(x-y\right)^2+\left(\dfrac{x}{3} + \dfrac{3}{y}\right)^2\left(y \neq 0\right)$，求$F\left(x, y\right)$的最小值.
    \begin{jieda}
        $F\left(x, y\right)$为点$A\left(x, \dfrac{x}{3}\right)$到点$B\left(y, -\dfrac{3}{y}\right)$的距离平方. \\
        由几何意义知应寻找双曲线上斜率为$\dfrac{1}{3}$的点，根据反比例函数的导数可得该点应为$\left(-3, 1\right)$或$\left(3, -1\right)$ \\
        根据两平行直线间的距离公式知$F\left(x, y\right)$的最小值为$\left(\dfrac{6}{\sqrt{10}}\right)^2 = \dfrac{18}{5}$
    \end{jieda}
\end{enumerate}



\section{诱导公式(1)}
\zsd 奇变偶不变，符号看象限——偶

A组：
\begin{enumerate}
\item 若 $P\left(-3,4\right)$ 为角 $a$ 终边上一点, 则 $\cos \left(2 \pi-a\right)=$ \blkda{$-\dfrac{3}{5}$}.
\item 已知角 $\theta$ 的终边过点 $\left(1,-2\right)$, 则 $\cos \left(\pi+\theta\right)=$ \blkda{$-\dfrac{\sqrt{5}}{5}$}.
\item 若 $\cos \alpha=-\frac{3}{5}, a$ 是第二象限角, 则 $\sin \left(\alpha+\pi\right)=$ \blkda{$-\dfrac{4}{5}$}.
\item 已知 $1-\cos \left(\pi-\alpha\right)=2 \sin \alpha$, 则 $\tan \alpha=$\blkda{0或$\dfrac{3}{4}$}.
\item 下列等式中不恒成立的是\dotfill{(\kh{D})}
\xx{$\sin \left(2 \pi-\alpha\right)=-\sin \alpha$}
{$\cos \left(2 \pi-\alpha\right)=\cos \alpha$}
{$\tan \left(2 \pi-\alpha\right)=-\tan \alpha$}
{$\cot \left(2 \pi-\alpha\right)=\cot \alpha$}

\item 已知角 $\alpha$ 和 $\beta$ 的终边关于 $y$ 轴对称，下列等式恒成立的是\dotfill{(\kh{A})}
\xx{$\sin \alpha=\sin \beta$}
{$\cos \alpha=\cos \beta$}
{$\tan \alpha=\tan \beta$}
{$\cot \alpha=\cot \beta$}

\item 已知 $\sin \alpha=\frac{3}{5}$, 且 $\textbf{}$ 是第二象限角, 则 $\cos \left(\pi-\alpha\right)+\sin \left(\pi+\alpha\right)$ 的值等于\blkda{$\dfrac{1}{5}$}.

\item 函数 $y=\frac{|\sin \alpha|}{\sin \alpha}+\frac{2 \cos \alpha}{|\cos \alpha|}+\frac{|\tan \alpha|}{\tan \alpha}+\frac{2 \cot \alpha}{|\cot \alpha|}$ 的值域为\blkda{$\{-4, -2, 0, 6\}$}.


\end{enumerate}

B组：
\begin{enumerate}
\item 若$\cos\left(\pi-\alpha\right)=\dfrac{\sqrt5}{3}$, 且$\alpha\in\left(\dfrac{\pi}{2},\pi\right)$, 则$\sin\left(\pi+\alpha\right)$=\blkda{$-\dfrac{2}{3}$}. 

\item 已知$\sin\left(\alpha-\pi\right)=\dfrac{2}{3}$, 且$\alpha\in\left(-\dfrac{\pi}{2},0\right)$, 则$\tan\alpha$=\blkda{$-\dfrac{2}{\sqrt{5}}$}. 

\item 已知$\sin\left(\alpha-\dfrac{2\pi}{3}\right)=\dfrac{1}{4}$, 则$\sin\left(\alpha+\dfrac{\pi}{3}\right)$=\blkda{$-\dfrac{1}{4}$}. 


\item 已知$\cos\left(\dfrac{\pi}{6}+\theta\right)=\dfrac{\sqrt{3}}{3}$, 则$\cos\left(\dfrac{5\pi}{6}-\theta\right)$=\blkda{$\dfrac{-\sqrt{3}}{3}$}. 

\item 已知$\cos\left(11\pi-3\right)=p$, 用$p$表示$\tan\left(-3\right)$=\blkda{$\dfrac{\sqrt{1-p^2}}{p}$}.  


\item 若$\tan\alpha=-2$, 则$\sin\left(\alpha-\pi\right)\cdot\cos\left(\pi+\alpha\right)$=\blkda{$-\dfrac{2}{5}$}. 

\item $\sin^2\left(\alpha+\pi\right)-\cos\left(\alpha+\pi\right)\cos\left(-\alpha\right)+1$的值是\blkda{2}.

\item 在$\triangle{ABC}$中, 给出下列四个式子: 
\ding{172} $\sin\left(A+B\right)+\sin{C}$, 
\ding{173} $\cos\left(A+B\right)+\cos{C}$, 
\ding{174} $\sin\left(2A+2B\right)+\sin2C$, 
\ding{175} $\cos\left(2A+2B\right)+\cos2C$, 其中恒为常数的是\blkda{\ding{173}\ding{174}}.


\item 函数式$\sqrt{1+2\sin\left(\pi-2\right)\cos\left(\pi+2\right)}$化简的结果是\dotfill{(\kh{A})}
\xx{$\sin2-\cos2$}
{$\pm\left(\sin2-\cos2\right)$}
{$\cos2-\sin2$}
{以上结论都不对}

\item 当$n\zz$时, 在
\ding{172} $\sin\left(n\pi+\frac{\pi}{3}\right)$; 
\ding{173} $\sin\left(2n\pi\pm\frac{\pi}{3}\right)$; 
\ding{174} $\sin\left(n\pi+\left(-1\right)^n\frac{\pi}{3}\right)$; 
\ding{175} $\cos\left(2n\pi\pm\frac{\pi}{6}\right)$; 
\ding{176} $\cos\left(n\pi+\left(-1\right)^n\frac{\pi}{6}\right)$中与$\sin\frac{\pi}{3}$相等的是\blkda{\ding{174}\ding{175}}.


\item 计算 $\cos 1^{\circ}+\cos 2^{\circ}+\cos 3^{\circ}+\cdots+\cos 177^{\circ}+\cos 178^{\circ}+\cos 179^{\circ}=$\blkda{0}.

\item 已知 $\cos 170^{\circ}=m$, 则 $\tan 10^{\circ}$ 的值为\blkda{$-\dfrac{\sqrt{1-m^2}}{m}$}.


\item 已知$\sin\left(-\pi+\alpha\right)+2\cos\left(3\pi-\alpha\right)=0$, 计算: 
\begin{enumerate}[(1)]
\item $\dfrac{2\sin\alpha-\cos\alpha}{\sin\alpha+3\cos\alpha}$; 
\item $\sin^2\alpha+\sin\alpha\cos\alpha-3\cos^2\alpha$. 
\end{enumerate}
\begin{jieda}
条件可化为$-\sin\alpha-2\cos\alpha=0$, 故$\tan\alpha=-2$. 
\begin{enumerate}[(1)]
\item $E=\dfrac{2\tan\alpha-1}{\tan\alpha+3}=-5$. 
\item $E=\dfrac{\sin^2\alpha+\sin\alpha\cos\alpha-3\cos^2\alpha}{\sin^2\alpha+\cos^2\alpha}=\dfrac{\tan^2\alpha+\tan\alpha-3}{\tan^2\alpha+1}=-\dfrac{1}{5}$
\end{enumerate}
\end{jieda}


\item 设$f\left(\theta\right)=\dfrac{2\cos^3\theta+\sin^2\left(2\pi-\theta\right)+\cos\left(-\theta\right)-3}{2+2\cos^2\left(\pi+\theta\right)+\cos\left(2\pi-\theta\right)}$, 求$f\left(\frac{\pi}{3}\right)$的值. 
\begin{jieda}
$f\left(\theta\right)=\dfrac{2\cos^3\theta+1-\cos^2\theta+\cos\theta-3}{2\cos^2\theta+\cos\theta-2}=\cos\theta-1$. \\
故$f\left(\frac{\pi}{3}\right)=-\dfrac{1}{2}$. 
\end{jieda}


\item 已知$\sin\left(3\pi-\alpha\right)=\sqrt2\sin\left(2\pi+\beta\right)$, $\sqrt3\cos\left(-\alpha\right)=-\sqrt2\cos\left(\pi+\beta\right)$, 
且$0<\alpha<\pi$, $0<\beta<\pi$, 求$\sin\alpha,\sin\beta$. 
\begin{jieda}
$\begin{cases}
\sin\alpha=\sqrt{2}\sin\beta\\
\sqrt{3}\cos\alpha=\sqrt{2}\cos\beta
\end{cases}$, 两式两边平方再相加可得$\sin^2\alpha+3\cos^2\alpha=2$, 故$\sin^2\alpha=\dfrac{1}{2}$. \\
又$\alpha\in\left(0,\pi\right)$, 故$\sin\alpha=\dfrac{\sqrt{2}}{2}$. 
代入第一式可得$\sin\beta=\dfrac{1}{2}$. 
% $\sin\alpha=\dfrac{\sqrt{2}}{2}$, $\sin\beta=\dfrac{1}{2}$. 
\end{jieda}

\item 已知$f\left(x\right)=\dfrac{\sin\left(n\pi-x\right)\cos\left(n\pi-x\right)}{\cos\left(\left(n+1\right)\pi-x\right)}\cdot\tan\left(x-n\pi\right)\cot\left(n\pi-x\right)\left(n\zz\right)$, 求$f\left(\frac{7\pi}{6}\right)$. 
\begin{jieda}
$n$为偶数时, $f\left(x\right)=-\sin{x}$; $n$为奇数时, $f\left(x\right)=\sin{x}$. \\
$f\left(\frac{7\pi}{6}\right)=\begin{cases}
-\frac{1}{2}, & n\text{为奇数}\\
\frac{1}{2}, & n\text{为偶数}
\end{cases}$. 
\end{jieda}

\end{enumerate}
C组：
\begin{enumerate}
    \item 已知函数$y = |sinx + cosx + tanx + cotx + secx + cscx|$，求函数的最小值.
    \begin{jieda}
        $y = |sinx + cosx + \dfrac{1}{sinx \cdot cosx} + \dfrac{sinx + cosx}{sinx \cdot cosx}|$，\\
        令$t = sinx + cosx, t \in \left[-\sqrt{2}, -1\right) \cup \left(-1, 1\right) \cup \left(1, \sqrt{2}\right]$，故$y = |t + \dfrac{1}{\frac{t^2 - 1}{2}} + \dfrac{t}{\frac{t^2 - 1}{2}}| = |\left(t-1\right) + \dfrac{2}{t-1} + 1| \geq 2\sqrt{2} - 1$，当$t = 1-\sqrt{2}$时取等号.
    \end{jieda}
    \item 已知$x, y, z \in \mathbf{R^+}$，满足$x+y+z = xyz$，求函数$f\left(x,y,z\right) = x^2\left(yz-1\right)+y^2\left(zx-1\right)+z^2\left(xy-1\right)$的最小值.
    \begin{jieda}
        $f\left(x,y,z\right) = x^2\left(yz-1\right)+y^2\left(zx-1\right)+z^2\left(xy-1\right) = \left(x+y+z\right)\left(xyz\right)-\left(x^2+y^2+z^2\right) \\
        = \left(x+y+z\right)^2-\left(x^2+y^2+z^2\right) = 2\left(xy+yz+zx\right)$ \\
        由柯西不等式知$\left(xy+yz+zx\right)\left(z+x+y\right) \geq \left(3\sqrt{xyz}\right)^2 = 9xyz = 9\left(x+y+z\right)$，故$2\left(xy+yz+zx\right) \geq 18$，当$x = y = z = \sqrt{3}$时取等.
    \end{jieda}
    
    \item 已知实数$x, y$满足：$17\left(x^2+y^2\right)-30xy-16 = 0$，求$\sqrt{16x^2 + 4y^2-16xy - 12x+6y + 9}$的最大值.
    \begin{jieda}
由 $17\left(x^2+y^2\right)-30 x y-16=0$ 可得, $\left(x+y\right)^2+16\left(x-y\right)^2=16$,
而
$$
\begin{aligned}
&\sqrt{16 x^2+4 y^2-16 x y-12 x+6 y+9} \\
& =\sqrt{\left(4 x-2 y\right)^2-3\left(4 x-2 y\right)+9} \\
& =\sqrt{[3\left(x-y\right)+\left(x+y\right)]^2-3[3\left(x-y\right)+\left(x+y\right)]+9} .
\end{aligned}
$$


令 $x+y=4 \cos \theta, x-y=\sin \theta$,
则 $3\left(x-y\right)+\left(x+y\right)=3 \sin \theta+4 \cos \theta=5 \sin \left(\theta+\varphi\right)$,
其中, $\sin \varphi=\frac{4}{5}, \cos \varphi=\frac{3}{5}, 3\left(x-y\right)+\left(x+y\right) \in[-5,5]$,\\
故当 $3\left(x-y\right)+\left(x+y\right)=-5$ 时,
$\sqrt{16 x^2+4 y^2-16 x y-12 x+6 y+9}=\sqrt{\left(-5\right)^2-3\left(-5\right)+9} =7$

    \end{jieda}
    
\end{enumerate}

\section{诱导公式(2)}
\zsd 奇变偶不变，符号看象限——奇

A组：
\begin{enumerate}
\item 已知 $\sin \alpha=\frac{2 \sqrt{5}}{5}$, 则 $\cos \left(\alpha-\frac{\pi}{2}\right)=$\blkda{$\frac{2\sqrt{5}}{5}$}.

\item 若 $\sin a=\frac{1}{3}, \alpha\in\left(\frac{\pi}{2}, \pi\right)$, 则 $\sin \left(\alpha-\frac{3 \pi}{2}\right)$ 的值为\blkda{$-\frac{2\sqrt{2}}{3}$}.

\item 已知点 $A$ 的坐标为 $\left(3,4\right)$, 将 $O A$ 绕坐标原点 $O$ 顺时针旋转 $\frac{\pi}{2}$ 至 $O A^{\prime}$, 则点 $A^{\prime}$ 的坐标为\blkda{$\left(4,-3\right)$}.
\item 下列等式中恒成立的是\dotfill{(\kh{A})}
\xx{$\sin \left(\frac{3 \pi}{2}-\alpha\right)=-\cos a$}
{$\cos \left(\frac{3 \pi}{2}-\alpha\right)=\sin a$}
{$\tan \left(\frac{3 \pi}{2}-\alpha\right)=-\cot \alpha$}
{$\cot \left(\frac{3 \pi}{2}-\alpha\right)=\cot \alpha$}


\item 已知角 $\alpha$ 的终边经过点 $P\left(-3,-4\right)$, 则 $\cos \left(\frac{\pi}{2}+\alpha\right)$ 的值为\blkda{$\frac{4}{5}$}.


\end{enumerate}

B组：
\begin{enumerate}
\item 已知$\sin\left(\frac{\pi}{2}-\alpha\right)=\frac{3}{5}$, 则$\cos\left(\pi+\alpha\right)$的值为\blkda{$-\dfrac{3}{5}$}. 

\item 已知$\cos\left(\frac{\pi}{2}+\varphi\right)=\dfrac{\sqrt3}{2}$, 且$\abs{\varphi}<\dfrac{\pi}{2}$, 则$\tan\varphi$=\blkda{$-\sqrt{3}$}. 
\begin{jieda}
$-\sin\varphi=\dfrac{\sqrt{3}}{2}$, 故$\sin\varphi=-\dfrac{\sqrt{3}}{2}$. \\
又$\abs{\varphi}<\dfrac{\pi}{2}$, 故$-\dfrac{\pi}{2}<\varphi<0$, 于是$\tan\varphi=-\sqrt{3}$. 
\end{jieda}

\item 如果$\cos\left(\pi+\alpha\right)=-\dfrac{1}{3}$, 那么$\sin\left(\frac{3\pi}{2}-\alpha\right)$=\blkda{$-\dfrac{1}{3}$}. 
\begin{jieda}
$-\dfrac{1}{3}=\cos\left(\pi+\alpha\right)=-\cos\alpha$, 故$\cos\alpha=\dfrac{1}{3}$. \\
于是$E=-\cos\alpha=-\dfrac{1}{3}$. 
\end{jieda}

\item 已知函数$f\left(x\right)=\sqrt2\cos\left(x-\frac{\pi}{12}\right)$, 若$\cos\theta=\frac{3}{5}$, $\theta\in\left(\frac{3\pi}{2},2\pi\right)$, 
则$f\left(\theta-\frac{5\pi}{12}\right)$=\blkda{$-\frac{4\sqrt{2}}{5}$}. 
\begin{jieda}
$f\left(\theta-\frac{5\pi}{12}\right)=\sqrt2\cos\left(\theta-\frac{5\pi}{12}-\frac{\pi}{12}\right)=\sqrt{2}\sin\theta=-\frac{4\sqrt{2}}{5}$
\end{jieda}


\item 已知$\sin\beta=\dfrac{1}{3}$, $\sin\left(\alpha+\beta\right)=1$, 则$\sin\left(2\alpha+\beta\right)$=\blkda{$\dfrac{1}{3}$}. 
\begin{jieda}
$\sin\left(\alpha+\beta\right)=1\iff\alpha+\beta=2k\pi+\dfrac{\pi}{2}\iff\alpha=\dfrac{\pi}{2}-\beta+2k\pi\left(k\zz\right)$. \\
故$\sin\left(2\alpha+\beta\right)=\sin\left(\pi-2\beta+4k\pi+\beta\right)=\sin\beta=\dfrac{1}{3}$. 
\end{jieda}


\item 设$f\left(n\right)=\sin\left(\frac{n\pi}{4}+\alpha\right)$, 则$f\left(n\right)\cdot f\left(n+4\right)+f\left(n+2\right)\cdot f\left(n+6\right)$=\blkda{$-1$}. 
\begin{jieda}
令$\beta=\dfrac{n\pi}{4}+\alpha$, 故$E=\sin\beta\sin\left(\pi+\beta\right)+\sin\left(\frac{\pi}{2}+\beta\right)\sin\left(\pi+\frac{\pi}{2}+\beta\right)=-\sin^2\beta-\cos^2\beta=-1$. 
\end{jieda}

\item 已知$\sin\left(\frac{\pi}{2}+\theta\right)+3\cos\left(\pi-\theta\right)=\sin\left(-\theta\right)$, 则$\sin\theta\cos\theta+\cos^2\theta=$\dotfill{(\kh{C})}
\xx{$\dfrac{1}{5}$}
{$\dfrac{2}{5}$}
{$\dfrac{3}{5}$}
{$\dfrac{\sqrt5}{5}$}
\begin{jieda}
$\cos\theta-3\cos\theta=-\sin\theta$, 故$\tan\theta=2$. \\
$E=\dfrac{\sin\theta\cos\theta+\cos^2\theta}{\sin^2\theta+\cos^2\theta}=\dfrac{\tan\theta+1}{\tan^2\theta+1}=\dfrac{3}{5}$. 
\end{jieda}

\item 已知点$P\left(\sin\left(\pi+\theta\right),\sin\left(\frac{3\pi}{2}-\theta\right)\right)$在第三象限, 则角$\theta$所在的象限是\dotfill{(\kh{A})}
\xx{第一象限}{第二象限}{第三象限}{第四象限}
\begin{jieda}
$\begin{cases}
\sin\left(\pi+\theta\right)<0\\
\sin\left(\frac{3\pi}{2}-\theta\right)<0
\end{cases}$, 
于是$\begin{cases}
\sin\theta>0\\
\cos\theta>0
\end{cases}$. 
\end{jieda}

%\dotfill\left($\dfrac{7\pi}{6}$\right)

\item 已知角$\alpha$的终边上一点的坐标为$\left(\sin\frac{4\pi}{3},\cos\frac{4\pi}{3}\right)$, 则角$\alpha$的最小正值为\blkda{$\dfrac{7\pi}{6}$}.
\begin{jieda}
法一: 点$\left(-\frac{\sqrt{3}}{2},-\frac{1}{2}\right)$, 故所求为$\dfrac{7\pi}{6}$. \\
法二: $f\left(\frac{\pi}{2}-\theta\right)=g\left(\theta\right)$. 
\end{jieda}

\item 已知角$\alpha$的顶点与原点$O$重合, 始边与$x$轴的正半轴重合, 将$\alpha$的终边按顺时针方向旋转$\frac{\pi}{2}$后, 
过点$P\left(\frac{3}{5},\frac{4}{5}\right)$, 则$\cos\alpha$等于\dotfill{(\kh{A})}
\xx{$-\dfrac{4}{5}$}
{$\dfrac{4}{5}$}
{$-\dfrac{3}{5}$}
{$\dfrac{3}{5}$}
\begin{jieda}
$\beta=\alpha-\frac{\pi}{2}$, 则$\alpha=\beta+\frac{\pi}{2}$. 
故$\cos\alpha=\cos\left(\beta+\frac{\pi}{2}\right)=-\sin\beta=-\frac{4}{5}$. 
\end{jieda}

\item 已知$f\left(x\right)=\dfrac{\cos\left(\frac{\pi}{2}+x\right)\cos\left(-x\right)\sin\left(\frac{3\pi}{2}-x\right)}{\sin\left(-\pi-x\right)\cos\left(2\pi-x\right)}$. 
\begin{enumerate}[(1)]
\item 化简$f\left(x\right)$; 
\item 若$x$是第三象限角, 且$\tan{x}=2$, 求$f\left(x\right)$的值. 
\end{enumerate}
\begin{jieda}
\begin{enumerate}[(1)]
\item $f\left(x\right)=\dfrac{-\sin{x}\cos{x}\left(-\cos{x}\right)}{\sin{x}\cos{x}}=\cos{x}$. 
\item $f\left(x\right)=-\dfrac{1}{\sqrt{5}}$. 
\end{enumerate}
\end{jieda}

\item 已知$\sin\alpha$是方程$3x^2-10x-8=0$的根, 且$\alpha$为第三象限角, \\
求$\dfrac{\sin\left(\alpha+\frac{3\pi}{2}\right)\sin\left(\frac{3\pi}{2}-\alpha\right)\tan^2\left(2\pi-\alpha\right)\tan\left(\pi-\alpha\right)}
{\cos\left(\frac{\pi}{2}-\alpha\right)\cos\left(\frac{\pi}{2}+\alpha\right)}$的值. 
\begin{jieda}
易得$\sin\alpha=-\dfrac{2}{3}$. 
故
\begin{align*}
E & = \dfrac{-\cos\alpha\left(-\cos\alpha\right)\tan^2\alpha\left(-\tan\alpha\right)}{\sin\alpha\left(-\sin\alpha\right)}\\
& = \tan\alpha\\
& = \dfrac{2}{\sqrt{5}}.
\end{align*}
\end{jieda}


\item 化简
\begin{enumerate}
    \item $\sin \left(21 \pi-\alpha\right)+\cos \left(\frac{9 \pi}{2}-\alpha\right)+\tan \left(\frac{9 \pi}{4}-\alpha\right)-\sin \left(-\alpha-19 \pi\right)-\cos \left(-\alpha-\frac{27}{2} \pi\right)-\tan \left(-\alpha-\frac{7 \pi}{4}\right)$.
    \item $\frac{\sin \left(31^{\circ}+\alpha\right)}{\tan \left(27^{\circ}+\alpha\right)} \cdot \frac{\tan \left(747^{\circ}+\alpha\right)}{\cos \left(36^{\circ}+\alpha\right)} \cdot \frac{\cos \left(1116^{\circ}+\alpha\right)}{\sin \left(751^{\circ}+\alpha\right)}$.
\end{enumerate}
\begin{jieda}
    \begin{enumerate}
\item $ y_s=\sin \left(\pi-\alpha\right)+\cos \left(\frac{\pi}{2}-\alpha\right)+\tan \left(\frac{\pi}{4}-\alpha\right)-\sin \left(\pi-\alpha\right)-\cos \left(\frac{\pi}{2}-\alpha\right)-\tan \left(\frac{\pi}{4}-\alpha\right)=0 $
\item $y_s=\frac{\sin \left(31^{\circ}+\alpha\right) \cdot \cos \left(27^{\circ}+\alpha\right)}{\sin \left(27^{\circ}+\alpha\right)} \cdot \frac{\sin \left(27^{\circ}+\alpha\right)}{\cos \left(36^{\circ}+\alpha\right) \cos \left(27^{\circ}+\alpha\right)} \cdot \frac{\cos \left(36^{\circ}+\alpha\right)}{\sin \left(31^{\circ}+\alpha\right)}=1$
\end{enumerate}
\end{jieda}



\item 已知 $\alpha$ 是第三象限角, $f\left(\alpha\right)=\frac{\sin \left(\alpha-\frac{\pi}{2}\right)+\cos \left(\frac{3 \pi}{2}+\alpha\right)}{\cos \left(-\alpha-2021 \pi\right)+\sin \left(2022 \pi-\alpha\right)}$.
\begin{enumerate}
    \item 若 $\alpha=-\frac{32 \pi}{3}$, 求 $f\left(\alpha\right)$ 的值;
    \item 若 $\tan \alpha=2$, 求 $f\left(\alpha\right)$ 的值.
\end{enumerate}
\begin{jieda}
$f_{\left(a\right)}=\frac{-\cos \alpha+\sin \alpha}{-\cos \alpha-\sin \alpha}$
    \begin{enumerate}
        \item $f\left(\alpha\right)=f\left(\frac{4}{3} \pi\right)=-2+\sqrt{3}$
        \item $f\left(\alpha\right)=\frac{-1+\tan \alpha}{-1-\tan \alpha}=-\frac{1}{3}$
    \end{enumerate}
\end{jieda}


\end{enumerate}

C组：
\begin{enumerate}
\item
\begin{enumerate}
\item 已知集合 $A=\left\{a \left\lvert\, a=k \pi+\frac{\pi}{6}\right., k \in \mathbb{Z}\right\}$.\\
\ding{172}是否存在 $B=[a, b]$, 使 $A \cap B=\left\{-\frac{11 \pi}{6},-\frac{5 \pi}{6}, \frac{\pi}{6}, \frac{7 \pi}{6}\right\}$ 成立? 如果存在, 求出 $a, b$的范围; 如果不存在, 说明理由.\\
\ding{173}是否存在 $B=[a, b]$, 使 $A \cap B$ 有且仅有 4 个元素? 如界存在, 求出 $b-a$ 的最大范围；如果不存在，说明遇由.
\item 所有能使 $\tan a=\tan 3$ 成立的 $a$ 组成集合 $A$, 请你写出一个集合 $B$, 使 $B \subseteq A$,且 $B$ 的元意有无限个。
\end{enumerate}
\begin{jieda}
    \begin{enumerate}
        \item \ding{172}  $a \in\left[-\frac{17}{6} \pi,-\frac{11}{6} \pi\right]$, $b \in\left[\frac{7}{6} \pi, \frac{13}{6} \pi\right)$ \\
        \ding{173} 存在满足条件的$B$, $b-a \in \left[3 \pi, 5 \pi\right)$
        \item $B=\{\alpha|\alpha=2k\pi+3,k\in\mathbb{Z}\}$
    \end{enumerate}
\end{jieda}

\item 当 $s$ 和 $t$ 取到所有实数时，求 $\left(s+7-|\cos t|\right)^2+\left(s-2|\sin t|\right)^2$ 的最小值.
\begin{jieda}
设点 $A\left(s+7, s\right)$, 点 $B\left(|\cos t|, 2|\sin t|\right)$, 则 $\left(s+7-|\cos t|\right)^2+$ $\left(s-2|\sin t|\right)^2=|A B|^2$.

注意到 $|\cos t| \geqslant 0,2|\sin t| \geqslant 0$, 不妨设 $t \in\left[0, \frac{\pi}{2}\right]$, 则点 $B$ 的坐标为 $\left(\cos t, 2 \sin t\right)$.
易知点 $A\left(s+7, s\right)$ 为直线 $x-y-7=0$ 上的动点, 点 $B\left(\cos t, 2 \sin t\right)$ 到直线 $x-y-7=0$的距离为
$$
d=\frac{|\cos t-2 \sin t-7|}{\sqrt{2}}=\frac{7+2 \sin t-\cos t}{\sqrt{2}} .
$$
又函数 $d=\frac{7+2 \sin t-\cos t}{\sqrt{2}}$ 在 $\left[0, \frac{\pi}{2}\right]$ 上单调递增, 故当 $t=0$ 时, $d$ 取得最小值 $\frac{6}{\sqrt{2}}$, 则 $|A B|^2$ 的最小值为 $\left(\frac{6}{\sqrt{2}}\right)^2=18$.
所以 $\left(s+7-|\cos t|\right)^2+\left(s-2|\sin t|\right)^2$ 的最小值为 18 .
\end{jieda}

\item 若 $a>0, b>0$, 求 $\min \left\{\max \left(a, b, \frac{1}{a^2}+\frac{1}{b^2}\right)\right\}$.
\begin{jieda}
设 $\left\{\max \left(a, b, \frac{1}{a^2}+\frac{1}{b^2}\right)\right\}=m$, 则 $a \leqslant m, b \leqslant m$, $\frac{1}{a^2}+\frac{1}{b^2} \leqslant m$, 因此 $m \geqslant \frac{2}{m^2}, m \geqslant \sqrt[3]{2}$.且当 $a=b=\sqrt[3]{2}$ 时, $m \geqslant \sqrt[3]{2}$ 所以 $m$ 的最小值为 $\sqrt[3]{2}$.所以 $\min \left\{\max \left(a, b, \frac{1}{a^2}+\frac{1}{b^2}\right)\right\}=\sqrt[3]{2}$.
\end{jieda}

\end{enumerate}


\section{求角}
\begin{enumerate}
\item 已知$\sin2x=-\dfrac{\sqrt2}{2}, x\in[-\pi,\pi]$, 则$x$=\blkda[3]{$-\dfrac{3\pi}{8},-\dfrac{\pi}{8},\dfrac{5\pi}{8},\dfrac{7\pi}{8}$}. 
\begin{jieda}
$2x\in[-2\pi,2\pi]$, 故$2x=-\dfrac{3\pi}{4},-\dfrac{\pi}{4},\dfrac{5\pi}{4},\dfrac{7\pi}{4}$, 
故$x=-\dfrac{3\pi}{8},-\dfrac{\pi}{8},\dfrac{5\pi}{8},\dfrac{7\pi}{8}$. 
\end{jieda}

\item 已知$\tan\left(2x+\dfrac{\pi}{6}\right)=-1$, 则$x$=\blkda[3]{$\dfrac{k\pi}{2}-\dfrac{5\pi}{24}\left(k\zz\right)$}. 
\begin{jieda}
$2x+\dfrac{\pi}{6}=k\pi-\dfrac{\pi}{4}$, 故$x=\dfrac{k\pi}{2}-\dfrac{5\pi}{24}\left(k\zz\right)$
\end{jieda}


\item 方程$2\sin^2x+3\cos{x}=0$, $x\in[-2\pi,2\pi]$的解集为\blkda[3]{$\left\{\pm\dfrac{4\pi}{3},\pm\dfrac{2\pi}{3}\right\}$} .
\begin{jieda}
$\iff 2-2\cos^2x+3\cos{x}=0$, 故$\cos{x}=-\dfrac{1}{2}$\left(2已舍去\right), 于是$x=2k\pi\pm\dfrac{2\pi}{3}\left(k\zz\right)$. \\
又$x\in[-2\pi,2\pi]$, 故$x=-\dfrac{4\pi}{3},-\dfrac{2\pi}{3},\dfrac{2\pi}{3},\dfrac{4\pi}{3}$. 
\end{jieda}

\item 已知$\sin\dfrac{\alpha}{2}=-\dfrac{\sqrt{3}}{2}$, 且$\alpha$是第二象限角, 则角$\alpha$=\blkda[3]{$4k\pi-\dfrac{4\pi}{3}\left(k\zz\right)$}. 
\begin{jieda}
$\dfrac{\alpha}{2}=2k\pi-\dfrac{\pi}{3}$或$2k\pi-\dfrac{2\pi}{3}$, 故$\alpha=4k\pi-\dfrac{2\pi}{3}$或$4k\pi-\dfrac{4\pi}{3}$. \\
又$\alpha$是第二象限角, 故$\alpha=4k\pi-\dfrac{4\pi}{3}\left(k\zz\right)$
\end{jieda}


\item 方程$\sin\left(x+\dfrac{\pi}{6}\right)-\cos{x}=0$的解集为\blkda[4]{$\left\{x\mid x=k\pi+\dfrac{\pi}{6},k\zz\right\}$}. 
\begin{jieda}
$\sin\alpha=\cos\beta\iff \sin\alpha=\sin\left(\dfrac{\pi}{2}-\beta\right)\iff
\alpha=\dfrac{\pi}{2}-\beta+2k\pi$或$\alpha=\dfrac{\pi}{2}+\beta+2k\pi$. 
\begin{itemize}
\item $x+\dfrac{\pi}{6}=\dfrac{\pi}{2}-x+2k\pi\iff x=k\pi+\dfrac{\pi}{6}\left(k\zz\right)$; 
\item $x+\dfrac{\pi}{6}=\dfrac{\pi}{2}+x+2k\pi\iff 2k\pi+\dfrac{\pi}{3}=0$, 无解. 
\end{itemize}
\end{jieda}


\item 已知集合$A=\left\{\alpha\mid 2\sin\alpha-1\geq 0\right\}$, $B=\left\{\alpha\mid\sqrt2\cos\alpha+1\geq 0\right\}$, \\
则$A\cap B$=\blkda[7]{$\left\{\alpha\mid 2k\pi+\dfrac{\pi}{6}\leq\alpha\leq 2k\pi+\dfrac{3\pi}{4},k\zz\right\}$}. 
\begin{jieda}
$\sin\alpha\geq\dfrac{1}{2}\iff 2k\pi+\dfrac{\pi}{6}\leq\alpha\leq 2k\pi+\dfrac{5\pi}{6}$, \\
$\cos\alpha\geq-\dfrac{1}{\sqrt{2}}\iff 2k\pi-\dfrac{3\pi}{4}\leq\alpha\leq 2k\pi+\dfrac{3\pi}{4}$. \\
于是$A\cap B=\left\{\alpha\mid 2k\pi+\dfrac{\pi}{6}\leq\alpha\leq 2k\pi+\dfrac{3\pi}{4},k\zz\right\}$. 
\end{jieda}


\item ``$\alpha=\dfrac{\pi}{6}+2k\pi\left(k\zz\right)$''是``$\cos2\alpha=\dfrac{1}{2}$''的\dotfill{(\kh{A})}
\xx{充分而不必要条件}{必要而不充分条件}{充分必要条件}{既不充分也不必要条件}


\item 若三角方程$\sin{x}=0$与$\sin2x=0$的解集分别为$E,F$, 则\dotfill{(\kh{A})}
\xx{$E\subsetneqq F$}{$E\supsetneqq F$}{$E=F$}{$E\cap F=\varnothing$}


\item 方程$\sqrt{1-\cos^2x}=\cos{x}$的解集是\dotfill{(\kh{D})}
\xx
{$\left\{x\mid x=k\pi\pm\dfrac{\pi}{4}, k\zz\right\}$}
{$\left\{x\mid x=k\pi+\dfrac{\pi}{4}, k\zz\right\}$} 
{$\left\{x\mid x=k\pi-\dfrac{\pi}{4}, k\zz\right\}$}
{$\left\{x\mid x=2k\pi\pm\dfrac{\pi}{4}, k\zz\right\}$}


\item 已知$\theta\in[0,2\pi]$, $\sin\theta$和$\cos\theta$是方程$x^2-kx+k+1=0$的两个根, 求$\theta$. 
\begin{jieda}
$\Delta\geq 0$, 则$k\in\left(-\infty,2-2\sqrt{2}]\cup[2+2\sqrt{2},+\infty\right)$. \\
于是$\begin{cases}
\sin\theta+\cos\theta=k\\
\sin\theta\cos\theta=k+1
\end{cases}$, 
而$\left(\sin\theta+\cos\theta\right)^2=1+2\sin\theta\cos\theta$, \\
故$k^2=1+2\left(k+1\right)$, 解得$k=-1$或3. \\
当$k=3$时, $\sin\theta\cos\theta=4>1$, 不符合条件, 故舍去; \\
当$k=-1$时, 解得$\left(\cos\theta,\sin\theta\right)=\left(0,-1\right)$或$\left(-1,0\right)$. 又$\theta\in[0,2\pi]$, 故$\theta=\dfrac{3\pi}{2},\pi$. 
\end{jieda}


\item 已知集合$A=\left\{-\dfrac{1}{2},3,\cos{x}\right\}$, $B=\left\{\dfrac{\sqrt{3}}{2},\sin{x}\right\}$. 
\begin{enumerate}[(1)]
\item 当$B\subsetneqq A$时, 求$x$的值;
\item 当$A\cap B=\{y\}$时, 求$x$和$y$的值. 
\end{enumerate}
\begin{jieda}
\begin{enumerate}[(1)]
\item 由题意可得$\begin{cases}
-\frac{1}{2}=\sin{x}\\
\cos{x}=\frac{\sqrt{3}}{2}
\end{cases}$, 故$x=2k\pi+\dfrac{11\pi}{6}\left(k\zz\right)$. 

\item 若$y=\dfrac{\sqrt{3}}{2}$, 则$\cos{x}=\dfrac{\sqrt{3}}{2}\land\sin{x}\neq-\dfrac{1}{2}$, 故$x=2k\pi+\dfrac{\pi}{6}\left(k\zz\right)$. \\
若$y=\sin{x}$, 则$y=\cos{x}$或$y=-\dfrac{1}{2}$, 即 (i)$x=k\pi+\dfrac{\pi}{4}$, 或(ii)$x=2k\pi+\dfrac{7\pi}{6}$或$x=2k\pi+\dfrac{11\pi}{6}$. \\
检验可得$x=2k\pi+\dfrac{11\pi}{6}$不符合条件. \\
综上, $\left(x,y\right)=\left(2k\pi+\frac{\pi}{6},\frac{\sqrt{3}}{2}\right)$, $\left(2k\pi+\frac{\pi}{4},\frac{\sqrt{2}}{2}\right)$, 
$\left(2k\pi+\frac{5\pi}{4},-\frac{\sqrt{2}}{2}\right)$, $\left(2k\pi+\frac{7\pi}{6},-\frac{1}{2}\right)$, 其中$k\zz$. 
\end{enumerate}
\end{jieda}

\item 已知关于$x$的方程$\sin^2x+\cos{x}+m=0$. 
\begin{enumerate}
\item 当$m=1$时, 求方程的解; 
\item 要使此方程有解, 试确定$m$的取值范围. 
\end{enumerate}
\begin{jieda}
\begin{enumerate}[(1)]
\item $1-\cos^2x+\cos{x}+1=0$, 解得$\cos{x}=-1$(2已舍去), 故$x=2k\pi+\pi\left(k\zz\right)$. 

\item $m=-\sin^2x-\cos{x}=\cos^2x-\cos{x}-1\in[-\frac{5}{4},1]$. 
\end{enumerate}
\end{jieda}


\item 已知角 $a$ 是第三象限角, $\tan a=\frac{1}{2}$.
\begin{enumerate}
    \item 求 $\sin a, \cos a$ 的值;
    \item 求 $\frac{1+2 \sin \left(\pi-a\right) \cos \left(-2 \pi-a\right)}{\sin ^2\left(-a\right)-\sin ^2\left(\frac{5 \pi}{2}-a\right)}$ 的值.
\end{enumerate}
\begin{jieda}
\begin{enumerate}
    \item $\tan \alpha=\frac{\sin \alpha}{\cos \alpha}=\frac{1}{2}, \sin ^2 \alpha+\cos ^2 \alpha=1$,

所以 $\left\{\begin{array}{l}\sin \alpha=\frac{\sqrt{5}}{5} \\ \cos \alpha=\frac{2 \sqrt{5}}{5}\end{array}\right.$ 或 $\left\{\begin{array}{l}\sin \alpha=-\frac{\sqrt{5}}{5} \\ \cos \alpha=-\frac{2 \sqrt{5}}{5}\end{array}\right.$ ，
而角 $\alpha$ 是第三象限角，则 $\sin \alpha<0, \cos \alpha<0$ ，故 $\left\{\begin{array}{l}\sin \alpha=-\frac{\sqrt{5}}{5} \\ \cos \alpha=-\frac{2 \sqrt{5}}{5}\end{array}\right.$

\item $\frac{1+2 \sin \left(\pi-\alpha\right) \cos \left(-2 \pi-\alpha\right)}{\sin ^2\left(-\alpha\right)-\sin ^2\left(\frac{5 \pi}{2}-\alpha\right)}=\frac{1+2 \sin \alpha \cos \left(-\alpha\right)}{\left(-\sin \alpha\right)^2-\sin ^2\left(\frac{\pi}{2}-\alpha\right)}$

$$
\begin{aligned}
& =\frac{1+2 \sin \alpha \cos \alpha}{\sin ^2 \alpha-\cos ^2 \alpha}=\frac{\sin ^2 \alpha+\cos ^2 \alpha+2 \sin \alpha \cos \alpha}{\sin ^2 \alpha-\cos ^2 \alpha} \\
& =\frac{\left(\sin \alpha+\cos \alpha\right)^2}{\left(\sin \alpha+\cos \alpha\right)\left(\sin \alpha-\cos \alpha\right)}=\frac{\sin \alpha+\cos \alpha}{\sin \alpha-\cos \alpha}=\frac{\tan \alpha+1}{\tan \alpha-1} .
\end{aligned}
$$
因为 $\tan \alpha=\frac{1}{2}$ ，所以原式 $=\frac{\frac{1}{2}+1}{\frac{1}{2}-1}=-3$.
\end{enumerate}
\end{jieda}


\end{enumerate}

C组：
\begin{enumerate}
    \item 已知函数 $f\left(x\right)$ 对于任意实数 $x, y$ 都有 $f\left(x+y\right)=f\left(x\right)+f\left(y\right)-3$, 当 $x>0$时 $f\left(x\right)<3$
    \begin{enumerate}
        \item $f\left(x\right)$ 在实数集 $R$ 上是否为单调函数? 并说明理由.
        \item 若 $f\left(6\right)=-9$, 求 $f\left(\left(\frac{1}{2}\right)^{2010}\right)$.
    \end{enumerate}
    \begin{jieda}
    \begin{enumerate}
        \item 对任意的 $x_1 、 x_2 \in R$ ，且 $x_1<x_2$ ，则 $f\left(x_2\right)=f\left(x_1+\left(x_2-x_1\right)\right)=f\left(x_1\right)+f\left(x_2-x_1\right)-3$.
因 $x_2-x_1>0$ ，所以， $f\left(x_2-x_1\right)<3$ ，即 $f\left(x_2\right)<f\left(x_1\right)$.
从而， $f\left(x\right)$ 在 $R$ 上为单调减函数.
\item 由 $f\left(6\right)=f\left(2\right)+f\left(4\right)-3=f\left(2\right)+\left(f\left(2\right)+f\left(2\right)-3\right)-3=3 f\left(2\right)-6=-9$得 $f\left(2\right)=-1$. 又由 $f\left(2\right)=f\left(1\right)+f\left(1\right)-3$, 得 $f\left(1\right)=1$.
对任意的 $a \in R$, 显然有 $f\left(a\right)=2 f\left(\frac{a}{2}\right)-3 \Rightarrow f\left(\frac{a}{2}\right)=\frac{f\left(a\right)+3}{2}$.
令 $a_1=f\left(1\right)=1, a_{n+1}=f\left(\left(\frac{1}{2}\right)^n\right)=\frac{f\left(\left(\frac{1}{2}\right)^{n-1}\right)+3}{2}=\frac{a_n+3}{2}$.
令 $b_n=a_{n+1}-a_n$. 则 $b_n=\frac{1}{2} b_{n-1}, b_1=a_2-a_1=1$.
从而， $b_n=\left(\frac{1}{2}\right)^{n-1}$ ，即 $a_{n+1}-a_n=\left(\frac{1}{2}\right)^{n-1}$.
故 $a_{n+1}=a_1+1+\frac{1}{2}+\cdots+\left(\frac{1}{2}\right)^{n-1}=3-\left(\frac{1}{2}\right)^{n-1}$.
所以， $f\left(\left(\frac{1}{2}\right)^{2010}\right)=a_{2011}=3-\left(\frac{1}{2}\right)^{2009}$.
    \end{enumerate}
    \end{jieda}

\item 设 $a, b$ 为正实数, $\frac{1}{a}+\frac{1}{b} \leq 2 \sqrt{2},\left(a-b\right)^2=4 a^3 b^3$, 求 $\log _{\mathrm{a}} b$.
\begin{jieda}
由 $\left(a-b\right)^2=4\left(a b\right)^3$ 得$\left(a+b\right)^2=4 a b+4\left(a b\right)^3 \geqslant 8\left(a b\right)^2,$
故 $a+b \geqslant 2 \sqrt{2} a b$,
又因为 $\frac{1}{a}+\frac{1}{b} \leqslant 2 \sqrt{2}$, 则 $a+b \leqslant 2 \sqrt{2} a b$,
所以 $a+b=2 \sqrt{2} a b$,
又 $\left(a+b\right)^2=4 a b+4\left(a b\right)^3$,

令 $a b=t$, 则 $8 t^2=4 t+4 t^3$, 得 $t-1=0, a b=1, \log _a b=-1$.

\end{jieda}
\end{enumerate}

\section{常用三角公式}

\subsection{两角和与差的余弦、正弦和正切公式(1)}
\zsd 两角和与差的余弦公式、两角和与差的正弦公式 \\
A组：
\begin{enumerate}
\item 若$\alpha$是第一象限角, $\sin\alpha=\dfrac{1}{4}$, 则$\cos\left(\alpha+\frac{\pi}{3}\right)$=\blkda[2]{$\dfrac{\sqrt{15}-\sqrt{3}}{8}$}. 
\item 化简: $\cos\left(\dfrac{\pi}{2}+\alpha\right)\tan\left(\dfrac{3\pi}{2}-\alpha\right)$=\blkda{$-\cos\alpha$}. 
\begin{jieda}
    $cos\left(\dfrac{\pi}{2}+\alpha\right) = -sin\alpha$，$tan\left(\dfrac{3\pi}{2}-\alpha\right) = tan\left(\dfrac{\pi}{2}-\alpha\right) = cot\alpha$ \\
    故原式$= -sin\alpha \cdot cot\alpha = -cos\alpha$ \\
    \ps 此处不应使用两角和（差）的公式.
\end{jieda}
\item 化简: $\cos\alpha+\cos\left(120^{\circ}-\alpha\right)+\cos\left(120^{\circ}+\alpha\right)$=\blkda{0}.
\begin{jieda}
    $\left \{\begin{array}{l}
        cos\left(\dfrac{2\pi}{3}+\alpha\right) = cos\left(\dfrac{2\pi}{3}\right)cos\alpha - sin\left(\dfrac{2\pi}{3}\right)sin\alpha \\
        cos\left(\dfrac{2\pi}{3}-\alpha\right) = cos\left(\dfrac{2\pi}{3}\right)cos\alpha + sin\left(\dfrac{2\pi}{3}\right)sin\alpha
    \end{array} \right. $ \\
    两式相加，即$cos\left(\dfrac{2\pi}{3}-\alpha\right)+cos\left(\dfrac{2\pi}{3}+\alpha\right) = 2cos\left(\dfrac{2\pi}{3}\right)cos\alpha = -cos\alpha$ \\
    \ps 本质上相当于推导了和差化积公式.
\end{jieda}
\item 化简$\sin\left(x+27^\circ\right)\cos\left(18^\circ-x\right)+\sin\left(18^\circ-x\right)\cos\left(x+27^\circ\right)$=\blkda{$\dfrac{\sqrt{2}}{2}$}

\item 已知$\cos\alpha=-\dfrac{1}{2}$, $\sin\beta=-\dfrac{\sqrt{3}}{2}$, 且$\alpha\in\left(\dfrac{\pi}{2},\pi\right)$, $\beta\in\left(\dfrac{3\pi}{2},2\pi\right)$,
则$\sin\left(\alpha+\beta\right)$=\blkda{$\dfrac{\sqrt{3}}{2}$}. 
\begin{jieda}
    $sin\alpha=\dfrac{\sqrt{3}}{2}$，$cos\beta = \dfrac{1}{2}$，$sin(\alpha+\beta) = \dfrac{\sqrt{3}}{2} \cdot \left(\dfrac{1}{2}\right) + \left(-\dfrac{1}{2}\right)\cdot\left(-\dfrac{\sqrt{3}}{2}\right) =  \dfrac{\sqrt{3}}{2}$ \\
    \ps 直接看出$\alpha = \dfrac{2\pi}{3}, \beta = \dfrac{5\pi}{3}$，故$sin\left(\alpha+\beta\right) = sin\left(\dfrac{7\pi}{3}\right) = sin\left(\dfrac{\pi}{3}\right) = \dfrac{\sqrt{3}}{2}$
\end{jieda}
\item 化简: $\dfrac{\sin\left(\pi-x\right)\cos\left(\frac{3\pi}{2}-x\right)\tan{x}}{\cos\left(2\pi-x\right)\cot\left(x-\frac{\pi}{2}\right)\sin\left(\frac{\pi}{2}-x\right)}=$\blkda{$\tan^2x$}
\begin{jieda}
    $\dfrac{\sin\left(\pi-x\right)\cos\left(\frac{3\pi}{2}-x\right)\tan{x}}{\cos\left(2\pi-x\right)\cot\left(x-\frac{\pi}{2}\right)\sin\left(\frac{\pi}{2}-x\right)}= \dfrac{sinx \cdot cos\left(-\dfrac{\pi}{2}-x\right)\cdot tanx}{cos(-x) \cdot \left(-cot\left(\dfrac{\pi}{2} - x\right)\right) \cdot cosx} \\
    = \dfrac{sinx \cdot cos\left(x+\dfrac{\pi}{2}\right)\cdot tanx}{cosx \cdot (-tanx) \cdot cosx} = \dfrac{sinx \cdot \left(-sinx\right)\cdot tanx}{cosx \cdot (-tanx) \cdot cosx} = tan^2x$
\end{jieda}
\item 已知$\sin\alpha=\dfrac{2}{3}$, $\cos\beta=-\dfrac{3}{4}$, 且两角在同一个象限内, 则$\cos\left(\alpha-\beta\right)$的值为\blkda[2]{$\dfrac{3\sqrt5+2\sqrt7}{12}$}
\begin{jieda}
因为两角在同一个象限内, 且$\sin\alpha>0$, $\cos\beta<0$, 故$\alpha,\beta$都是第二象限角. \\
于是$\cos\alpha=-\dfrac{\sqrt{5}}{3}$, $\sin\beta=\dfrac{\sqrt{7}}{4}$, 故$\cos\left(\alpha-\beta\right)=\dfrac{3\sqrt{5}+2\sqrt{7}}{12}$. 
\end{jieda}

\end{enumerate}
B组：
\begin{enumerate}

\item 若$\triangle{ABC}$中, $\cos{A}=-\dfrac{5}{13}$, $\sin{B}=\dfrac{4}{5}$, 则$\cos{C}$=\blkda{$\dfrac{63}{65}$}. 
\begin{jieda}
$\cos{B}=\pm\frac{3}{5}$, 若$\cos{B} = -\frac{3}{5}$，则$\sin{\left(A+B\right)} = \frac{12}{13} \cdot \left(-\frac{3}{5}\right) + \left(-\frac{5}{13}\right) \cdot \frac{4}{5} < 0$，故舍掉，因此$\cos{B}=\frac{3}{5}$，故$\cos{C} = -\cos{\left(A+B\right)} = -\left(-\frac{5}{13} \cdot \frac{3}{5} - \frac{12}{13} \cdot \frac{4}{5}\right) = \frac{63}{65}$
\end{jieda}
\item 若$\triangle{ABC}$中, 若$\cos{A}\cos{B}>\sin{A}\sin{B}$, 则$\triangle{ABC}$是\blkda{钝角}三角形. 

\item 已知$\cos\left(\alpha-\beta\right)=\dfrac{5}{13}$, $\sin\alpha=-\dfrac{4}{5}$, $\alpha\in\left(-\dfrac{\pi}{2},\dfrac{\pi}{2}\right)$, 
$\beta\in\left(0,\pi\right)$, 则$\cos\beta$=\blkda{$\dfrac{63}{65}$}. 
\begin{jieda}
$\sin\alpha=-\dfrac{4}{5}$, $\alpha\in\left(-\dfrac{\pi}{2},\dfrac{\pi}{2}\right)$, 则$\alpha\in\left(-\dfrac{\pi}{2},0\right)$, 且$\cos\alpha=\dfrac{3}{5}$. \\
于是$\alpha-\beta\in\left(-\dfrac{3\pi}{2},0\right)$. \\
又$\cos\left(\alpha-\beta\right)=\dfrac{5}{13}>0$, 故$\alpha-\beta\in\left(-\dfrac{\pi}{2},0\right)$, 且$\sin\left(\alpha-\beta\right)=-\dfrac{12}{13}$. \\
于是$\cos\beta=\cos\left(-\beta\right)=\cos\left(\alpha-\beta-\alpha\right)=\dfrac{63}{65}$. \\
\ps 三角比求值, 麻烦在于确定角的范围. 象限角与三角比的符号相结合是常见的判断角度. 
\end{jieda}
\item 若$\alpha,\beta$均为锐角, $P=\sin\left(\alpha+\beta\right)$, $Q=\sin\alpha+\sin\beta$, 则$P$与$Q$的大小关系是\blkda{$P<Q$}. 
\begin{jieda}
法一: $\alpha,\beta$均为锐角, 故$\cos\alpha,\sin\alpha,\cos\beta,\sin\beta\in\left(0,1\right)$. \\
于是$P=\sin\alpha\cos\beta+\cos\alpha\sin\beta<\sin\alpha\cdot 1+1\cdot\sin\beta=Q$, 即$P<Q$. \\
法二: $P=2\sin\frac{\alpha+\beta}{2}\cos\frac{\alpha+\beta}{2}$, $Q=2\sin\frac{\alpha+\beta}{2}\cos\frac{\alpha-\beta}{2}$. 
\end{jieda}

\item 命题甲: $3\sin\beta=\sin\left(2\alpha+\beta\right)$是命题乙: $2\tan\alpha=\tan\left(\alpha+\beta\right)$成立的\blkda[2.5]{必要不充分}条件. 
\begin{jieda}
甲$\iff3\sin\left(\alpha+\beta-\alpha\right)=\sin\left(\alpha+\beta+\alpha\right)\iff \sin\left(\alpha+\beta\right)\cos\alpha=2\cos\left(\alpha+\beta\right)\sin\alpha$. \\
若$\cos\alpha=0$, 则$\sin\alpha\neq 0$, 故$\cos\left(\alpha+\beta\right)=0$; 同理, 若$\cos\left(\alpha+\beta\right)=0$, 则$\cos\alpha=0$. \\
此时$\tan\alpha,\tan\left(\alpha+\beta\right)$无意义($\alpha=m\pi+\frac{\pi}{2},\beta=n\pi\left(m,n\zz\right)$). \\
若$\cos\alpha\cos\left(\alpha+\beta\right)\neq 0$, 则易得乙成立. \\
综上, 甲是乙的``必要不充分''条件. 
\end{jieda}

\item 若$\sin2x\sin3x=\cos2x\cos3x$, 则$x$的一个值是\dotfill{(\kh{A})}
\xx{$18^{\circ}$}{$36^{\circ}$}{$45^{\circ}$}{$30^{\circ}$}

\item 已知$\cos\alpha=-\dfrac{8}{17}$, $\cos\beta=-\dfrac{12}{13}$, 且两角为不同象限的角, 则$\cos\left(\alpha-\beta\right)$的值为\dotfill{(\kh{D})}
\xx{$\dfrac{171}{221}$}
{$-\dfrac{21}{221}$}
{$\pm\dfrac{21}{221}$}
{$\dfrac{21}{221}$}
\begin{jieda}
因为两角为不同象限的角, 且余弦都小于零, \\
故两者一为第二象限角, 一为第三象限角, 于是$\sin\alpha\sin\beta<0$. \\
故$\cos\left(\alpha-\beta\right)=\cos\alpha\cos\beta+\sin\alpha\sin\beta=\dfrac{8\cdot12-15\cdot5}{17\cdot13}=\dfrac{21}{221}$. 
\end{jieda}

\item 若$\sin\left(\alpha+\beta\right)\sin\left(\beta-\alpha\right)=m$, 则$\sin^2\beta-\sin^2\alpha$为\dotfill{(\kh{B})}
\xx{$-m$}{$m$}{$-4m$}{$4m$}
\begin{jieda}
\begin{align*}
\sin\left(\alpha+\beta\right)\sin\left(\alpha-\beta\right) 
& =\sin^2\alpha\cos^2\beta-\cos^2\alpha\sin^2\beta\\
& =\sin^2\alpha\left(1-\sin^2\beta\right)-\left(1-\sin^2\alpha\right)\sin^2\beta\\
& =\sin^2\alpha-\sin^2\beta\\
& =\cos^2\beta-\cos^2\alpha
\end{align*}
同理$\cos\left(\alpha+\beta\right)\cos\left(\alpha-\beta\right)=\cos^2\alpha-\sin^2\beta=\cos^2\beta-\sin^2\alpha$. 
\end{jieda}


\item 已知$\triangle{ABC}$中, 已知$\sin{A}\sin{B}+\sin{A}\cos{B}+\cos{A}\sin{B}+\cos{A}\cos{B}=2$, 则$\triangle{ABC}$是\dotfill{(\kh{C})}
\xx{等腰三角形}{直角三角形}{等腰直角三角形}{等腰或直角三角形}
\begin{jieda}
$\iff 2=\cos\left(A-B\right)+\sin\left(A+B\right)\iff 1=\cos\left(A-B\right)=\sin\left(A+B\right)$. 
\end{jieda}


\item 下列四个命题中, 假命题是\dotfill{(\kh{B})}
\xx
{存在这样的$\alpha,\beta$值, 使得$\sin\left(\alpha+\beta\right)=\sin\alpha\cos\beta-\cos\alpha\sin\beta$}
{不存在无穷多的$\alpha,\beta$值, 使得$\sin\left(\alpha+\beta\right)=\sin\alpha\cos\beta-\cos\alpha\sin\beta$}
{对于任意的$\alpha,\beta$值, 都有$\sin\left(\alpha+\beta\right)=\sin\alpha\cos\beta+\cos\alpha\sin\beta$}
{不存在$\alpha,\beta$值, 使得$\sin\left(\alpha+\beta\right)\neq\sin\alpha\cos\beta+\cos\alpha\sin\beta$}
\begin{jieda}
对于B选项, 只要$\cos\alpha\sin\beta=0$, 则必然有$\sin\left(\alpha+\beta\right)=\sin\alpha\cos\beta-\cos\alpha\sin\beta$. \\
C与D等价. 
\end{jieda}

\item 已知$\cos\left(\alpha+\beta\right)=\dfrac{4}{5}$, $\cos\left(\alpha-\beta\right)=-\dfrac{5}{13}$, 
$\dfrac{3\pi}{2}<\alpha+\beta<2\pi$, $\dfrac{\pi}{2}<\alpha-\beta<\pi$, 求$\cos2\alpha$和$\cos2\beta$的值. 
\begin{jieda}
结合角的范围, 易得$\sin\left(\alpha+\beta\right)=-\dfrac{3}{5}$, $\sin\left(\alpha-\beta\right)=\dfrac{12}{13}$. \\
$\cos2\alpha=\cos\left(\alpha+\beta+\alpha-\beta\right)=\dfrac{16}{65}$. \\
$\cos2\beta=\cos\left(\alpha+\beta-\left(\alpha-\beta\right)\right)=-\dfrac{56}{65}$. 
\end{jieda}


\item 已知$\cos\left(2\alpha-\beta\right)=-\dfrac{11}{14}$, $\sin\left(\alpha-2\beta\right)=\dfrac{4\sqrt{3}}{7}$, 
$0<\beta<\dfrac{\pi}{4}<\alpha<\dfrac{\pi}{2}$, 求$\alpha+\beta$的值. 
\begin{jieda}
$\dfrac{\pi}{4}<2\alpha-\beta<\pi$, 故$\sin\left(2\alpha-\beta\right)=\dfrac{5\sqrt{3}}{14}$. \\
$-\dfrac{\pi}{4}<\alpha-2\beta<\dfrac{\pi}{2}$, 故$\cos\left(\alpha-2\beta\right)=\dfrac{1}{7}$. \\
于是$\cos\left(\alpha+\beta\right)=\cos\left(\left(2\alpha-\beta\right)-\left(\alpha-2\beta\right)\right)=\dfrac{1}{2}$, 且$\alpha+\beta\in\left(\dfrac{\pi}{4},\dfrac{3\pi}{4}\right)$, \\
故$\alpha+\beta=\dfrac{\pi}{3}$. 
\end{jieda}

\item 若$\sin\alpha+\sin\beta+\sin\gamma=0$, $\cos\alpha+\cos\beta+\cos\gamma=0$, 求$\cos\left(\alpha-\beta\right)$. 
\begin{jieda}
$1=\sin^2\gamma+\cos^2\gamma=\left(\sin\alpha+\sin\beta\right)^2+\left(\cos\alpha+\cos\beta\right)^2=2+2\cos\left(\alpha-\beta\right)$, \\
展开化简易得$\cos\left(\alpha-\beta\right)=-\dfrac{1}{2}$. \\
\ps 几何上, $\alpha,\beta,\gamma$终边之间的夹角为$\ang{120}$. 
\end{jieda}

\item 已知$8\cos\left(2\alpha+\beta\right)+5\cos\beta=0$, 求$\tan\left(\alpha+\beta\right)\cdot\tan\alpha$的值. 
\begin{jieda}
$\iff 8\cos\left(\left(\alpha+\beta\right)+\alpha\right)+5\cos\left(\left(\alpha+\beta\right)-\alpha\right)=0\iff 13\cos\left(\alpha+\beta\right)\cos\alpha=3\sin\left(\alpha+\beta\right)\sin\alpha$. \\
(1) 若$\cos\alpha=0$, 则$\sin\alpha=\pm1\neq 0$, 于是$\sin\left(\alpha+\beta\right)=0$, \\
故$0=\sin\alpha\cos\beta+\cos\alpha\sin\beta$, 故$\cos\beta=0$. \\
此时$\tan\left(\alpha+\beta\right)\cdot\tan\alpha$无意义. \\
(2) 若$\cos\left(\alpha+\beta\right)=0$, 则$\sin\left(\alpha+\beta\right)=\pm1\neq0$, 于是$\sin\alpha=0$. \\
而$0=\cos\left(\alpha+\beta\right)=\cos\alpha\cos\beta-\sin\alpha\sin\beta$, 故$\cos\beta=0$. \\
此时$\tan\left(\alpha+\beta\right)\cdot\tan\alpha$无意义. \\
% 于是$\alpha=m\pi+\dfrac{\pi}{2},\alpha+\beta=n\pi$, 
% 即$\alpha=m\pi+\dfrac{\pi}{2},\beta=n\pi+\dfrac{\pi}{2}$. 
(3) 若$\cos\left(\alpha+\beta\right)\cos\alpha\neq 0$, 则$\tan\left(\alpha+\beta\right)\tan\alpha=\dfrac{13}{3}$. \\
综上, 当$\alpha=\frac{m\pi}{2}$, $\beta=n\pi+\frac{\pi}{2}$时, $\tan\left(\alpha+\beta\right)\cdot\tan\alpha$无意义; \\
当$\alpha\neq\frac{m\pi}{2}$, $\beta \neq n\pi+\frac{\pi}{2}$时, $\tan\left(\alpha+\beta\right)\tan\alpha=\dfrac{13}{3}$. 
\end{jieda}

\end{enumerate}

\clearpage

\subsection{两角和与差的余弦、正弦和正切公式(2)}
\zsd 两角和与差的正弦公式、两角和与差的正切公式 \\
A组：
\begin{enumerate}
\item $\tan\alpha=-2$, 则$\tan\left(\alpha-\frac{\pi}{3}\right)$=\blkda{$\dfrac{8+5\sqrt{3}}{11}$}. 
 
\item 化简: $\dfrac{\tan\left(20^{\circ}+\alpha\right)+\tan\left(40^{\circ}-\alpha\right)}{1-\tan\left(20^{\circ}+\alpha\right)\tan\left(40^{\circ}-\alpha\right)}$=\blkda{$\sqrt{3}$}. 
\item 已知$cos\theta = -\dfrac{3}{5}$，$\theta \in \left(0, \pi\right)$，则$sin\left(\theta + \dfrac{\pi}{4}\right) = $\blkda{$\dfrac{\sqrt{2}}{10}$}
\begin{jieda}
    $sin\theta = \dfrac{4}{5}$，$sin\left(\theta + \dfrac{\pi}{4}\right) = \dfrac{\sqrt{2}}{2}\left(sin\theta + cos\theta\right) = \dfrac{\sqrt{2}}{10}$
\end{jieda}
\item 已知$\alpha, \beta \in \left(0, \dfrac{\pi}{2}\right)$，且$sin\alpha = \dfrac{1}{3}, cos\beta = \dfrac{1}{4}$，则$sin\left(\alpha + \beta\right) = $\blkda[3]{$\dfrac{1 + 2\sqrt{30}}{12}$}
\begin{jieda}
    $cos\alpha = \dfrac{2\sqrt{2}}{3}, sin\beta = \dfrac{\sqrt{15}}{4}$，故$sin\left(\alpha + \beta\right) = \dfrac{1}{3} \cdot \dfrac{1}{4} + \dfrac{\sqrt{15}}{4} \cdot \dfrac{2\sqrt{2}}{3}  = \dfrac{1 + 2\sqrt{30}}{12}$
\end{jieda}
\item 已知$tan\alpha = 2$，求$tan\left(\alpha + \dfrac{\pi}{4}\right)$，$tan\left(\dfrac{\pi}{6} - \alpha\right)$，$cot\left(\alpha - \dfrac{\pi}{3}\right)$的值
\begin{jieda}
    \begin{dingautolist}{172}
        \item $tan\left(\alpha + \dfrac{\pi}{4}\right) = \dfrac{1+tan\alpha}{1-tan\alpha} = -3$
        \item $tan\left(\dfrac{\pi}{6} - \alpha\right) = \dfrac{\dfrac{\sqrt{3}}{3} - tan\alpha}{1 + \dfrac{\sqrt{3}}{3} \cdot tan\alpha} = \dfrac{1-2\sqrt{3}}{\sqrt{3}+2} = 8-5\sqrt{3}$
        \item $cot\left(\alpha - \dfrac{\pi}{3}\right) = \dfrac{1+\sqrt{3} \cdot tan\alpha}{tan\alpha - \sqrt{3}} = 8+5\sqrt{3}$ 
    \end{dingautolist}
\end{jieda}
\end{enumerate}
B组：
\begin{enumerate}
\item $\tan\alpha=-3$, $\tan\left(\alpha-\beta\right)=\frac{1}{2}$, 则$\tan\beta$=\blkda{7}
\item 已知$\tan\alpha=\dfrac{1}{7}$, $\tan\beta=\dfrac{1}{3}$, $\alpha,\beta$均为锐角, 则$\alpha+2\beta$=\blkda{$\dfrac{\pi}{4}$}
\item 已知$\tan\left(\alpha+\beta\right)=\dfrac{2}{5}$, $\tan\left(\beta-\dfrac{\pi}{4}\right)=\dfrac{1}{4}$, 则$\tan\left(\alpha+\dfrac{\pi}{4}\right)$=\blkda{$\dfrac{3}{22}$}
\item 已知$\alpha$为第三象限角, $\sin\alpha=-\dfrac{\sqrt3}{3}$, 则$\dfrac{\tan\left(\alpha+\beta\right)-\tan\beta}{1+\tan\left(\alpha+\beta\right)\tan\beta}$=\blkda{$\dfrac{\sqrt{2}}{2}$}. 
\item 若$\tan\alpha=2$, 则$\dfrac{1-\tan\left(\alpha-45^{\circ}\right)}{1+\tan\left(\alpha-45^{\circ}\right)}$的值为\blkda{$\dfrac{1}{2}$}
\begin{jieda}
$\dfrac{1-\tan\left(\alpha-45^{\circ}\right)}{1+\tan\left(\alpha-45^{\circ}\right)}=\tan\left(\ang{45}-\left(\alpha-\ang{45}\right)\right)=\cot\alpha=\dfrac{1}{2}$. 
\end{jieda}


\item 若$\tan\alpha,\tan\beta$是方程$x^2-px+q=0$的两根, $\cot\alpha,\cot\beta$是方程$x^2-rx+s=0$的两根, 
则 (1) $ps=r$; (2) $qs=1$; (3) $qr=p$; (4) $r\left(1-q\right)=p\left(s-1\right)$中成立的式子有\blkda{4}个.
\begin{jieda}
$\tan\alpha\cot\alpha=\tan\beta\cot\beta=1$. \\
于是方程$x^2-px+q=0$与$sx^2-rx+1=0$是同解方程, 且$s\neq 0$, $\Delta\geq 0$. \\
于是$\dfrac{1}{s}=\dfrac{-p}{-r}=\dfrac{q}{1}$, 即$\dfrac{1}{s}=\dfrac{p}{r}=q$. (若分母为零, 则改写为乘积的形式)\\
易得(1)(2)(3)成立, 将(1)(2)(3)代入(4)中(4)也成立. 
\end{jieda}
\item 若$\sin\alpha\sin\beta+\cos\alpha\cos\beta=0$, 那么$\sin\alpha\cos\alpha+\sin\beta\cos\beta$的值等于\blkda{0}
\begin{jieda}
$\sin\alpha\sin\beta+\cos\alpha\cos\beta=0\iff\cos\left(\alpha-\beta\right)=0\iff\alpha-\beta=k\pi+\dfrac{\pi}{2}(k \in \mathbf{Z})$. \\
于是$\sin\alpha\cos\alpha+\sin\beta\cos\beta = \sin\left(\beta+k\pi+\frac{\pi}{2}\right)\cos\left(\beta+k\pi+\frac{\pi}{2}\right)+\sin\beta\cos\beta=0$. (对$k$分奇偶讨论一下)
\end{jieda}
\item 化简: $\dfrac{\tan\left(\frac{\pi}{4}+\alpha\right)-\tan\left(\frac{\pi}{4}-\alpha\right)}{\tan2\alpha}-\tan\left(\frac{\pi}{4}+\alpha\right)\tan\left(\frac{\pi}{4}-\alpha\right)$. 
\begin{jieda}
$2\alpha=\dfrac{\pi}{4}+\alpha-\left(\dfrac{\pi}{4}-\alpha\right)$, 于是只需要展开$\tan2\alpha$即可. \\
\da 1. 
\end{jieda}
\item 求证: $\tan5\alpha-\tan3\alpha-\tan2\alpha=\tan5\alpha\tan3\alpha\tan2\alpha$. 
\begin{jieda}
$\iff \tan5\alpha=\dfrac{\tan3\alpha+\tan2\alpha}{1-\tan3\alpha\tan2\alpha}$.
\end{jieda}


\item 已知关于$x$的方程$mx^2+\left(2m-3\right)x+\left(m-2\right)=0\left(m\neq 0\right)$的两个根为$\tan\alpha,\tan\beta$, 求$\tan\left(\alpha+\beta\right)$的最小值.
\begin{jieda}
$\Delta\geq 0\land m\neq 0\iff m\leq\dfrac{9}{4}$且$ m\neq 0$. 由韦达定理有$\left \{ \begin{array}{l}
         tan\alpha + tan\beta = \dfrac{3-2m}{m}  \\
         tan\alpha \cdot tan\beta = \dfrac{m-2}{m}
        \end{array}\right.$ \\
$\tan\left(\alpha+\beta\right)=\dfrac{\tan\alpha+\tan\beta}{1-\tan\alpha\tan\beta}=\dfrac{\dfrac{3-2m}{m}}{1-\dfrac{m-2}{m}} = \dfrac{3}{2}-m\geq-\dfrac{3}{4}$. \\
故$\min=-\dfrac{3}{4}$, $m=\dfrac{9}{4}$. 
\end{jieda}

\item 设$A,B,C$是$\triangle{ABC}$的三内角, 已知$\dfrac{\tan{A}-\tan{B}}{\tan{A}+\tan{B}}=\dfrac{\sin{C}-\sin{B}}{\sin{C}}$, 
求$\cos\dfrac{B+C}{2}$的值. 
\begin{jieda}
切割化弦, 通分得
\begin{align*}
\dfrac{\sin{C}-\sin{B}}{\sin{C}}
& =\dfrac{\sin{A}\cos{B}-\sin{B}\cos{A}}{\sin{A}\cos{B}+\sin{B}\cos{A}}\\
& =\dfrac{\sin\left(A-B\right)}{\sin\left(A+B\right)}\\
& =\dfrac{\sin\left(A-B\right)}{\sin{C}}
\end{align*}
故$\sin{C}-\sin{B}=\sin\left(A-B\right)$, 即$\sin{B}=\sin\left(A+B\right)-\sin\left(A-B\right)=2\cos{A}\sin{B}$. \\
又$B\in\left(0,\pi\right)$, 则$\sin{B}\neq 0$, 故$\cos{A}=\dfrac{1}{2}$, 即$A=\dfrac{\pi}{3}$. \\
故$\cos\dfrac{B+C}{2}=\cos\dfrac{\pi}{3}=\dfrac{1}{2}$. 
\end{jieda}
\item 已知$2\sin\beta=\sin\left(2\alpha+\beta\right)$, 求$\tan\left(\alpha+\beta\right):\tan\alpha$的值. 
\begin{jieda}
$\iff 2\sin\left(\alpha+\beta-\alpha\right)=\sin\left(\alpha+\beta+\alpha\right)\iff\sin\left(\alpha+\beta\right)\cos\alpha=3\cos\left(\alpha+\beta\right)\sin\alpha$. \\
若$\cos\alpha=0$, 则$\cos\left(\alpha+\beta\right)=0$; 若$\cos\left(\alpha+\beta\right)=0$, 则$\cos\alpha=0$. \\
故当$\cos\alpha\cos\left(\alpha+\beta\right)=0$时, $\tan\left(\alpha+\beta\right):\tan\alpha$不存在; \\
当$\cos\alpha\cos\left(\alpha+\beta\right)\neq 0$时, $\tan\left(\alpha+\beta\right):\tan\alpha=3:1$. 
\end{jieda}

\item 已知$0^\circ<\alpha<\beta<90^\circ$, 且$\sin\alpha,\sin\beta$是方程$x^2-\left(\sqrt2\cos40^\circ\right)x+\cos^240^{\circ}-\frac{1}{2}=0$的两个实根, 
求$\sin\left(\beta-5\alpha\right)$和$\sin\left(\beta-2\alpha\right)$的值. 
\begin{jieda}
法一: 方程的解为$x=\dfrac{\sqrt{2}}{2}\left(\cos\ang{40}\pm\sin\ang{40}\right)$, 即$x=\sin\ang{85}$或$\sin\ang{5}$. \\
于是结合条件可得$\alpha=\ang{5},\beta=\ang{85}$. \\
法二: $\begin{cases}
\sin\alpha+\sin\beta=\sqrt{2}\cos\ang{40}\\
\sin\alpha\sin\beta=\cos^2\ang{40}-\frac{1}{2}
\end{cases}$, \\
于是$\left(\sin\alpha+\sin\beta\right)^2=2\sin\alpha\sin\beta+1$, 故$\sin^2\alpha+\sin^2\beta=1$. \\
于是$\sin^2\alpha=\cos^2\beta$. 又$0^\circ<\alpha<\beta<90^\circ$, 故$\sin\alpha=\cos\beta$, 即$\alpha+\beta=\ang{90}$. \\
代入第一式可得, $\sin\alpha+\cos\alpha=\sqrt{2}\sin\left(\alpha+\ang{45}\right)=\sqrt{2}\cos\ang{40}$, 即$\alpha+\ang{45}=\ang{50},\ang{130}$. \\
结合条件可得$\alpha=\ang{5}$, $\beta=\ang{85}$. \\
于是$\sin\left(\beta-5\alpha\right)=\dfrac{\sqrt{3}}{2}$, $\sin\left(\beta-2\alpha\right)=\dfrac{\sqrt{6}+\sqrt{2}}{4}$. \\
\ps 由方程可得$\cos\ang{40}=\dfrac{\sqrt{2}}{2}\left(x\pm\sqrt{1-x^2}\right)$, 一般化为$\sqrt{2}\cos\theta=x\pm\sqrt{1-x^2}$, 对照上述过程可知$x=\sin\left(\ang{45}\pm\theta\right)$. \\
法三: $\left(\sin\beta-\sin\alpha\right)^2=\left(\sin\beta+\sin\alpha\right)^2-4\sin\alpha\sin\beta=2\sin^2\ang{40}$, 于是$\sin\beta-\sin\alpha=\sqrt{2}\sin\ang{40}$. \\
于是$\sin\alpha=\sin\ang{5}$, $\sin\beta=\sin\ang{85}$. 
\end{jieda}

\item 已知$\sin2\left(\alpha+\gamma\right)=n\sin2\beta$, 求$\dfrac{\tan\left(\alpha+\beta+\gamma\right)}{\tan\left(\alpha-\beta+\gamma\right)}$的值. 
\begin{jieda}
法一: $\iff\sin\left(\left(\alpha+\beta+\gamma\right)+\left(\alpha-\beta+\gamma\right)\right)=n\sin\left(\left(\alpha+\beta+\gamma\right)-\left(\alpha-\beta+\gamma\right)\right)$. \\
展开合并化简易得$\dfrac{\tan\left(\alpha+\beta+\gamma\right)}{\tan\left(\alpha-\beta+\gamma\right)}=\dfrac{n+1}{n-1}$. \\
法二: $\dfrac{\tan{x}+\tan{y}}{\tan{x}-\tan{y}}=\dfrac{\sin{x}\cos{y}+\cos{x}\sin{y}}{\sin{x}\cos{y}-\cos{x}\sin{y}}=\dfrac{\sin\left(x+y\right)}{\sin\left(x-y\right)}$. \\
令$a=\tan\left(\alpha+\beta+\gamma\right), b=\tan\left(\alpha-\beta+\gamma\right)$, \\
则$\dfrac{a+b}{a-b}=\dfrac{\sin2\left(\alpha+\gamma\right)}{\sin2\beta}=n$, 
故$\dfrac{a}{b}=\dfrac{n+1}{n-1}$. 
\end{jieda}
\end{enumerate}

\clearpage

\subsection{两角和与差的余弦、正弦和正切公式(3)}
\zsd 辅助角公式 \\
A组：
\begin{enumerate}
\item 用辅助角公式化简
\begin{enumerate}
    \item $\dfrac{\sqrt{3}}{2}\sin\alpha+\dfrac{1}{2}\cos\alpha$=\blkda[3]{$\sin\left(\alpha+\dfrac{\pi}{6}\right)$}. 
    
    \item $\dfrac{\sqrt{3}}{2}\sin\alpha-\dfrac{1}{2}\cos\alpha$=\blkda[3]{$\sin\left(\alpha-\dfrac{\pi}{6}\right)$}. 
    
    \item $\sin\alpha-\sqrt{3}\cos\alpha$=\blkda[3]{$2\sin\left(\alpha-\dfrac{\pi}{3}\right)$}. 
    
    \item $2\cos\alpha+2\sqrt{3}\sin\alpha$=\blkda[3]{$4\sin\left(\alpha+\dfrac{\pi}{6}\right)$}. 
    
    \item $\sin\alpha+\cos\alpha$=\blkda[3]{$\sqrt{2}\sin\left(\alpha+\dfrac{\pi}{4}\right)$}. 
    
    \item $\sin\alpha-\cos\alpha$=\blkda[3]{$\sqrt{2}\sin\left(\alpha-\dfrac{\pi}{4}\right)$}. 
    
    \item $-\sin\alpha-\cos\alpha$=\blkda[3]{$-\sqrt{2}\sin\left(\alpha+\dfrac{\pi}{4}\right)$}. 
    
    \item $3\sin\alpha-4\cos\alpha$=\blkda[4]{$5\sin\left(\alpha-\arccos\dfrac{3}{5}\right)$}. 
    \begin{jieda}
    或$5\sin\left(\alpha-\arctan\dfrac{4}{3}\right)$，$\ang{53}$不可以, 因为$\ang{53}\neq\arccos\frac{3}{5}$. 
    \end{jieda}
\end{enumerate}
\end{enumerate}
B组：
\begin{enumerate}
    \item $\sqrt{3}\sin\left(\alpha+\dfrac{\pi}{6}\right)-3\cos\left(\alpha+\dfrac{\pi}{6}\right)$=\blkda[3]{$2\sqrt{3}\sin\left(\alpha-\dfrac{\pi}{6}\right)$}. 

    \item $\dfrac{\sqrt{2}}{6}\sin\left(\dfrac{\pi}{3}-\alpha\right)+\dfrac{\sqrt{6}}{6}\cos\left(\dfrac{\pi}{3}-\alpha\right)$
    =\blkda[3]{$\dfrac{\sqrt{2}}{3}\sin\left(\alpha+\dfrac{\pi}{3}\right)$}. 
    \item 将下列各式化为$A\sin\left(\alpha+\varphi\right)\left(A>0\right)$的形式: 
    \begin{enumerate}[(1)]
    \item $\cos\left(\dfrac{\pi}{3}+x\right)+\cos{x}$; 
    
    \item $\sin\left(\dfrac{\pi}{3}-2x\right)+\cos\left(\dfrac{\pi}{3}+2x\right)$. 
    \end{enumerate}
    \begin{jieda}
    \begin{enumerate}[(1)]
    \item $\cos\left(\dfrac{\pi}{3}+x\right)+\cos{x}=\cos\dfrac{\pi}{3}\cos{x}-\sin\dfrac{\pi}{3}\sin{x}+\cos{x}=\dfrac{3}{2}\cos{x}-\dfrac{\sqrt{3}}{2}\sin{x}=\sqrt{3}\sin\left(x+\dfrac{2\pi}{3}\right)$. 
    
    \item $\sin\left(\dfrac{\pi}{3}-2x\right)+\cos\left(\dfrac{\pi}{3}+2x\right)=\sin\dfrac{\pi}{3}\cos2x-\cos\dfrac{\pi}{3}\sin2x+\cos\dfrac{\pi}{3}\cos2x-\sin\dfrac{\pi}{3}\sin2x\\
    =\dfrac{\sqrt{6}+\sqrt{2}}{2}\sin\left(2x+\dfrac{3\pi}{4}\right)$. 
    \end{enumerate}
    \end{jieda}
    
    \item 已知$m=\dfrac{\sqrt{5}\sin\alpha+1}{\cos\alpha+2}$, 求$m$的取值范围. 
    \begin{jieda}
    $m\cos\alpha+2m=\sqrt{5}\sin\alpha+1$, 即$\sqrt{5}\sin\alpha-m\cos\alpha=2m-1$, \\
    故$\sqrt{5+m^2}\sin\left(\alpha-\varphi\right)=2m-1$. \\
    于是$-1\leq \dfrac{2m-1}{\sqrt{5+m^2}}\leq 1$, 解得$m\in\left[-\dfrac{2}{3},2\right]$. 
    \end{jieda}
    
    \item 已知$n=3\sin\left(x+\ang{20}\right)+5\sin\left(x+\ang{80}\right)$, 求$n$的最大值. 
    \begin{jieda}
    令$t=x+\ang{20}$, 则
    \begin{align*}
    n & =3\sin{t}+5\sin\left(t+\ang{60}\right)\\
    & = 3\sin{t}+5\sin{t}\cos\ang{60}+5\cos{t}\sin\ang{60}\\
    & = \dfrac{11}{2}\sin{t}+\dfrac{5\sqrt{3}}{2}\cos{t}\\
    & = 7\sin\left(t+\arccos\dfrac{11}{14}\right)
    \end{align*}
    故$n_{\max}=7$. 
    \end{jieda}
    
   \item 已知$5\sqrt{3}sin{x}+5cos{x}=6$, $\sqrt{2}sin{y}+\sqrt{6}cos{y}=1$, 
    且$x\in\left(0,\dfrac{\pi}{3}\right)$, $y\in\left(\dfrac{\pi}{2},\pi\right)$, 
    求$cos(x+y)$. 
    \begin{jieda}
    根据辅助角公式有$\begin{cases}
    sin(x+\frac{\pi}{6})=\frac{3}{5}\\
    sin(y+\frac{\pi}{3})=\frac{1}{2\sqrt{2}}
    \end{cases}$. \\
    又$x+\frac{\pi}{6}\in(\frac{\pi}{6},\frac{\pi}{2})$, $y+\frac{\pi}{3}\in(\frac{5\pi}{6}, \frac{4\pi}{3})$, 
    故$\begin{cases}
    cos(x+\frac{\pi}{6})=\frac{4}{5}\\
    cos(y+\frac{\pi}{3})=-\frac{\sqrt{7}}{2\sqrt{2}}
    \end{cases}$. \\
    故$cos(x+y)=cos(x+\frac{\pi}{6}+y+\frac{\pi}{3}-\frac{\pi}{2})=sin(x+\frac{\pi}{6}+y+\frac{\pi}{3})=\dfrac{4\sqrt{2}-3\sqrt{14}}{20}$. 
    \end{jieda}

    \item 已知$\dfrac{\pi}{4} < \alpha < \dfrac{3\pi}{4}, 0 < \beta < \dfrac{\pi}{4}$，$cos\left(\dfrac{\pi}{4} + \alpha\right) = -\dfrac{3}{5}, sin\left(\dfrac{3\pi}{4} + \beta\right) = \dfrac{5}{13}$，求$sin\left(\alpha + \beta\right)$
    \begin{jieda}
        $\dfrac{\pi}{4} < \alpha < \dfrac{3\pi}{4} \Rightarrow \dfrac{\pi}{4} + \alpha \in \left(\dfrac{\pi}{2}, \pi \right)$， 因此$sin\left(\dfrac{\pi}{4}+\alpha \right) = \dfrac{4}{5}$ \\
        $ 0 < \beta < \dfrac{\pi}{4} \Rightarrow \dfrac{3\pi}{4} + \beta \in \left(\dfrac{3\pi}{4}, \pi \right)$，因此$cos\left(\dfrac{3\pi}{4} + \beta\right) = -\dfrac{12}{13}$ \\
        $sin\left(\alpha + \beta\right) = -sin\left(\alpha + \beta + \pi\right) = -sin\left(\dfrac{\pi}{4}+\alpha + \dfrac{3\pi}{4} + \beta\right) = -\left(\dfrac{4}{5} \cdot \left(-\dfrac{12}{13}\right) + \left(-\dfrac{3}{5}\right)\cdot\dfrac{5}{13}\right) = \dfrac{63}{65}$
    \end{jieda}

    \item 已知$sin\left(\alpha + \beta\right) = \dfrac{2}{3}, sin\left(\alpha - \beta\right) = \dfrac{2}{5}$，求$\dfrac{tan\alpha}{tan\beta}$.
    \begin{jieda}
        $\left \{ \begin{array}{l}
             sin\alpha \cdot cos\beta + cos\alpha \cdot sin\beta = \dfrac{2}{3}  \\
             sin\alpha \cdot cos\beta - cos\alpha \cdot sin\beta = \dfrac{2}{5}
        \end{array} \right. \Rightarrow \left \{ \begin{array}{l}
             sin\alpha \cdot cos\beta = \dfrac{8}{15}  \\
             cos\alpha \cdot sin\beta = \dfrac{2}{15}
        \end{array} \right.$，故$\dfrac{tan\alpha}{tan\beta} = 4$
    \end{jieda}

    \item 若$sin\left(\dfrac{\pi}{4} - x\right) = \dfrac{5}{13}$，且$x \in \left(0, \dfrac{\pi}{4}\right)$，求$cos\left(x+\dfrac{\pi}{4}\right)$，$\dfrac{cos2x}{cos\left(x+\dfrac{\pi}{4}\right)}$.
    \begin{jieda}
        $sin\left(\dfrac{\pi}{4} - x\right) = \dfrac{5}{13} \Rightarrow cos\left(x+\dfrac{\pi}{4}\right) = sin\left(\dfrac{\pi}{2} - \left(x+\dfrac{\pi}{4}\right)\right) = sin\left(\dfrac{\pi}{4} - x\right) = \dfrac{5}{13}$ \\
        $sin\left(x+\dfrac{\pi}{4}\right) = \dfrac{12}{13}$，
        $sin\left(x-\dfrac{\pi}{4}\right) = -\dfrac{5}{13}$，$cos\left(x-\dfrac{\pi}{4}\right) = \dfrac{12}{13}$ \\
        $cos2x = cos\left(x+\dfrac{\pi}{4} + x - \dfrac{\pi}{4}\right) = \dfrac{5}{13} \cdot \dfrac{12}{13} + \dfrac{12}{13} \cdot \dfrac{5}{13} = \dfrac{120}{169} $，因此$\dfrac{cos2x}{cos\left(x+\dfrac{\pi}{4}\right)} = \dfrac{24}{13}$

        \ps 若用二倍角公式则更方便：\\
        $cos2x = sin(2x+\dfrac{\pi}{2}) = 2sin(x+\dfrac{\pi}{4})cos(x+\dfrac{\pi}{4})$，故$\dfrac{cos2x}{cos(x+\dfrac{\pi}{4})} = 2sin(x+\dfrac{\pi}{4})$ \\
        因为$cos(x+\dfrac{\pi}{4}) = sin(\dfrac{\pi}{4} - x) = \dfrac{5}{13}$，$x \in (0, \dfrac{\pi}{4})$，故$sin(x+\dfrac{\pi}{4}) = \dfrac{12}{13}$，故$\dfrac{cos2x}{cos(x+\dfrac{\pi}{4})} = \dfrac{24}{13}$
    \end{jieda}
    \item 已知$\dfrac{\pi}{2} <\beta < \alpha < \dfrac{3\pi}{4}$，$cos\left(\alpha - \beta\right) = \dfrac{12}{13}, sin\left(\alpha + \beta\right) = -\dfrac{3}{5}$，求$sin2\alpha$
    \begin{jieda}
        $sin\left(\alpha - \beta\right) = \dfrac{5}{13}$，$cos\left(\alpha + \beta\right) = -\dfrac{4}{5}$，故
        $sin2\alpha = \left(-\dfrac{3}{5}\right) \cdot \dfrac{12}{13} + \left(-\dfrac{4}{5}\right) \cdot \dfrac{5}{13} = -\dfrac{56}{65}$
    \end{jieda}
\end{enumerate}

\clearpage
\subsection{两角和与差的余弦、正弦和正切公式(4)}
\zsd 综合应用


\begin{enumerate}
\item 在$\triangle{ABC}$中, 已知$\sin{A}=\dfrac{5}{13}$, $\cos{B}=\dfrac{4}{5}$, 则$\sin{C}$的值为\blkda{$\dfrac{56}{65}$}. 
\begin{jieda}
由题意可得$\cos{A}=\pm\dfrac{12}{13}$, $\sin{B}=\dfrac{3}{5}$. \\
$\sin{C}=\sin(\pi-A-B)=\sin(A+B)=\sin{A}\cos{B}+\cos{A}\sin{B}$. \\
当$\cos{A}=\dfrac{12}{13}$时, $\sin{C}=\dfrac{56}{65}$; \\
当$\cos{A}=-\dfrac{12}{13}$时, $\sin{C}=-\dfrac{16}{65}<0$, 故舍去. 
\end{jieda}


\item 已知$\alpha,\beta$为锐角, 且$\sin\alpha=\dfrac{8}{17}$, $\cos(\alpha-\beta)=\dfrac{21}{29}$, 
则$\cos\beta$的值为\blkda{$\dfrac{155}{493}$}.  % $\dfrac{415}{493}$不符合条件. 
\begin{jieda}
$\cos\alpha=\dfrac{15}{17}$, $\sin(\alpha-\beta)=\pm\dfrac{20}{29}$. \\
$\cos\beta=\cos(-\beta)=\cos(\alpha-\beta-\alpha)=\cos(\alpha-\beta)\cos\alpha+\sin(\alpha-\beta)\sin\beta$. \\
当$\sin(\alpha-\beta)=\dfrac{20}{29}$时, $\sin\beta=-\dfrac{132}{493}<0$, 故舍去; \\ % $\cos\beta=\dfrac{415}{493}$; \\
当$\sin(\alpha-\beta)=-\dfrac{20}{29}$时, $\sin\beta=\dfrac{468}{493}$, $\cos\beta=\dfrac{155}{493}$. \\
\ps $\alpha=\arcsin\dfrac{8}{17}$, $\alpha-\beta=\pm\arccos\dfrac{21}{29}$. \\
若$\alpha-\beta=\arccos\dfrac{21}{29}$, 则$\beta=\arcsin\dfrac{8}{17}-\arccos\dfrac{21}{29}$, 计算器可知$\beta<0$, 故舍去. \\
若$\alpha-\beta=-\arccos\dfrac{21}{29}$, 则$\beta=\arcsin\dfrac{8}{17}+\arccos\dfrac{21}{29}$, 计算器可知$\beta$为锐角, 符合条件. 
\end{jieda}

\item 已知: $\sin\alpha+\sin\beta=\dfrac{1}{2}$, $\cos\alpha+\cos\beta=\dfrac{1}{3}$, 则$\cos(\alpha-\beta)$=\blkda{$-\dfrac{59}{72}$}. 

\item 化简: $2\sin\left(\dfrac{\pi}{4}+\alpha\right)\sin\left(\dfrac{\pi}{4}-\alpha\right)$=\blkda{$\cos2\alpha$}. 

\item 设$A,B,C$为$\triangle{ABC}$的三内角, 则在下列表达式中, 一定是常数的是\blkda{\ding{172}\ding{174}\ding{175}}: \\
\ding{172} $\cos(A+B)+\cos{C}$;
\ding{173} $\sin(A+B)+\sin{C}$;
\ding{174} $\tan\dfrac{A+B}{2}\tan\dfrac{C}{2}$;
\ding{175} $\cot\dfrac{A+B}{2}\cot\dfrac{C}{2}+\cot\dfrac{A+C}{2}\cot\dfrac{B}{2}+\cot\dfrac{B+C}{2}\cot\dfrac{A}{2}$. 

\item 若$\alpha+\beta=\dfrac{\pi}{4}$, $\alpha,\beta\in\left(0,\dfrac{\pi}{2}\right)$, 
则$(1+\tan\alpha)(1+\tan\beta)$的值为\blkda{2}. 
\begin{jieda}
$1=\tan(\alpha+\beta)=\dfrac{\tan\alpha+\tan\beta}{1-\tan\alpha\tan\beta}$, 故$1-\tan\alpha\tan\beta=\tan\alpha+\tan\beta$, \\
于是$(1+\tan\alpha)(1+\tan\beta)=2$. \\
\ps $\alpha+\beta=\dfrac{\pi}{4}\iff (1+\tan\alpha)(1+\tan\beta)=2$, 只要正切有意义就是成立的. 
\end{jieda}


\item 已知关于$x$的方程$x^2+3ax+3a+1=0(a>2)$的两根分别是$\tan\alpha,\tan\beta$, 且$\alpha,\beta\in\left(-\dfrac{\pi}{2},\dfrac{\pi}{2}\right)$, 
则$\alpha+\beta$=\blkda{$-\dfrac{3\pi}{4}$}. 
\begin{jieda}
$\Delta\geq 0\land a>2\iff a>2$. \\
故$\begin{cases}
\tan\alpha+\tan\beta=-3a<0\\
\tan\alpha\tan\beta=3a+1>0
\end{cases}$, 
故$\tan\alpha,\tan\beta<0$. \\
又$\alpha,\beta\in\left(-\dfrac{\pi}{2},\dfrac{\pi}{2}\right)$, 故$\alpha,\beta\in\left(-\dfrac{\pi}{2},0\right)$, 则$\alpha+\beta\in(-\pi,0)$. \\
则$\tan(\alpha+\beta)=\dfrac{\tan\alpha+\tan\beta}{1-\tan\alpha\tan\beta}=\dfrac{-3a}{1-(3a+1)}=1$, 故$\alpha+\beta=-\dfrac{3\pi}{4}$. 
\end{jieda}


\item 已知$\cos\alpha\cos\beta=\dfrac{1}{5}$, 则$\sin\alpha\sin\beta$的取值范围是\blkda{$\left[-\dfrac{4}{5},\dfrac{4}{5}\right]$}. 
\begin{jieda}
令$t=\sin\alpha\sin\beta$, 则
$\begin{cases}
\frac{1}{5}+t=\cos(\alpha-\beta)\in[-1,1]\\
\frac{1}{5}-t=\cos(\alpha+\beta)\in[-1,1]
\end{cases}$, 
解得$t\in[-\frac{4}{5},\frac{4}{5}]$. \\
令$\alpha=\beta=\arccos\frac{1}{\sqrt{5}}$, 则$t=\frac{4}{5}$; 令$\alpha=\arccos\frac{1}{\sqrt{5}},\beta=-\arccos\frac{1}{\sqrt{5}}$, 则$t=-\frac{4}{5}$. 
\end{jieda}

\item 在$\triangle{ABC}$中, 若$\tan{A}\cdot\tan{B}>1$, 那么$\triangle{ABC}$是\dotfill(\kh{A})
\xx{锐角三角形}{直角三角形}{钝角三角形}{以上三种情况均有可能}

\item 已知$\theta$是第一象限角, 那么恒有\dotfill(\kh{D})
\xx{$\sin\dfrac{\theta}{2}>0$}
{$\sin\dfrac{\theta}{2}<\cos\dfrac{\theta}{2}$}
{$\sin\dfrac{\theta}{2}>\cos\dfrac{\theta}{2}$}
{$\tan\dfrac{\theta}{2}<1$}
\begin{jieda}
如图所示：\\
\begin{minipage}{\textwidth}
\centering
\includegraphics[scale=0.8]{./bodyxb/src/main42chap05renyijiaodesanjiaobi01.pdf}
\end{minipage}
 $\sin\dfrac{\theta}{2}$可正可负, $\sin\dfrac{\theta}{2}-\cos\dfrac{\theta}{2}$可正可负, $0<\tan\dfrac{\theta}{2}<1$. 
\end{jieda}

\item 把下列各式化为$Asin(\alpha + \varphi)$的形式. \\
\begin{enumerate*}
    \item $\dfrac{\sqrt{2}}{2}sin\alpha + \dfrac{\sqrt{2}}{2}cos\alpha$
    \item $sin\alpha - \sqrt{3}cos\alpha$
    \item $a\cdot sin\alpha + b \cdot cos\alpha$
\end{enumerate*}
\begin{jieda}
    \begin{enumerate}
        \item $sin\left(\alpha + \dfrac{\pi}{4}\right)$
        \item $2sin\left(\alpha - \dfrac{\pi}{3}\right)$
        \item $\sqrt{a^2+b^2} sin(\alpha + \varphi) \quad \left(cos \varphi = \dfrac{a}{\sqrt{a^2+b^2}}, sin\varphi = \dfrac{b}{\sqrt{a^2+b^2}}\right)$
    \end{enumerate}
\end{jieda}

\item 已知$\begin{cases}
1+\cos\alpha-\sin\beta+\sin\alpha\sin\beta=0\\
1-\cos\alpha-\cos\beta+\sin\alpha\cos\beta=0
\end{cases}$, 求$\sin\alpha$的值. 
\begin{jieda}
$\begin{cases}
1+\cos\alpha=(1-\sin\alpha)\sin\beta\\
1-\cos\alpha=(1-\sin\alpha)\cos\beta
\end{cases}$, 
两式两边平方再相加可得$2+2\cos^2\alpha=(1-\sin\alpha)^2$. \\
即$3\sin^2\alpha-2\sin\alpha-3=0$, 解得$\sin\alpha=\dfrac{1-\sqrt{10}}{3}$(另一个大于1, 已舍去). 
\end{jieda}


\item 已知$\cos\left(\alpha-\dfrac{\beta}{2}\right)=\dfrac{12}{13}$, $\sin\left(\beta-\dfrac{\alpha}{2}\right)=\dfrac{4}{5}$, 
$\alpha,\beta\in\left(0,\dfrac{\pi}{2}\right)$, 求$\cos(\alpha+\beta)$的值. 
\begin{jieda}
$\alpha-\dfrac{\beta}{2},\beta-\dfrac{\alpha}{2}\in\left(-\dfrac{\pi}{4},\dfrac{\pi}{2}\right)$. \\
因为$\sin\left(\beta-\dfrac{\alpha}{2}\right)=\dfrac{4}{5}>\dfrac{\sqrt{2}}{2}$, 故$\beta-\dfrac{\alpha}{2}\in\left(\dfrac{\pi}{4},\dfrac{\pi}{2}\right)$, 
且$\cos\left(\beta-\dfrac{\alpha}{2}\right)=\dfrac{3}{5}$. \\
因为$\cos\left(\alpha-\dfrac{\beta}{2}\right)=\dfrac{12}{13}$, 故$\sin\left(\alpha-\dfrac{\beta}{2}\right)=\pm\dfrac{5}{13}$. \\
又$\dfrac{\alpha+\beta}{2}=\alpha-\dfrac{\beta}{2}+\beta-\dfrac{\alpha}{2}$, 且$\cos(\alpha+\beta)=2\cos^2\dfrac{\alpha+\beta}{2}-1$, 故
\begin{itemize}
\item 当$\sin\left(\alpha-\dfrac{\beta}{2}\right)=\dfrac{5}{13}$时, $\cos\dfrac{\alpha+\beta}{2}=\dfrac{16}{65}$, 故$\cos(\alpha+\beta)=-\dfrac{3713}{4225}$; 
\item 当$\sin\left(\alpha-\dfrac{\beta}{2}\right)=-\dfrac{5}{13}$时, $\cos\dfrac{\alpha+\beta}{2}=\dfrac{56}{65}$, 故$\cos(\alpha+\beta)=\dfrac{2047}{4225}$. 
\end{itemize}
\ps $\alpha-\dfrac{\beta}{2}=\pm\arccos\dfrac{12}{13}$, $\beta-\dfrac{\alpha}{2}=\arcsin\dfrac{4}{5}$. \\
若$\alpha-\dfrac{\beta}{2}=\arccos\dfrac{12}{13}$, 
则$\alpha=\dfrac{2}{3}\left(\arcsin\dfrac{4}{5}+2\arccos\dfrac{12}{13}\right)$, 
$\beta=\dfrac{2}{3}\left(2\arcsin\dfrac{4}{5}+\arccos\dfrac{12}{13}\right)$, 计算器可知符合条件. \\
若$\alpha-\dfrac{\beta}{2}=-\arccos\dfrac{12}{13}$, 
则$\alpha=\dfrac{2}{3}\left(\arcsin\dfrac{4}{5}-2\arccos\dfrac{12}{13}\right)$, 
$\beta=\dfrac{2}{3}\left(2\arcsin\dfrac{4}{5}-\arccos\dfrac{12}{13}\right)$, 计算器可知符合条件. 
\end{jieda}

\item 已知$\tan(\alpha-\beta)=\dfrac{1}{2}$, $\tan\beta=-\dfrac{1}{7}$, 且$\alpha,\beta\in(0,\pi)$, 求$2\alpha-\beta$的大小. 
\begin{jieda}
$\tan\beta=-\dfrac{1}{7}>-1\land\beta\in(0,\pi)$, 则$\beta\in\left(\dfrac{3\pi}{4},\pi\right)$. \\
$\tan\alpha=\tan(\alpha-\beta+\beta)=\dfrac{\frac{1}{2}-\frac{1}{7}}{1+\frac{1}{2}\cdot\frac{1}{7}}=\dfrac{1}{3}<1$, 
又$\alpha\in(0,\pi)$, 故$\alpha\in\left(0,\dfrac{\pi}{4}\right)$. \\
于是$2\alpha-\beta\in\left(-\pi,-\dfrac{\pi}{4}\right)$. \\
则$\tan(2\alpha-\beta)=\tan(\alpha-\beta+\alpha)=\dfrac{\frac{1}{2}+\frac{1}{3}}{1-\frac{1}{2}\cdot\frac{1}{3}}=1$, 
故$2\alpha-\beta=-\dfrac{3\pi}{4}$. 
\end{jieda}


\item
\begin{enumerate}[(1)]
\item 已知$\triangle{ABC}$不是直角三角形, 求证: $\tan{nA}+\tan{nB}+\tan{nC}=\tan{nA}\tan{nB}\tan{nC}(n\zz)$; 
\item 在$\triangle{ABC}$中, 求证: $\tan\dfrac{A}{2}\tan\dfrac{B}{2}+\tan\dfrac{B}{2}\tan\dfrac{C}{2}+\tan\dfrac{C}{2}\tan\dfrac{A}{2}=1$. 
\end{enumerate}
\begin{jieda}
\begin{enumerate}[(1)]
\item $n(A+B+C)=n\pi$, 故$nC=n\pi-nA-nB$. \\
于是$\tan{nC}=\tan(n\pi-nA-nB)=-\tan(nA+nB)=-\dfrac{\tan{nA}+\tan{nB}}{1-\tan{nA}\tan{nB}}$, \\
则$\tan{nC}(1-\tan{nA}\tan{nB})=-\tan{nA}-\tan{nB}$, 整理即得. 

\item $\iff\tan\dfrac{C}{2}\cdot\left(\tan\dfrac{A}{2}+\tan\dfrac{B}{2}\right)=1-\tan\dfrac{A}{2}\tan\dfrac{B}{2}
\iff\tan\dfrac{C}{2}=\dfrac{1-\tan\frac{A}{2}\tan\frac{B}{2}}{\tan\frac{A}{2}+\tan\frac{B}{2}}$. \\
$\LHS=\tan\dfrac{\pi-A-B}{2}=\cot\dfrac{A+B}{2}=\dfrac{1}{\tan\frac{A+B}{2}}=\RHS$. 
\end{enumerate}
\end{jieda}
\end{enumerate}
\clearpage

\subsection{二倍角公式}
\zsd 二倍角正弦、余弦和正切公式 \\
A组：
\begin{enumerate}
\item 化简：
\begin{enumerate*}
    \item $cos^4\dfrac{\alpha}{2} - sin^4\dfrac{\alpha}{2}$
    \item $\dfrac{1}{1-tan\alpha} - \dfrac{1}{1+tan\alpha}$
    \item $1+2cos^2\theta - cos2\theta$
\end{enumerate*}
\begin{jieda}
    \begin{enumerate}
        \item $cos^4\dfrac{\alpha}{2} - sin^4\dfrac{\alpha}{2} = \left(cos^2\dfrac{\alpha}{2} + sin^2\dfrac{\alpha}{2}\right) \left(cos^2\dfrac{\alpha}{2} - sin^2\dfrac{\alpha}{2}\right) = cos\alpha$
        \item $\dfrac{1}{1-tan\alpha} - \dfrac{1}{1+tan\alpha} = \dfrac{2tan\alpha}{1-tan^2\alpha} = tan2\alpha$
        \item $1+2cos^2\theta - cos2\theta = 2cos^2\theta + 2sin^2\theta = 2$
    \end{enumerate}
\end{jieda}
\item 已知$\alpha \in (\dfrac{\pi}{2}, \pi), sin\alpha = \dfrac{12}{13}$，求$sin2\alpha, cos2\alpha, tan2\alpha$
\begin{jieda}
    $cos\alpha = -\dfrac{5}{13}$
    \begin{enumerate}
        \item $sin2\alpha = 2sin\alpha cos\alpha = -\dfrac{120}{169}$
        \item $cos2\alpha = cos^2\alpha - sin^2\alpha = -\dfrac{119}{169}$
        \item $tan2\alpha = \dfrac{sin2\alpha}{cos2\alpha} = \dfrac{120}{119}$
    \end{enumerate}
\end{jieda}
\item 试用$cos\theta$表示$cos3\theta$
\begin{jieda}
    \begin{align*}
    cos3\theta & = cos(2\theta+\theta) \\
    & = cos(2\theta)cos\theta - sin(2\theta) sin\theta \\
    & = (2cos^2\theta - 1)cos\theta - 2sin^2\theta cos\theta \\
    & = 2cos^3\theta - cos\theta - 2(1-cos^2\theta)cos\theta \\
    & = 4cos^3\theta - 3cos\theta.
    \end{align*}
\end{jieda}
\item 已知$\cos(\pi+\alpha)=\dfrac{1}{3}(\pi<\alpha<2\pi)$, 则$\sin2\alpha$=\blkda{$\dfrac{4\sqrt{2}}{9}$}. 
\begin{jieda}
$\cos\alpha=-\dfrac{1}{3}$, $\pi<\alpha<2\pi$, 故$\alpha$是第三象限角. \\
法一: $\sin2\alpha=2\sin\alpha\cos\alpha=2\cdot\dfrac{-2\sqrt{2}}{3}\cdot\dfrac{-1}{3}$. \\
法二: $\alpha=2\pi-\arccos(-\frac{1}{3})=\pi+\arccos\frac{1}{3}$, 然后计算器可以解决. 
\end{jieda}
\end{enumerate}
B组：
\begin{enumerate}
    \item $\sqrt{2}\sin\dfrac{\pi}{32}\cos\dfrac{\pi}{32}\cos\dfrac{\pi}{16}\cos\dfrac{\pi}{8}$=\blkda{$\dfrac{1}{8}$}. 
    
    \item 若$\alpha\in(0,\pi)$, $\sin\alpha+\cos\alpha=\dfrac{1}{3}$, 则$\cos2\alpha$=\blkda{$-\dfrac{\sqrt{17}}{9}$}. 
    \begin{jieda}
    法一: $\dfrac{1}{9}=(\sin\alpha+\cos\alpha)^2=1+2\sin\alpha\cos\alpha$, 则$2\sin\alpha\cos\alpha=-\dfrac{8}{9}<0$. \\
    又$\alpha\in(0,\pi)$, 故$\sin\alpha>0>\cos\alpha$. \\
    于是$\cos\alpha-\sin\alpha=-\sqrt{(\cos\alpha-\sin\alpha)^2}=-\sqrt{1-2\cos\alpha\sin\alpha}=-\dfrac{\sqrt{17}}{3}$. \\
    故$\cos^2\alpha=(\cos\alpha-\sin\alpha)(\cos\alpha+\sin\alpha)=-\dfrac{\sqrt{17}}{9}$. \\
    法二: $\begin{cases}
    \sin\alpha+\cos\alpha=\frac{1}{3}\\
    \sin^2\alpha+\cos^2\alpha=1
    \end{cases}$且$\alpha\in(0,\pi)$, 解得$\cos\alpha=\frac{1-\sqrt{17}}{6},\sin\alpha=\frac{1+\sqrt{17}}{6}$. \\
    法三: $\sqrt{2}\sin(\alpha+\frac{\pi}{4})=\frac{1}{3}$, 且$\alpha\in(0,\pi)$, 
    故$\alpha+\frac{\pi}{4}=\pi-\arcsin\frac{1}{3\sqrt{2}}$($\arcsin\frac{1}{3\sqrt{2}}$舍去).\\
    于是$\alpha=\frac{3\pi}{4}-\arcsin\frac{1}{3\sqrt{2}}$. 
    \end{jieda}
    
    \item 若$\theta\in\left(\dfrac{\pi}{4},\dfrac{\pi}{2}\right)$, $\sin2\theta=\dfrac{1}{16}$, 则$\cos\theta-\sin\theta$ =\blkda{$-\dfrac{\sqrt{15}}{4}$}. 
    \begin{jieda}
    因为$\theta\in\left(\dfrac{\pi}{4},\dfrac{\pi}{2}\right)$, 故$\cos\theta<\sin\theta$. \\
    于是$\cos\theta-\sin\theta=-\sqrt{(\cos\theta-\sin\theta)^2}=-\sqrt{1-\sin2\theta}=-\dfrac{\sqrt{15}}{4}$. 
    \end{jieda}
    \item 已知角$\theta$的顶点与坐标原点重合, 始边与$x$轴的正半轴重合, 终边在射线$y=2x(x\geq 0)$上, 则$\cos2\theta$=\blkda{$-\dfrac{3}{5}$}. 
    \begin{jieda}
    $\tan\theta=2$, 则$\cos2\theta=\dfrac{1-\tan^2\theta}{1+\tan^2\theta}=-\dfrac{3}{5}$
    \end{jieda}
    
    
    \item 已知$\dfrac{3\pi}{2}<\theta<2\pi$, 化简: $\sqrt{\dfrac{1}{2}-\dfrac{1}{2}\sqrt{\dfrac{1}{2}+\dfrac{1}{2}\cos2\theta}}$ =\blkda{$\sin\dfrac{\theta}{2}$}. 
    
    \item 若$\dfrac{\sin\alpha+\cos\alpha}{\sin\alpha-\cos\alpha}=\dfrac{1}{2}$, 则$\tan2\alpha$=\blkda{$\dfrac{3}{4}$}. 
    \begin{jieda}
    齐次化后得$\tan\alpha=-3$, 故$\tan2\alpha=\dfrac{2\tan\alpha}{1-\tan^2\alpha}=\dfrac{3}{4}$. 
    \end{jieda}
    
    \item 已知$\sin\alpha\cos\alpha=-\dfrac{1}{8}$, 且$\dfrac{\pi}{2}<\alpha<\dfrac{3\pi}{2}$, 则$\sin\alpha-\cos\alpha$=\blkda{$\dfrac{\sqrt{5}}{2}$}. 
    \begin{jieda}
    因为$\sin\alpha\cos\alpha<0$, 且$\dfrac{\pi}{2}<\alpha<\dfrac{3\pi}{2}$, 故$\alpha\in\left(\dfrac{\pi}{2},\pi\right)$, 则$\sin\alpha>0>\cos\alpha$. 
    \end{jieda}
    \item 若$sin(\dfrac{\pi}{4} - x) = \dfrac{5}{13}$，且$x \in (0, \dfrac{\pi}{4})$，则$\dfrac{cos2x}{cos(x+\dfrac{\pi}{4})} = $\blkda{$\dfrac{24}{13}$}
    \begin{jieda}
        $cos2x = sin(2x+\dfrac{\pi}{2}) = 2sin(x+\dfrac{\pi}{4})cos(x+\dfrac{\pi}{4})$，故$\dfrac{cos2x}{cos(x+\dfrac{\pi}{4})} = 2sin(x+\dfrac{\pi}{4})$ \\
        因为$cos(x+\dfrac{\pi}{4}) = sin(\dfrac{\pi}{4} - x) = \dfrac{5}{13}$，$x \in (0, \dfrac{\pi}{4})$，故$sin(x+\dfrac{\pi}{4}) = \dfrac{12}{13}$，故$\dfrac{cos2x}{cos(x+\dfrac{\pi}{4})} = \dfrac{24}{13}$
    \end{jieda}
    \item 已知$\alpha,\beta$都是锐角, 且$3\sin^2\alpha+2\sin^2\beta=1$, $3\sin2\alpha-2\sin2\beta=0$, 那么$\alpha,\beta$之间的关系是\dotfill(\kh{D})
    \xx{$\alpha+\beta=\dfrac{\pi}{4}$}{$\alpha-\beta=\dfrac{\pi}{4}$}{$\alpha+2\beta=\dfrac{\pi}{4}$}{$\alpha+2\beta=\dfrac{\pi}{2}$}
    \begin{jieda}
    \begin{enumerate}
        \item 法一: $3\sin^2\alpha=1-2\sin^2\beta=\cos2\beta$, $3\sin\alpha\cos\alpha=\sin2\beta$, 故两式相除有$\tan2\beta=\cot\alpha$. 
        \item 法二: $3\sin^2\alpha=1-2\sin^2\beta=\cos2\beta$, $3\sin\alpha\cos\alpha=\sin2\beta$,\\
        故$\left(3\sin^2\alpha\right)^2+\left(3\sin\alpha\cos\alpha\right)^2=\cos^22\beta+\sin^22\beta=1$, \\
        解得$\sin\alpha=\dfrac{1}{3}=\cos2\beta$. 
    \end{enumerate}
    \end{jieda}
    
    \item 设$\sin\theta+\cos\theta=\dfrac{\sqrt{2}}{3}$, $\dfrac{\pi}{2}<\theta<\pi$, 求$\tan\theta-\cot\theta$的值. 
    \begin{jieda}
    由条件可得$\sin\theta\cos\theta=-\dfrac{7}{18}$, $\sin\theta-\cos\theta=\dfrac{4}{3}$. 
    \begin{align*}
    \tan\theta-\cot\theta & = \dfrac{\sin\theta}{\cos\theta}-\dfrac{\cos\theta}{\sin\theta}\\
    & = \dfrac{\sin^2\theta-\cos^2\theta}{\sin\theta\cos\theta}\\
    & = \dfrac{(\sin\theta-\cos\theta)(\sin\theta+\cos\theta)}{\sin\theta\cos\theta}\\
    & = -\dfrac{8\sqrt{2}}{7}. 
    \end{align*}
    \ps $2\sin\theta\cos\theta=(\sin\theta+\cos\theta)^2-1=1-(\sin\theta-\cos\theta)^2$. 
    \end{jieda}
    
    
    \item 已知$\tan\left(\dfrac{\pi}{4}+\alpha\right)=\dfrac{1}{2}$. 
    \begin{enumerate}[(1)]
    \item 求$\tan\alpha$的值; 
    \item 求$\dfrac{\sin2\alpha-\cos^2\alpha}{1+\cos2\alpha}$的值. 
    \end{enumerate}
    \begin{jieda}
    \begin{enumerate}[(1)]
    \item $\tan\alpha=\tan\left(\alpha+\dfrac{\pi}{4}-\dfrac{\pi}{4}\right)=-\dfrac{1}{3}$. 
    
    \item 齐次化: $\dfrac{\sin2\alpha-\cos^2\alpha}{1+\cos2\alpha}=\dfrac{2\sin\alpha\cos\alpha-\cos^2\alpha}{1+\cos^2\alpha-\sin^2\alpha}=\tan\alpha-\dfrac{1}{2}=-\dfrac{5}{6}$. 
    \end{enumerate}
    \end{jieda}
    
    
    \item 已知$6\sin^2\alpha+\sin\alpha\cos\alpha-2\cos^2\alpha=0$, $\alpha\in\left[\dfrac{\pi}{2},\pi\right]$, 求$\sin\left(2\alpha+\dfrac{\pi}{3}\right)$的值. 
    \begin{jieda}
    $\iff 6\tan^2\alpha+\tan\alpha-2=0$. 又$\alpha\in\left[\dfrac{\pi}{2},\pi\right]$, 故$\tan\alpha=-\dfrac{2}{3}$. \\
    故$\sin\left(2\alpha+\dfrac{\pi}{3}\right)=\sin2\alpha\cos\dfrac{\pi}{3}+\cos2\alpha\sin\dfrac{\pi}{3}=f(\tan\alpha)=\dfrac{5\sqrt{3}-12}{26}$. 
    \end{jieda}
    
\end{enumerate}
C组：
\begin{enumerate}
    \item 对于集合$S=\{\theta_1,\theta_2,\cdots,\theta_n\}$和常数$\theta_0$, \\
    定义$\mu=\dfrac{\cos^2(\theta_1-\theta_0)+\cos^2(\theta_2-\theta_0)+\cdots+\cos^2(\theta_n-\theta_0)}{n}$为集合$S$相对于$\theta_0$的``余弦方差'', 
    求证: 集合$\left\{\dfrac{\pi}{3},\dfrac{2\pi}{3},\pi\right\}$相对于任何常数$\alpha$的``余弦方差''是一个与$\alpha$无关的常数. 
    \begin{jieda}
    等价于证明$\cos^2\left(\dfrac{\pi}{3}-\alpha\right)+\cos^2\left(\dfrac{2\pi}{3}-\alpha\right)+\cos^2(\pi-\alpha)$是与$\alpha$无关的常数. \\
    降次, 展开, 化简, 立得所求为$\dfrac{1}{2}$. 
    \end{jieda}
    
    
    \item 已知$f(\theta)=\sin^2\theta+\sin^2(\theta+\alpha)+\sin^2(\theta+\beta)$, 
    其中$\alpha,\beta$是满足$0<\alpha<\beta<\pi$的常数, 试问: 是否存在这样的$\alpha,\beta$, 使得$f(\theta)$的值是与$\theta$无关的定值? 
    若存在, 求出所有的$\alpha$和$\beta$; 若不存在, 请说明理由. 
    \begin{jieda}
    法一: 
    $\begin{cases}
    f(0)=\sin^2\alpha+\sin^2\beta\\
    f(\frac{\pi}{2})=1+\cos^2\alpha+\cos^2\beta
    \end{cases}$, 于是$f(0)+f(\frac{\pi}{2})=3$, 故$f(0)=f(\frac{\pi}{2})=\frac{3}{2}$. \\
    又$f(-\alpha)=\sin^2\alpha+\sin^2(\beta-\alpha)$, $f(-\beta)=\sin^2\beta+\sin^2(\alpha-\beta)$, 
    故$\sin^2\beta=\sin^2\alpha=\dfrac{3}{4}$. \\
    结合$0<\alpha<\beta<\pi$可得$\alpha=\dfrac{\pi}{3}$, $\beta=\dfrac{2\pi}{3}$. \\
    代入检验可知成立. \\
    法二: $f(\theta)=\frac{3}{2}+\frac{1}{2}((\sin2\alpha+\sin2\beta)\sin2\theta-(1+\cos2\alpha+\cos2\beta)\cos2\theta)$. \\
    于是$\begin{cases}
    \sin2\alpha+\sin2\beta=0\\
    1+\cos2\alpha+\cos2\beta=0
    \end{cases}$. 
    \end{jieda}
\end{enumerate}

\subsection{三角变换的应用(1)}
\zsd 半角公式、万能公式 \\
A组：
\begin{enumerate}
\item $\sin{\dfrac{\alpha}{2}}=$\blkda[3]{$\pm\sqrt{\dfrac{1-\cos{\alpha}}{2}}$}; $\cos{\dfrac{\alpha}{2}}$=\blkda[3]{$\pm\sqrt{\dfrac{1+\cos\alpha}{2}}$};\\
$\tan{\dfrac{\alpha}{2}}$=\blkda[3]{$\pm\sqrt{\dfrac{1-\cos\alpha}{1+\cos\alpha}}$}=\blkda[3]{$\dfrac{1-\cos\alpha}{\sin\alpha}$}=\blkda[3]{$\dfrac{\alpha}{1+\cos{\alpha}}$}.

\item 在$\triangle{ABC}$中, 已知$\cos{A}=-\dfrac{3}{5}$, 则$\sin\dfrac{A}{2}$=\blkda{$\dfrac{2\sqrt{5}}{5}$}. 

\item 已知$\tan\alpha=\dfrac{4}{3}$, 且$\pi<\alpha<\dfrac{3\pi}{2}$, 则$\cos\dfrac{\alpha}{2}$=\blkda{$-\dfrac{\sqrt{5}}{5}$}. 

\item 已知$\sin\theta=-\dfrac{3}{5}$, $3\pi<\theta<\dfrac{7\pi}{2}$, 则$\tan\dfrac{\theta}{2}$=\blkda{$-3$}. 

\item 不用计算器计算: $\tan\ang{15}+\cot\ang{15}=$\blkda{4}.(写过程)
\end{enumerate}
B组：
\begin{enumerate}

\item 已知$\sin\theta=\dfrac{3}{5}$, $\dfrac{5\pi}{2}<\theta<3\pi$, 则$\tan\dfrac{\theta}{2}+2\cos{\dfrac{\theta}{2}}$=\blkda{$3-\dfrac{\sqrt{10}}{5}$}. 
\begin{jieda}
    $cos\theta = -\dfrac{4}{5}$，$\dfrac{\theta}{2} \in (\dfrac{5\pi}{4}, \dfrac{3\pi}{2})$，$tan\dfrac{\theta}{2} = \dfrac{sin\theta}{1+cos\theta} = 3$, $cos\dfrac{\theta}{2} = -\sqrt{\dfrac{1+cos\theta}{2}} = -\dfrac{\sqrt{10}}{10}$，故$\tan\dfrac{\theta}{2}+2\cos{\dfrac{\theta}{2}} = \dfrac{15-\sqrt{10}}{5}$
\end{jieda}
\item 化简: $\dfrac{\sin2x}{1+\cos2x}\cdot\dfrac{\cos{x}}{1+\cos{x}}\cdot\dfrac{\sin{x}}{1-\cos{x}}$=\blkda{$1$}. 

\item 已知$25\sin^2\theta+\sin\theta-24=0$, $\theta$为第二象限角, 则$\cos\dfrac{\theta}{2}$=\blkda{$\pm\dfrac{3}{5}$}. 
\begin{jieda}
虽然$\sin\theta$只有一解, 且$\theta$的终边唯一确定, 但是$\frac{\theta}{2}$的终边有两种情形. 
\end{jieda}

\item 已知$2\sin{x}=1+\cos{x}$, 则$\tan\dfrac{x}{2}$的值为\dotfill(\kh{B})
\xx{$\dfrac{1}{2}$}{$\dfrac{1}{2}$或不存在}{2或$\dfrac{1}{2}$}{不存在}
\begin{jieda}
$\begin{cases}
2\sin{x}=1+\cos{x}\\
\sin^2x+\cos^2x=1
\end{cases}$, 
故$(\cos{x},\sin{x})=(-1,0),(\frac{3}{5},\frac{4}{5})$. \\
\ps 有(由)方程组, 解得. 
\end{jieda}

\item  已知$\cos(\pi-A)=\dfrac{12}{13}$, $A$是第二象限的角, 则$\cot\dfrac{3\pi-A}{2}=$\blkda{5}
\begin{jieda}
$\cot\dfrac{3\pi-A}{2}=\tan\dfrac{A}{2}=\dfrac{1-\cos{A}}{\sin{A}}$，$cosA = -\dfrac{12}{13}, sinA = \dfrac{5}{13}$，故$\cot\dfrac{3\pi-A}{2}=5$
\end{jieda}

\item 若$\dfrac{2\sin\theta+\cos\theta}{\sin\theta-3\cos\theta}=-5$, 则$3\cos2\theta+4\sin2\theta$=\blkda{$\dfrac{7}{5}$}. 
\begin{jieda}
$-5=\dfrac{2\tan\theta+1}{\tan\theta-3}$, 故$\tan\theta=2$. \\
于是$E=3\cdot\dfrac{1-\tan^2\theta}{1+\tan^2\theta}+4\cdot\dfrac{2\tan\theta}{1+\tan^2\theta}=\dfrac{7}{5}$. 
\end{jieda}

\item 已知$\dfrac{\tan\alpha+1}{1-\tan\alpha}=2007$, 则$\dfrac{1+\sin2\alpha}{\cos2\alpha}$的值为\blkda{2007}
\begin{jieda}
$E=\dfrac{(\cos\alpha+\sin\alpha)^2}{\cos^2\alpha-\sin^2\alpha}=\dfrac{\cos\alpha+\sin\alpha}{\cos\alpha-\sin\alpha}=\dfrac{1+\tan\alpha}{1-\tan\alpha}$. 
\end{jieda}

\item 化简下列各式
\begin{enumerate}
    \item $\sin 50^{\circ}\left(1+\sqrt{3} \tan 10^{\circ}\right)$
    \item $\frac{1-\tan ^2\left(\frac{\pi}{4}-a\right)}{1+\tan ^2\left(\frac{\pi}{4}-a\right)}$
    \item $\sin ^2\left(a-\frac{\pi}{6}\right)+\sin ^2\left(a+\frac{\pi}{6}\right)-\sin ^2 a$
    \item $
    \frac{(1+\sin a+\cos a)\left(\sin \frac{a}{2}-\cos \frac{a}{2}\right)}{\sqrt{2+2 \cos a}}(\pi<a<2 \pi) $
\end{enumerate}
\begin{jieda}
\begin{enumerate}
    \item $\sin 50^{\circ}\left(1+\sqrt{3} \tan 10^{\circ}\right)=\sin 50^{\circ} \frac{\cos 10^{\circ}+\sqrt{3} \sin 10^{\circ}}{\cos 10^{\circ}}=\sin 50^{\circ} \frac{2\left(\frac{1}{2} \cos 10^{\circ}+\frac{\sqrt{3}}{2} \sin 10^{\circ}\right)}{\cos 10^{\circ}}$
$
=\sin \left(90^{\circ}-40^{\circ}\right) \frac{2 \sin \left(30^{\circ}+10^{\circ}\right)}{\cos \left(90^{\circ}-80^{\circ}\right)}=\cos 40^{\circ} \frac{2 \sin 40^{\circ}}{\sin 80^{\circ}}=\frac{\sin 80^{\circ}}{\sin 80^{\circ}}=1
$
\item $\frac{1-\tan ^2\left(\frac{\pi}{4}-a\right)}{1+\tan ^2\left(\frac{\pi}{4}-a\right)}=\frac{1-\frac{\sin ^2\left(\frac{\pi}{4}-a\right)}{\cos ^2\left(\frac{\pi}{4}-a\right)}}{1+\frac{\sin ^2\left(\frac{\pi}{4}-a\right)}{\cos ^2\left(\frac{\pi}{4}-a\right)}}=\frac{\cos ^2\left(\frac{\pi}{4}-a\right)-\sin ^2\left(\frac{\pi}{4}-a\right)}{\cos ^2\left(\frac{\pi}{4}-a\right)+\sin ^2\left(\frac{\pi}{4}-a\right)}=\frac{\cos 2\left(\frac{\pi}{4}-a\right)}{1}$
$=\cos \left(\frac{\pi}{2}-2 a\right)=\sin 2 a$


\item $
\begin{aligned}
&  \sin ^2\left(a-\frac{\pi}{6}\right)+\sin ^2\left(a+\frac{\pi}{6}\right)-\sin ^2 a=\frac{1-\cos 2\left(a-\frac{\pi}{6}\right)}{2}+\frac{1-\cos 2\left(a+\frac{\pi}{6}\right)}{2}-\sin ^2 a \\
& =1-\sin ^2 a-\frac{1}{2}\left[\cos \left(2 a-\frac{\pi}{3}\right)+\cos \left(2 a+\frac{\pi}{3}\right)\right]=\cos ^2 a-\frac{1}{2}\left(2 \cos 2 a \cos \frac{\pi}{3}\right)
\end{aligned}$

$
=\cos ^2 a-\frac{1}{2} \cos 2 a=\cos ^2 a-\frac{1}{2}\left(2 \cos ^2 a-1\right)=\frac{1}{2}
$


\item $\dfrac{(1+\sin a+\cos a)\left(\sin \dfrac{a}{2}-\cos \dfrac{a}{2}\right)}{\sqrt{2+2 \cos a}}=\dfrac{\left(1+2 \sin \dfrac{a}{2} \cos \dfrac{a}{2}+2 \cos ^2 \dfrac{a}{2}-1\right)\left(\sin \dfrac{a}{2}-\cos \dfrac{a}{2}\right)}{\sqrt{2+2\left(2 \cos ^2 \dfrac{a}{2}-1\right)}}$

\end{enumerate} 
\end{jieda}

\item 在$\triangle{ABC}$中, $\tan\dfrac{A}{2}+\cot\dfrac{A}{2}=\dfrac{10}{3}$, $\cos{B}=\dfrac{5}{13}$. 
\begin{enumerate}[(1)]
\item 求$\cos(A-B)$的值; 
\item 求$\cos\dfrac{A-B}{2}$的值. 
\end{enumerate}
\begin{jieda}
$\dfrac{10}{3}=\tan\dfrac{A}{2}+\cot\dfrac{A}{2}=\dfrac{2}{\sin{A}}$, 故$\sin{A}=\dfrac{3}{5}$. \\
又$\sin{B}=\dfrac{12}{13}>\sin{A}$, 于是$0<A<B<\dfrac{\pi}{2}$. \\
\begin{enumerate*}[(1)]
\item $\cos(A-B)=\dfrac{56}{65}$;  
\item $\cos\dfrac{A-B}{2}=\dfrac{11\sqrt{130}}{130}$. 
\end{enumerate*}
\end{jieda}

\item 已知$\sin\left(\dfrac{\pi}{4}+2\alpha\right)\cdot\sin\left(\dfrac{\pi}{4}-2\alpha\right)=\dfrac{1}{4}$, $\alpha\in\left(\dfrac{\pi}{4},\dfrac{\pi}{2}\right)$, 
求$2\sin^2\alpha+\tan\alpha-\cot\alpha-1$的值. 
\begin{jieda}
法一: 因为$\sin(x+y)\sin(x-y)=\sin^2x-\sin^2y$, 
故$\sin^2\dfrac{\pi}{4}-\sin^22\alpha=\dfrac{1}{4}$, 即$\sin^22\alpha=\dfrac{1}{4}$. \\
又$\alpha\in\left(\dfrac{\pi}{4},\dfrac{\pi}{2}\right)$, 于是$\sin2\alpha=\dfrac{1}{2}$, 故$2\alpha=\dfrac{5\pi}{6}$. \\
故$E=-\cos2\alpha-2\cot2\alpha=\dfrac{5\sqrt{3}}{2}$. \\
法二: 因为$\dfrac{\pi}{4}+2\alpha+\dfrac{\pi}{4}-2\alpha=\dfrac{\pi}{2}$, \\
故$2\sin\left(\dfrac{\pi}{4}+2\alpha\right)\cdot\sin\left(\dfrac{\pi}{4}-2\alpha\right)=\sin\left(2\left(\dfrac{\pi}{4}+2\alpha\right)\right)=\cos4\alpha$. \\
于是$\cos4\alpha=\dfrac{1}{2}$, 且$4\alpha\in\left(\pi,2\pi\right)$, 故$4\alpha=\dfrac{5\pi}{3}$, 即$\alpha=\dfrac{5\pi}{12}$. 
\end{jieda}

\item 已知$-\dfrac{\pi}{4}<\alpha<\dfrac{\pi}{4}$, $\dfrac{\pi}{4}<\beta<\dfrac{3\pi}{4}$, 
且$\sin\left(\dfrac{3\pi}{4}+\alpha\right)=\dfrac{5}{13}$, $\cos\left(\dfrac{\pi}{4}-\beta\right)=\dfrac{3}{5}$, 
求$\sin\dfrac{\alpha-\beta}{2}$的值. 
\begin{jieda}
法一: $\cos(\alpha-\beta)=\cos\left(\dfrac{3\pi}{4}+\alpha+\dfrac{\pi}{4}-\beta-\pi\right)=\dfrac{16}{65}$. \\
$\dfrac{\pi}{2}<\alpha+\dfrac{3\pi}{4}<\pi$, $-\dfrac{\pi}{2}<\dfrac{\pi}{4}-\beta<0$, 故$\alpha-\beta\in(-\pi,0)$, \\
于是$\sin\dfrac{\alpha-\beta}{2}=-\dfrac{7\sqrt{130}}{130}$. \\
法二: 因为$\dfrac{\pi}{2}<\alpha+\dfrac{3\pi}{4}<\pi$, $-\dfrac{\pi}{2}<\dfrac{\pi}{4}-\beta<0$, \\
故$\alpha+\dfrac{3\pi}{4}=\pi-\arcsin\dfrac{5}{13}$, $\dfrac{\pi}{4}-\beta=-\arccos\dfrac{3}{5}$. \\
于是$\alpha-\beta=\dfrac{\pi}{2}-\arcsin\dfrac{5}{13}-\arccos\dfrac{3}{5}$. 剩下的是计算器. 
\end{jieda}




\item 上海某地区想设计建造一个自然保护区, 如图所示, 在一块矩形绿地$ABCD$中(其中$AB$为30米, $AD$为15米), 过道$EF$将其分为两个主要区域($E,F$分别是$AB,CD$边上动点), 
监测区为以$D$为圆心、$AD$为半径的四分之一圆, 古树区为四边形$BEFC$, 且$EF$与圆弧相切, 记切点为$G$, 且点在$ABCD$内部. \\
\begin{minipage}{\textwidth}
\centering
\includegraphics[scale=1]{./bodyxb/src/main42chap22banjiaogongshi01.png}
\end{minipage}
\begin{enumerate}[(1)]
\item 若$\angle{ADE}=20^{\circ}$, 求$EF$的长(结果精确到0.1); 
\item $E$点在线段$AB$上何处时, 才能使古树区的面积最大, 并求出最大值(结果精确到0.01). 
\end{enumerate}
\begin{jieda}
设$\angle{ADE}=\theta\in(0,\frac{\pi}{4}]$, 由条件可知$\angle{ADE}=\angle{GDE}$, 故$\angle{CDG}=\frac{\pi}{2}-2\theta$. \\
在$\triangle{FDG}$中, $\angle{DGF}=\frac{\pi}{2}$, 故$GF=15\tan(\frac{\pi}{2}-2\theta)=15\cot2\theta$. 同理$GE=15\tan\theta$. 
\begin{enumerate}
\item 由条件可得
\begin{align*}
EF & =EG+GF \\
& =15\tan\theta+15\cot2\theta\\
& =15\left(\tan\theta+\dfrac{1-\tan^2\theta}{2\tan\theta}\right)\\
& =\dfrac{15}{2}\left(\tan\theta+\dfrac{1}{\tan\theta}\right)\\
& \approx 23.3
\end{align*}

\item $AE=GE=15\tan\theta$, $DF=\dfrac{15}{\cos(\frac{\pi}{2}-2\theta)}=\dfrac{15}{\sin2\theta}$. \\
于是$S_{AEFD}=\dfrac{1}{2}\cdot\left(15\tan\theta+\dfrac{15}{\sin2\theta}\right)\cdot15
=\dfrac{225}{4}\left(3\tan\theta+\dfrac{1}{\tan\theta}\right)\geq\dfrac{225\sqrt{3}}{2}$. \\
当$\tan\theta=\frac{1}{\sqrt{3}}$, 即$\theta=\ang{30}$时等号成立, 此时$AE=5\sqrt{3}$, 古树区面积的最大值为$255.14$平方米. 
\end{enumerate}
\end{jieda}

\end{enumerate}

C组：
\begin{enumerate}
\item 设 $\dfrac{1+\sin x}{\cos x}=\dfrac{22}{7}$ 且 $\dfrac{1+\cos x}{\sin x}=\dfrac{m}{n}$ ，其中 $\dfrac{m}{n}$ 为最简分数，则 $m+n=$\blkda{44}．

\begin{jieda}
    \begin{enumerate}
        \item 法一 \\
            $\dfrac{1+\cos x}{\sin x}=\dfrac{m}{n}$，即$tan\dfrac{x}{2} = \dfrac{n}{m}$，
            $\dfrac{1+\sin x}{\cos x}=\dfrac{22}{7}$，即$tan(\dfrac{\pi}{4} - \dfrac{x}{2}) = \dfrac{7}{22}$，因此有$\dfrac{n}{m} = tan(\dfrac{\pi}{4} - (\dfrac{\pi}{4} - \dfrac{x}{2})) = \dfrac{1-tan(\dfrac{\pi}{4} - \dfrac{x}{2})}{1+tan(\dfrac{\pi}{4}-\dfrac{x}{2})} = \dfrac{15}{29}$， 因为$\dfrac{m}{n}$ 为最简分数，所以$m+n=44$.
        \item 法二 \\
            由题设，知 $\sin x=$ $2 \sin \dfrac{x}{2} \cos \dfrac{x}{2}, \cos x=\cos ^2 \frac{x}{2}-\sin ^2 \dfrac{x}{2}$ ，且 $1=$ $\cos ^2 \dfrac{x}{2}+\sin ^2 \dfrac{x}{2}$ ．故 由 $\dfrac{1+\sin x}{\cos x}=\dfrac{22}{7}$ ，得 $\dfrac{1+\sin x}{\cos x}=\dfrac{\cos ^2 \dfrac{x}{2}+\sin ^2 \dfrac{x}{2}+2 \sin \dfrac{x}{2} \cos \dfrac{x}{2}}{\cos ^2 \dfrac{x}{2}-\sin ^2 \dfrac{x}{2}}=\dfrac{\cos\dfrac{x}{2}+\sin\dfrac{x}{2}}{\cos\dfrac{x}{2}-\sin\dfrac{x}{2}}=\dfrac{22}{7}$, 即$15\cos\dfrac{x}{2}=29\sin\dfrac{x}{2}$. 于是$\dfrac{m}{n}=\dfrac{1+\cos x}{\sin x}=\dfrac{1+2\cos ^2\dfrac{x}{2}-1}{2\sin\dfrac{x}{2}\cos\dfrac{x}{2}}=\dfrac{\cos \dfrac{x}{2}}{\sin\dfrac{x}{2}}=\dfrac{29}{15}$.从而$m+n=29+15=44$
    \end{enumerate}

\end{jieda}


\item 若三个角 $x, y, z$ 成等差数列，公差为 $\dfrac{\pi}{3}$ ，则 $\tan x \tan y+\tan y \tan z+\tan z \tan x=$ $\qquad$ ．
\begin{jieda}
根据 $x=y-\frac{\pi}{3}, z=y+\frac{\pi}{3}$ ，则

$$
\tan x=\frac{\tan y-\sqrt{3}}{1+\sqrt{3} \tan y}, \tan z=\frac{\tan y+\sqrt{3}}{1-\sqrt{3} \tan y}
$$


所以

$$
\begin{aligned}
& \tan x \tan y=\frac{\tan ^2 y-\sqrt{3} \tan y}{1+\sqrt{3} \tan y} \\
& \tan y \tan z=\frac{\tan ^2 y+\sqrt{3} \tan y}{1-\sqrt{3} \tan y} \\
& \tan z \tan x=\frac{\tan ^2 y-3}{1-3 \tan ^2 y}
\end{aligned}
$$


则 $\tan x \tan y+\tan y \tan z+\tan z \tan x=\frac{9 \tan ^2 y-3}{1-3 \tan ^2 y}=-3$ ．
\end{jieda}
\item 计算 $\cos \frac{2 \pi}{7} \cos \frac{4 \pi}{7} \cos \frac{6 \pi}{7}$ 的值为 $\qquad$ ．
\begin{jieda}
    记 $S=\cos \frac{2 \pi}{7} \cos \frac{4 \pi}{7} \cos \frac{6 \pi}{7}$ ，则

$$
\begin{aligned}
& S=-\cos \frac{\pi}{7} \cos \frac{2 \pi}{7} \cos \frac{4 \pi}{7} \\
& -S \cdot 8 \sin \frac{\pi}{7}=8 \sin \frac{\pi}{7} \cos \frac{\pi}{7} \cos \frac{2 \pi}{7} \cos \frac{4 \pi}{7} \\
& =4 \sin \frac{2 \pi}{7} \cos \frac{2 \pi}{7} \cos \frac{4 \pi}{7} \\
& =2 \sin \frac{4 \pi}{7} \cos \frac{4 \pi}{7} \\
& =\sin \frac{8 \pi}{7} \\
& =-\sin \frac{\pi}{7} .
\end{aligned}
$$


所以，$S=\frac{1}{8}$ ．
\end{jieda}
\end{enumerate}

\clearpage
\subsection{三角变换的应用(3)}
\zsd 积化和差与和差化积公式

B组:
\begin{enumerate}
\item 已知$\alpha-\beta=\dfrac{\pi}{3}$, 则$\sin\alpha\sin\beta$的最大值为\blkda{$\dfrac{3}{4}$}. 

\item 已知$\cos3\theta=\dfrac{1}{4}$, 则$12\cos\theta\cos\left(\dfrac{\pi}{3}+\theta\right)\cos\left(\dfrac{\pi}{3}-\theta\right)$=\blkda{$\dfrac{3}{4}$}. 
\begin{jieda}
$\cos3\theta=4\cos^3\theta-3\cos\theta$, $\sin3\theta=3\sin\theta-4\sin^3\theta$\\
$\cos3\theta=4\cos\theta\cos\left(\dfrac{\pi}{3}+\theta\right)\cos\left(\dfrac{\pi}{3}-\theta\right)$\\
$\sin3\theta=4\sin\theta\sin\left(\dfrac{\pi}{3}+\theta\right)\sin\left(\dfrac{\pi}{3}-\theta\right)$\\
$\tan\theta=\tan\theta\tan\left(\dfrac{\pi}{3}+\theta\right)\tan\left(\dfrac{\pi}{3}-\theta\right)=\dfrac{3\tan\theta-\tan^3\theta}{1-3\tan^2\theta}$
\end{jieda}


\item 已知$\sin\alpha\cos\beta=\dfrac{1}{2}$, 则$\cos\alpha\cos\beta$的取值范围是\blkda[2]{$\left[-\frac{\sqrt{3}}{2},\frac{\sqrt{3}}{2}\right]$}. 
\begin{jieda}
$\iff 1=\sin(\alpha+\beta)+\sin(\alpha-\beta)$. \\
令$t=\cos\alpha\cos\beta$, 则$2t=\cos(\alpha+\beta)+\cos(\alpha-\beta)$. \\
于是$1^2+(2t)^2=\cdots=2+2\cos2\beta$, 即$4t^2=1+2\cos2\beta\in[-1,3]$, 故$t\in[-\frac{\sqrt{3}}{2},\frac{\sqrt{3}}{2}]$. 
\end{jieda}

\item 已知$\dfrac{1}{\sin(\ang{120}-x)}+\dfrac{1}{\sin(\ang{120}+x)}=\dfrac{4\sqrt{3}}{3}$, 则$\cos{x}$=\blkda{1或$-\dfrac{1}{4}$}. 
\begin{jieda}
\begin{align*}
\LHS & = \dfrac{\sin(\ang{120}+x)+\sin(\ang{120}-x)}{\sin(\ang{120}+x)\sin(\ang{120}-x)}\\
& = \dfrac{\sqrt{3}\cos{x}}{\sin^2\ang{120}-\sin^2x}\\
& = \dfrac{\sqrt{3}\cos{x}}{\frac{3}{4}-\sin^2x}
\end{align*}
于是$3-4\sin^2x=3\cos{x}$, 故$4\cos^2x-3\cos{x}-1=0$, 解得$\cos{x}=1$或$-\dfrac{1}{4}$. 
\end{jieda}

\item 若$\cos(\alpha+\beta)\cos(\alpha-\beta)=\dfrac{1}{3}$, 则$\cos^2\alpha-\sin^2\beta$=\blkda{$\dfrac{1}{3}$}
\begin{jieda}
$\cos(\alpha+\beta)\cos(\alpha-\beta)=\cos^2\alpha-\sin^2\beta$. 
\end{jieda}

\item 若$\alpha,\beta$为锐角, 则下列各式中, 正确的是\dotfill(\kh{C})
\xx{$\cos(\alpha+\beta)>\cos\alpha+\cos\beta$}
{$\cos(\alpha+\beta)>\sin\alpha+\sin\beta$}
{$\cos(\alpha+\beta)<\cos\alpha+\cos\beta$}
{$\cos(\alpha+\beta)<\sin\alpha+\sin\beta$}
\begin{jieda}
令$\alpha=\beta=\frac{\pi}{3}$, 可以否定AB选项; \\
令$\alpha=\beta\to0$, 可以否定D选项; \\
对于C选项, 注意到$\alpha,\beta$为锐角, 则$\cos(\alpha+\beta)<\cos\alpha<\cos\alpha+\cos\beta$. 
\end{jieda}


\item 已知 $\sin \theta \sin \left(\theta+\dfrac{\pi}{4}\right)=\dfrac{3 \sqrt{2}}{5}$ ．
\begin{enumerate*}
    \item 求 $\cos \left(2 \theta+\dfrac{\pi}{4}\right)$ 的值；
    \item 求 $\tan \theta$ 的值．
\end{enumerate*}


\begin{jieda}
    \begin{enumerate}
        \item $\sin{\theta}\sin{(\theta+\dfrac{\pi}{4})}=\dfrac{1}{2}[\cos{\dfrac{\pi}{4}}-\cos{(2\theta+\dfrac{\pi}{4})}]=\dfrac{\sqrt{2}}{4}-\dfrac{1}{2}\cos{(2\theta+\dfrac{\pi}{4}}=\dfrac{3\sqrt{2}}{5}$, 可得$\cos{(2\theta+\dfrac{\pi}{4}})=-\dfrac{7\sqrt{2}}{10}$.
        \item 因为 $\sin \theta \sin \left(\theta+\frac{\pi}{4}\right)=\frac{\sqrt{2}}{2} \sin \theta(\sin \theta+\cos \theta)=\frac{3 \sqrt{2}}{5}$ ，
        所以 $\sin ^2 \theta+\sin \theta \cos \theta=\frac{6}{5}$ ，则 $\frac{\sin ^2 \theta+\sin \theta \cos \theta}{\sin ^2 \theta+\cos ^2 \theta}=\frac{\tan ^2 \theta+\tan \theta}{\tan ^2 \theta+1}=\frac{6}{5}$ ，解得 $\tan \theta=2$ 或 3 ．
    \end{enumerate}
\end{jieda}


\item 已知$\sin{x}+\sin{y}=\sqrt{2}$, $\cos{x}+\cos{y}=\dfrac{2\sqrt{3}}{3}$, 求$\cos(x+y)$及$\tan{x}+\tan{y}$的值. 
\begin{jieda}
两式两边平方再相加可得$\cos(x-y)=\frac{2}{3}$. \\
又$\begin{cases}
2\sin\frac{x+y}{2}\cos\frac{x-y}{2}=\sqrt{2}\\
2\cos\frac{x+y}{2}\cos\frac{x-y}{2}=\frac{2\sqrt{3}}{3}
\end{cases}$, 两式相除, 得$\tan\frac{x+y}{2}=\sqrt{\frac{3}{2}}$. \\
于是由万能公式有$\cos(x+y)=-\frac{1}{5}$, $\sin(x+y)=\frac{2\sqrt{6}}{5}$. 
\begin{align*}
\tan{x}+\tan{y} & = \dfrac{\sin{x}}{\cos{x}}+\dfrac{\sin{y}}{\cos{y}}\\
& = \dfrac{\sin(x+y)}{\cos{x}\cos{y}}\\
& = \dfrac{\sin(x+y)}{\frac{1}{2}(\cos(x+y)+\cos(x-y))}\\
& = \dfrac{12\sqrt{6}}{7}
\end{align*}
\ps 另外, 两式两边平方再相减, 得
\begin{align*}
\frac{4}{3}-2
& = \cos^2x+\cos^2y+2\cos{x}\cos{y}-(\sin^2x+\sin^2y+2\sin{x}\sin{y})\\
& = \cos2x+\cos2y+2\cos(x+y)\\
& = 2\cos(x+y)\cos(x-y)+2\cos(x+y)
\end{align*}
于是在$\cos(x-y)=\frac{2}{3}$的前提下, 亦可得到$\cos(x+y)=-\frac{1}{5}$. 
\end{jieda}

\end{enumerate}


C组：

\begin{enumerate}

\item 设函数$f(x)=\dfrac{\sin3x\sin^3x+\cos3x\cos^3x}{\cos^22x}+\sin2x$, $x\in\left(-\dfrac{\pi}{4},\dfrac{\pi}{4}\right)$的最大值为\blkda{$\sqrt{2}$}
\begin{jieda}
\begin{align*}
f(x)
& = \dfrac{(\sin(3x)\sin x)\sin^2x+(\cos(3x)\cos x)\cos^2x}{\cos^22x}+\sin2x \\
& = \dfrac{(-\dfrac{1}{2}(\cos(4x)-\cos(2x))\sin^2x+(\dfrac{1}{2}(\cos(4x)+\cos(2x))\cos^2x}{\cos^22x}+\sin2x \\
& = \dfrac{cos2x+cos2x \cdot cos4x}{2cos^22x} +\sin2x \\
& = \dfrac{1 + cos4x}{2cos^2x} +\sin2x \\
& = \cos2x+\sin2x\\
& = \sqrt{2}\sin\left(2x+\dfrac{\pi}{4}\right)
\end{align*}
因为$x\in\left(-\dfrac{\pi}{4},\dfrac{\pi}{4}\right)$, 故$2x+\dfrac{\pi}{4}\in\left(-\dfrac{\pi}{4},\dfrac{3\pi}{4}\right)$, 
于是$f(x)\in(-1,\sqrt{2}]$，故最大值为$\sqrt{2}$.
\end{jieda}


\item 设$\alpha,\beta,\gamma$均为锐角, 则等式$\cos^2\alpha+\cos^2\beta+\cos^2\gamma+2\cos\alpha\cos\beta\cos\gamma=1$成立的充分必要条件是$\alpha+\beta+\gamma=\pi$. 

\begin{jieda}
法一: $\iff\dfrac{1+\cos2\alpha}{2}+\dfrac{1+\cos2\beta}{2}+\cos^2\gamma+2\cos\alpha\cos\beta\cos\gamma=1$\\
$\iff\dfrac{\cos2\alpha+\cos2\beta}{2}+\cos^2\gamma+2\cos\alpha\cos\beta\cos\gamma=0$\\
$\iff\cos(\alpha+\beta)\cos(\alpha-\beta)+\cos^2\gamma+(\cos(\alpha+\beta)+\cos(\alpha-\beta))\cos\gamma=0$\\
$\iff(\cos\gamma+\cos(\alpha+\beta))(\cos\gamma+\cos(\alpha-\beta))=0$\\
因为$\alpha,\beta,\gamma$均为锐角, 故$\cos\gamma>0$, $\cos(\alpha-\beta)>0$, \\
于是$\cos\gamma+\cos(\alpha+\beta)=0$, 故$\gamma=\pi-\alpha-\beta$, 即$\alpha+\beta+\gamma=\pi$. \\
法二: $\iff(\cos\alpha+\cos\beta\cos\gamma)^2=\sin^2\beta\sin^2\gamma$. \\
由于$\alpha,\beta,\gamma$均为锐角, 故$\cos\alpha+\cos\beta\cos\gamma=\sin\beta\sin\gamma$, 即$\cos\alpha+\cos(\beta+\gamma)=0$. \\
法三: 注意到$\alpha,\beta,\gamma$均为锐角, 则$\cos\gamma=-\cos\alpha\cos\beta+\sqrt{(1-\cos^2\alpha)(1-\cos^2\beta)}$. \\
于是$\cos\gamma=-\cos\alpha\cos\beta+\sin\alpha\sin\beta=-\cos(\alpha+\beta)$. 
\end{jieda}


\item 若 $\tan 4 x=\frac{\sqrt{3}}{3}$ ，则 $\frac{\sin 4 x}{\cos 8 x \cos 4 x}+\frac{\sin 2 x}{\cos 4 x \cos 2 x}+\frac{\sin x}{\cos 2 x \cos x}+\frac{\sin x}{\cos x}=\blkda{$\sqrt{3}$}$.

\begin{jieda}
   根据题意，有 $\frac{\sin 4 x}{\cos 8 x \cos 4 x}+\frac{\sin 2 x}{\cos 4 x \cos 2 x}+\frac{\sin x}{\cos 2 x \cos x}+\frac{\sin x}{\cos x}=$ $(\tan 8 x-\tan 4 x)+(\tan 4 x-\tan 2 x)+(\tan 2 x-\tan x)+\tan x=\tan 8 x=\sqrt{3}$.
\end{jieda}



\item 设 $f(x)=\frac{1}{x+1}+\frac{2}{x^2+1}$ ，则 $f\left(\tan 1^{\circ}\right)+f\left(\tan 2^{\circ}\right)+\cdots+f\left(\tan 89^{\circ}\right)=$\blkda{$\dfrac{267}{2}$}．
\begin{jieda}
例序相加法

$$
\begin{aligned}
& f(\tan \theta)+f\left[\tan \left(90^{\circ}-\theta\right)\right] \\
& =\frac{1}{\tan \theta+1}+\frac{2}{\tan ^2 \theta+1}+\frac{1}{\cot \theta+1}+\frac{2}{\cot ^2 \theta+1} \\
& =\frac{1}{\tan \theta+1}+\frac{2}{\tan ^2 \theta+1}+\frac{\tan \theta}{1+\tan \theta}+\frac{2 \tan ^2 \theta}{1+\tan ^2 \theta} \\
& =1+2=3
\end{aligned}
$$

记 $S=f\left(\tan 1^{\circ}\right)+f\left(\tan 2^{\circ}\right)+\cdots+f\left(\tan 89^{\circ}\right)$ ，
则 $S=f\left(\tan 89^{\circ}\right)+f\left(\tan 88^{\circ}\right)+\cdots+f\left(\tan 1^{\circ}\right)$ ．
两式相加可得： $2 S=3 \times 89$ ，可得原式的值为 $\frac{267}{2}$ .
\end{jieda}

\item 若 $\tan x \tan y=2, \sin x \sin y=\frac{1}{3}$ ，则 $x-y=$ \blkda[4]{$2 k \pi \pm \frac{\pi}{3}, k \in Z$}．
\begin{jieda}
    由题意可得
$$
\tan x \tan y=\frac{\sin x \sin y}{\cos x \cos y}=\frac{1}{3 \cos x \cos y}=2,
$$
解得 $\cos x \cos y=\frac{1}{6}$ ，故
$$
\cos (x-y)=\cos x \cos y+\sin x \sin y=\frac{1}{6}+\frac{1}{3}=\frac{1}{2}
$$
故 $x-y=2 k \pi \pm \frac{\pi}{3}, k \in Z$.
\end{jieda}

\end{enumerate}

\section{解三角形}

\subsection{解三角形(1)}
\zsd 正弦定理\\
A组：

正弦定理
\begin{table}[h]
    \centering
    \begin{tabular}{|c|c|}
        \hline 
        分类 & 内容\\
       \hline
       定理  & \hphantom{11111111} \blkda{$\dfrac{a}{\sin{A}}=\dfrac{b}{\sin{B}}=\dfrac{c}{\sin{C}}=2R $} \hphantom{11111111}（R是$\triangle{ABC}$外接圆的半径） \\
       \hline
       变形公式  & 
       \makecell{\ding{172}$a=2R\sin{A}$,$b=2R\sin{B},c=2R\sin{C}$ \\
       \ding{173}$\sin{A}:\sin{B}:\sin{C}=a:b:c$ \\
        \ding{174}$\sin{A}=\dfrac{a}{2R},\sin{B}=\dfrac{b}{2R},\sin{C}=\dfrac{c}{2R}$}
       
      \\
       \hline
       解决的问题&

       \makecell{\ding{172}已知两角和任意一边，求其他两边和另一角 \\ \ding{173}已知两边和其中一边的对角，求另一边的对角}  
       
       \\
       \hline
    \end{tabular}
\end{table}

\begin{enumerate}

\item 在 $\triangle{ABC}$中，若 $b=3, c=2$, $\sin C=\dfrac{\sqrt{2}}{3}$，则$B$=\blkda[4]{$\dfrac{\pi}{4}$或$\dfrac{3\pi}{4}$}
\begin{jieda}
    利用正弦定理可解得，但要主要最后解的个数的判断.
\end{jieda}

\item 在$\triangle{ABC}$中，$a=6,B=\ang{30},C=\ang{120}$，则$\triangle{ABC}$的面积是\blkda[3]{$9\sqrt{3}$}
\begin{jieda}
    利用面积公式可解.
\end{jieda}


\item  在$\triangle{ABC}$中, $a=5$, $B=\ang{105}$, $C=\ang{15}$, 则此三角形的最大边的长为\blkda[3]{$\dfrac{5\sqrt{6}+15\sqrt{2}}{6}$}.
\begin{jieda}
$A=\ang{60}$, $b=\dfrac{a}{\sin{A}}\sin{B}$. 
\end{jieda}

\item 在 $\triangle A B C$ 中，$A=60^{\circ}, a=7, b=8$ ，则三角形\dotfill(\kh{B})
\xx{有一解}{有两解}{无解}{不确定}


\end{enumerate}


B组：
\begin{enumerate}

\item 一船自西向东匀速航行, 上午10时到达一座灯塔$P$的南偏西$\ang{75}$距塔68海里的$M$处, 下午2时到达这座灯塔的东南方向的$N$处, 
则这只船的航行速度为\blkda{$\dfrac{17\sqrt{6}}{2}$} 海里/小时. 
\begin{jieda}
$\dfrac{68}{\sin\ang{45}}=\dfrac{MN}{\sin\ang{120}}$, 故$MN=34\sqrt{6}$, $v=\dfrac{MN}{4}$. 
\end{jieda}

\item 在$\triangle{ABC}$中, 已知$\cos{A}=\dfrac{5}{13}$, $\sin{B}=\dfrac{3}{5}$, 则$\cos{C}$的值为\blkda{$\dfrac{16}{65}$}. 
\begin{jieda}
$\sin{A}=\dfrac{12}{13}$, 则$\sin{A}>\sin{B}$, 故$A>B$. \\
又$A$为锐角, 故$B$亦为锐角, 于是$\cos{B}=\dfrac{4}{5}$. 
\end{jieda}

\item 在$\triangle{ABC}$中, 若$b=2a\sin{B}$, 则这个三角形中角$A$的值是\blkda{$\dfrac{\pi}{6}$或$\dfrac{5\pi}{6}$}. 
\begin{jieda}
$\iff \sin{B}=2\sin{A}\sin{B}\iff \sin{A}=\frac{1}{2}$. 
\end{jieda}

\item 在$\triangle{ABC}$中，若$a=2,b=2\sqrt{3},A=\ang{30}$,则C的值为\blkda[3]{$\dfrac{\pi}{6}$或$\dfrac{\pi}{2}$}
\begin{jieda}
    $\dfrac{a}{sinA} = \dfrac{b}{sinB}$，因此$sinB = \dfrac{\sqrt{3}}{2}$，故$B = \dfrac{\pi}{3}$或$\dfrac{2\pi}{3}$，因此$C = $$\dfrac{\pi}{6}$或$\dfrac{\pi}{2}$
\end{jieda}


\item 在$\triangle{ABC}$中, 若$A=2B$, 则$a$等于\dotfill(\kh{D})
\xx{$2b\sin{A}$}{$2b\cos{A}$}{$2b\sin{B}$}{$2b\cos{B}$}
\begin{jieda}
$\dfrac{a}{\sin{A}}=\dfrac{b}{\sin{B}}$, 则$\dfrac{a}{\sin2B}=\dfrac{b}{\sin{B}}$, 故$a=2b\cos{B}$. 
\end{jieda}

\item 在$\triangle{ABC}$中, 若$\tan{A}=\dfrac{1}{3}$, $C=\ang{105}$, $BC=1$, 则$AB$=\blkda{$\dfrac{\sqrt{10}}{2}$}. 
\begin{jieda}
$\dfrac{BC}{\sin{A}}=\dfrac{AB}{\sin{C}}$. 
\end{jieda}


\item 某林场为了及时发现火情，设立了两个观测点 $A$ 和 $B$ ．某日两个观测点的林场人员都观测到 $C$ 处出现火情．在 $A$ 处观测到火情发生在北偏西 $40^{\circ}$ 方向，而在 $B$ 处观测到火情在北偏西 $60^{\circ}$ 方向．已知 $B$ 在 $A$ 的正东方向 10 km 处，那么火场 $C$与 $A$ 距离约为 \blkda{$14.6$} km．（结果精确到 0.1 km ）
\begin{minipage}{\textwidth}
\centering
\includegraphics[scale=0.8]{./bodyxb/src/afslasdjflsgag.jpg}
\end{minipage}
\begin{jieda}
    在 $\triangle{ABC}$ 中，由正弦定理可得 $\frac{A B}{\sin \angle A C B}=\frac{A C}{\sin \angle A B C}$ ，
即 $\frac{10}{\sin 20^{\circ}}=\frac{A C}{\sin 30^{\circ}}$ ，
所以 $A C=\frac{10 \sin 30^{\circ}}{\sin 20^{\circ}} \approx 14.6 \mathrm{~km}$ ，

\end{jieda}




\item 某货轮在 $A$ 处看灯塔 $S$ 在北偏东 $30^{\circ}$ 方向，它以每小时 18 海里的速度向正北方向航行，经过 40 分钟航行到 $B$ 处，看灯塔 $S$ 在北偏东 $75^{\circ}$ 方向，则此时货轮到灯塔 $S$ 的距离为 \blkda{$6\sqrt{2}$}．


\item 在$\triangle{ABC}$中, 若$2\cos{B}\sin{A}=\sin{C}$, 则$\triangle{ABC}$的形状一定是\dotfill(\kh{C})
\xx{等腰直角三角形}{直角三角形}{等腰三角形}{等边三角形}
\begin{jieda}
$\iff 2\cos{B}a=c\iff 2a\cos{B}=a\cos{B}+b\cos{A}\iff a\cos{B}=b\cos{A}\iff \sin{A}\cos{B}=\sin{B}\cos{A}\iff\sin(A-B)=0$. 
\end{jieda}

\item 在 $\triangle A B C$ 中，角 $A, B,C$ 的对边分别是 $a, b, c, A=30^{\circ}, a=3$，若这个三角形有两解，则 $b$ 的取值范围是 \blkda{$3<b<6$}.

\item 在 $\triangle A B C$ 中，$B C=6, A C=8, \angle A=40^{\circ}$ ，则 $\angle B$ 的解的个数是 \blkda{$2$} 个．

\item 在 $\triangle A B C$ 中，由已知条件解三角形，其中有两解的是\dotfill(\kh{C})
\xx{$b=20, A=45^{\circ}, C=80^{\circ}$}{$a=30, c=28, B=60^{\circ}$}{$a=14, b=16, A=45^{\circ}$}{$a=12, c=15, A=120^{\circ}$}

\item 在 $\triangle A B C$ 中，$A C=8, A=45^{\circ}$ ，要使 $\triangle A B C$ 被唯一确定，那么 $B C$ 的取值范围是 \blkda[3]{$\{4 \sqrt{2}\} \cup[8,+\infty)$}.

\item 在 $\triangle A B C$ 中，角 $A, B, C$ 所对的边分别为 $a, b, c$ ，若 $a=\sqrt{3}$ ，且 $c-2 b+2 \sqrt{3} \cos C=0$ ，则该三角形外接圆的半径为 \blkda{$1$}．

\begin{jieda}
$\because a=\sqrt{3}, \therefore c-2 b+2 a \cos C=0, \therefore \sin C-2 \sin B+2 \sin A \cos C=0$ ．

$\therefore \sin C-2 \sin (A+C)+2 \sin A \cos C=0$

$\therefore \sin C-2 \sin A \cos C-2 \sin C \cos A+2 \sin A \cos C=0$

$\therefore \sin C-2 \sin C \cos A=0, \because \sin C>0, \therefore \cos A=\frac{1}{2}, A \in(0, \pi)$ ，

$\therefore A=\frac{\pi}{3}$ ，设该三角形外接圆的半径为 $r$

由正弦定理得$\frac{a}{\sin A}=\frac{\sqrt{3}}{\frac{\sqrt{3}}{2}}=2=2 r, \therefore r=1$
\end{jieda}


\item 在 $\triangle A B C$ 中，$B C=6, A C=5, D$ 为边 $A B$ 的中点，则 $\triangle A C D$ 的外接圆面积 $S_1$与 $\triangle B C D$ 的外接圆面积 $S_2$ 之比为\blkda{$\frac{25}{36}$}.
\begin{jieda}
    设 $\triangle A C D$ 的外接圆为 $R_1, \triangle B C D$ 的外接圆半径为 $R_2$ ，在 $\triangle A C D$ 由正弦定理得 $\frac{A C}{\sin \angle A D C}=2 R_1$ ，在 $\triangle B C D$ 中由正弦定理得 $\frac{B C}{\sin \angle B D C}=2 R_2$ ，又 $\because \angle A D C+\angle B D C=\pi, \therefore \sin \angle A D C=\sin (\pi-\angle B D C)=\sin \angle B D C$ ，$\therefore \frac{S_1}{S_2}=\frac{\pi R_1^2}{\pi R_2^2}=\left(\frac{5}{6}\right)^2=\frac{25}{36},$故答案为：$\frac{25}{36}$ ．
\end{jieda}

\item 已知在三角形 ABC 中，若边 $c=12, A=45^{\circ}, C=30^{\circ}$ ，求 $a, b$ 和 $B$ ．
\begin{jieda}
    根据正弦定理，得$a=\frac{c \sin A}{\sin C}=\frac{12 \times \sin 45^{\circ}}{\sin 30^{\circ}}=12 \sqrt{2},$

又 $B=180^{\circ}-(A+C)=180^{\circ}-\left(45^{\circ}+30^{\circ}\right)=$ $105^{\circ}$ ．

所以 $b=\frac{c \sin B}{\sin C}=\frac{12 \times \sin 105^{\circ}}{\sin 30^{\circ}}=24 \sin 75^{\circ}$ $=24 \times \frac{\sqrt{6}+\sqrt{2}}{4}=6(\sqrt{6}+\sqrt{2})$.

故 $a=12 \sqrt{2},\quad b=6(\sqrt{6}+\sqrt{2}),\quad B=105^{\circ}$.

\end{jieda}

\item 在 $\triangle ABC$ 中，已知 $a=\sqrt{3}, b=\sqrt{2}, B=45^{\circ}$ ，求 $A, C$ 及 $c$ ．
\begin{jieda}
由正弦定理得 $\sin A=\frac{a \sin B}{b}=\frac{\sqrt{3} \times \frac{\sqrt{2}}{2}}{\sqrt{2}}=\frac{\sqrt{3}}{2}$ ，

又 $a>b, A>B$ ，则 $A=60^{\circ}$ 或 $120^{\circ}$ ．

由 $\sin 75^{\circ}=\sin \left(45^{\circ}+30^{\circ}\right)=\frac{\sqrt{6}+\sqrt{2}}{4}$ ，

$
\sin 15^{\circ}=\sin \left(45^{\circ}-30^{\circ}\right)=\frac{\sqrt{6}-\sqrt{2}}{4}
$

当 $A=60^{\circ}$ 时，$C=75^{\circ}$ 时，

$
c=\frac{a \sin C}{\sin A}=\frac{\sqrt{3}}{\frac{\sqrt{3}}{2}} \times \frac{\sqrt{6}+\sqrt{2}}{4}=\frac{\sqrt{6}+\sqrt{2}}{2} .
$

当 $A=120^{\circ}$ 时，$C=15^{\circ}$ 时，

$
c=\frac{a \sin C}{\sin A}=\frac{\sqrt{3}}{\frac{\sqrt{3}}{2}} \times \frac{\sqrt{6}-\sqrt{2}}{4}=\frac{\sqrt{6}-\sqrt{2}}{2} .
$

\end{jieda}

\item $\triangle A B C$ 中若边 $b=2 \sqrt{3}, B=60^{\circ}, c=2$ ，求 $a$ 和 $A, C$
\begin{jieda}
    $C=\frac{\pi}{6}, A=\frac{\pi}{2}, a=4$
\end{jieda}

\item 在$\triangle A B C$中，已知 $BC=a, AC=b$ ，且 $a, b$ 是方程 $x^2-2 \sqrt{3} x+2=0$ 的两根， $2 \cos (A+B)=1$
\begin{enumerate}
    \item 求 $\angle C$ 的度数；
    \item 求 $\triangle A B C$ 的面积.
\end{enumerate}
\begin{jieda}
\begin{enumerate}
    \item $2 \cos (A+B)=1$ 知 $\cos C=-\frac{1}{2}$ 则 $C=120^{\circ}$                                                                                                                                                                                                                                                             
    \item $S_{\triangle A B C}=\frac{1}{2} a b \sin C=\frac{\sqrt{3}}{2}$
\end{enumerate}
\end{jieda}

\item 在 $\triangle A B C$ 中， $\sin A+\cos A=\frac{\sqrt{2}}{2}, A C=2, A B=3$ ，求 $\tan A$ 的值和 $\triangle A B C$ 的面积．
\begin{jieda}
先解三角方程，求出角 $A$ 的值．

$\because \sin A+\cos A=\sqrt{2} \cos \left(A-45^{\circ}\right)=\frac{\sqrt{2}}{2}, $

$\therefore \cos \left(A-45^{\circ}\right)=\frac{1}{2}.$

又$0^{\circ}<A<180^{\circ}, \therefore A-45^{\circ}=60^{\circ}, A=105^{\circ}. $

$\therefore \tan A=\tan \left(45^{\circ}+60^{\circ}\right)=\frac{1+\sqrt{3}}{1-\sqrt{3}}=-2-\sqrt{3} . $

$\sin A=\sin 105^{\circ}=\sin \left(45^{\circ}+60^{\circ}\right)=\sin 45^{\circ} \cos 60^{\circ}+\cos 45^{\circ} \sin 60^{\circ}=\frac{\sqrt{2}+\sqrt{6}}{4}$

$\therefore S_{\triangle A B C}=\frac{1}{2} A C \cdot A B \sin A=\frac{1}{2} \times 2 \times 3 \times \frac{\sqrt{2}+\sqrt{6}}{4}=\frac{3}{4}(\sqrt{2}+\sqrt{6}) .$

\end{jieda}

\item 在 $\triangle A B C$ 中，记角 $A, ~ B, ~ C$ 所对的边分别为 $a, b, c$ ，已知 $2 \sin ^2 \frac{A+B}{2}=1+\cos 2 C$ ．
\begin{enumerate}
    \item 求 $C$ 的大小;
    \item 若 $c^2=2 b^2-2 a^2$ ，求 $\cos 2 A-\cos 2 B$ 的值.
\end{enumerate}
\begin{jieda}
    \begin{enumerate}
        \item 由 $2 \sin ^2 \frac{A+B}{2}=1+\cos 2 C$ ，可得 $1-2 \sin ^2 \frac{A+B}{2}+\cos 2 C=0$ ，则 $\cos (A+B)+\cos 2 C=0$ ，即 $2 \cos ^2 C-\cos C-1=0$ ，所以 $\cos C=-\frac{1}{2}$ 或 $\cos C=1$ ．因为 $0<C<\pi$ ，所以 $\cos C=-\frac{1}{2}$ ，所以 $C=\frac{2 \pi}{3}$ ．
        \item 由 $c^2=2 b^2-2 a^2$ 及正弦定理可得 $\sin ^2 C=2 \sin ^2 B-2 \sin ^2 A$ ，即 $\frac{3}{4}=1-\cos 2 B-(1-\cos 2 A)$．故 $\cos 2 A-\cos 2 B=\frac{3}{4}$ ．
    \end{enumerate}
\end{jieda}

\end{enumerate}

\clearpage
\subsection{解三角形(2)}
\zsd 正余弦定理 \\
A组：\\

余弦定理 \\
\begin{table}[h]
    \centering
    \begin{tabular}{|c|c|}
       \hline
       公式表达  & $a^2 = $\blkda[5]{$b^2+c^2-2bccosA$} \\
       \hline
       推论  & $cosA = $\blkda[5]{$\dfrac{b^2+c^2-a^2}{2ab}$} \\
       \hline
    \end{tabular}
\end{table}
\begin{enumerate}
    \item 在$\triangle ABC$中，已知$b = 40, c = 32, A = 60^{\circ}$，则$a$ = \blkda{$8\sqrt{21}$}
    \begin{jieda}
        $a^2 = b^2+c^2-2bccosA = 1600 + 1024 - 1280 = 1344$，故$a = 8\sqrt{21}$
    \end{jieda}
    \item 在$\triangle{ABC}$内，已知三边之比为$2:3:4$，则该三角形的最大角的余弦值为\blkda{$-\dfrac{1}{4}$}
    \begin{jieda}
        设$a = 2k, b = 3k, c = 4k (k > 0)$，则根据大角对大边，最大角应为$C$，故$cosC = \dfrac{a^2 + b^2 - c^2}{2ab} = -\dfrac{1}{4}$
    \end{jieda}
    \item 在$\triangle{ABC}$内，三边$a, b, c$满足$a^2=b^2+c^2+bc$，则$a$的对角$A = $\blkda{$\dfrac{2\pi}{3}$}
    \begin{jieda}
        根据余弦定理有$a^2=b^2+c^2-2bccosA$，故$cosA = -\dfrac{1}{2}$，故$A = \dfrac{2\pi}{3}$
    \end{jieda}
    \item 若$\triangle{ABC}$的三个内角满足$sinA : sinB : sinC = 5: 11: 13$，则$\triangle{ABC}$ \dotfill{（\kh{C}）}
    \xx{一定是锐角三角形}{一定是直角三角形}{一定是钝角三角形}{可能是锐角三角形，也可能是钝角三角形}
    \begin{jieda}
        根据正弦定理有$a: b : c = 5 : 11 : 13$，故$cosC = \dfrac{25 + 121 - 169}{110} < 0$，故$\triangle{ABC}$为钝角三角形
    \end{jieda}
\end{enumerate}

B组：
\begin{enumerate}
\item 在$\triangle{ABC}$中, $(2a-c)\cos{B}=b\cos{C}$, 则$B$=\blkda{$\dfrac{\pi}{3}$}. 
\begin{jieda}
$\iff 2a\cos{B}=c\cos{B}+b\cos{C}=a\iff\cos{B}=\frac{1}{2}$. 
\end{jieda}

\item 某人在$C$点测得某塔在南偏西$\ang{80}$, 塔顶仰角为$\ang{45}$, 此人沿南偏东$\ang{40}$方向前进10米到$D$, 测得塔顶$A$的仰角为$\ang{30}$, 则塔高为\blkda{10}米.
\begin{jieda}
    设塔高为$h$，塔底为$B$则根据仰角可知$|BC| = h$，$|BD| = \sqrt{3}h$，由此在$\triangle BCD$中可以根据余弦定理有：
    $3h^2=h^2+100-2h\cdot10\cdot\cos\ang{120}$. 因此$h = 10$
    \end{jieda}

\item 若满足条件$C=\dfrac{\pi}{3}$, $AB=\sqrt{3}$, $BC=a$的$\triangle{ABC}$有两个, 那么$a$的取值范围是\blkda{$(\sqrt{3},2)$}
\begin{jieda}
法一: 画图. \\
法二: 余弦定理. 设$AC=x$, 则$3=x^2+a^2-2ax\cos\ang{60}$. \\
于是$x^2-ax+a^2-3=0$, 该方程有两个不等正根. \\
法三： $\dfrac{c}{sinC} = \dfrac{a}{sinA}$. 当$A$有两解时，其取值范围是$(\dfrac{\pi}{3}, \dfrac{\pi}{2}) \cup (\dfrac{\pi}{2}, \dfrac{2\pi}{3})$，$sinA \in (\dfrac{\sqrt{3}}{2}, 1)$，故$a \in (\sqrt{3}, 2)$.
\end{jieda}

\item 在$\triangle{ABC}$中, $(a^2+c^2-b^2)\tan{B}=\sqrt{3}ac$, 则角$B$的值为\blkda{$\dfrac{\pi}{3}$或$\dfrac{2\pi}{3}$}

\begin{jieda}
$(a^2+c^2-b^2)\tan{B}=\sqrt{3}ac$即$\sqrt{3}cotB = \dfrac{a^2+c^2-b^2}{ac}$，故根据余弦定理有$\dfrac{\sqrt{3}}{2}cotB = cosB$，故$\sin{B}=\frac{\sqrt{3}}{2}$，因此$B = \dfrac{\pi}{3}$或$\dfrac{2\pi}{3}$
\end{jieda}

\item 在$\triangle{ABC}$中, $b^2+c^2=a^2+\sqrt{2}bc$, 且$a=\sqrt{2}b$, 则$\angle{C}$=\blkda{$\dfrac{7\pi}{12}$}. 

\item 在$\triangle{ABC}$中, $\cos\dfrac{A}{2}=\dfrac{2\sqrt{5}}{5}$, $\triangle{ABC}$的面积为2, 且$b+c=6$, 则$a$=\blkda{$2\sqrt{5}$}. 
\begin{jieda}
    $cos\dfrac{A}{2} = \dfrac{2\sqrt{5}}{5} \Rightarrow sin\dfrac{A}{2} = \dfrac{\sqrt{5}}{5} \Rightarrow sinA = \dfrac{4}{5}$，$cosA = cos^2\dfrac{A}{2} - sin^2\dfrac{A}{2} = \dfrac{3}{5}$ \\
    $S = \dfrac{1}{2}bcsinA = 2$，因此$bc = 5$，故$cosA = \dfrac{b^2+c^2-a^2}{2bc} = \dfrac{(b+c)^2-a^2}{2bc} - 1 = \dfrac{36-a^2}{10} - 1$，故$a = 2\sqrt{5}$
\end{jieda}

\item 在锐角$\triangle{ABC}$中, $a=1$, $b=2$, 则边长$c$的取值范围是\blkda{$(\sqrt{3},\sqrt{5})$}. 
\begin{jieda}
    根据两边之和大于第三边，$c \in (1, 3)$ \\
    根据$cosC = \dfrac{a^2+b^2-c^2}{2ab} > 0$可得$c < \sqrt{5}$ \\
    根据$cosB = \dfrac{a^2 + c^2 - b^2}{2ac} > 0$可得$c > \sqrt{3}$ \\
    综上，$c \in (\sqrt{3}, \sqrt{5})$
\end{jieda}
    
\item 在$\triangle{ABC}$中, ``$c\cos{B}=b\cos{C}$''是``$\triangle{ABC}$是等腰三角形''的\dotfill(\kh{A})
\xx{充分不必要条件}{必要不充分条件}
{充分必要条件}{既不充分也不必要条件}
\begin{jieda}
条件$\iff c\cdot\dfrac{a^2+c^2-b^2}{2ac}=b\cdot\dfrac{a^2+b^2-c^2}{2ab}\iff a^2+c^2-b^2=a^2+b^2-c^2\iff b=c$. \\
等腰$\iff a=b\lor b=c\lor c=a$. 
\end{jieda}

\item 在锐角$\triangle{ABC}$中, $a=2b\sin{A}$.
\begin{enumerate}
\item 求$B$的大小; 
\item 若$a=3\sqrt{3}$, $c=5$, 求$b$. 
\end{enumerate}
\begin{jieda}
\begin{enumerate}
\item $B=\dfrac{\pi}{6}$; 
\item $b=\sqrt{7}$. 
\end{enumerate}
\end{jieda}
\item 在锐角$\triangle{ABC}$中, $\sqrt{3}a=2c\sin{A}$. 
\begin{enumerate}
\item 求角$C$的大小;
\item 若$c=\sqrt{7}$, 且$\triangle{ABC}$的面积为$\dfrac{3\sqrt{3}}{2}$, 求$a+b$的值. 
\end{enumerate}
\begin{jieda}
\begin{enumerate}
\item $\dfrac{\pi}{3}$

\item $a+b=5$
\end{enumerate}
\end{jieda}


\item 已知半径为$R$的圆外接于$\triangle{ABC}$, 且$2R(\sin^2A-\sin^2C)=(\sqrt{3}a-b)\sin{B}$. 
\begin{enumerate}
\item 求角$C$; 
\item 求$\triangle{ABC}$面积的最大值. 
\end{enumerate}
\begin{jieda}
\begin{enumerate}
\item $\iff a^2-c^2=(\sqrt{3}a-b)b\iff a^2+b^2-c^2=\sqrt{3}ab$, 于是$2ab\cos{C}=\sqrt{3}ab$. \\
于是$\cos{C}=\dfrac{\sqrt{3}}{2}$, 故$C=\dfrac{\pi}{6}$. 

\item 因为面积必须是长度的平方, 条件中又没有边长, 故$S$用$R$表示. \\
$S=2R^2\sin{A}\sin{B}\sin{C}=R^2\sin{C}(\cos(A-B)-\cos(A+B))=\frac{1}{2}R^2(\cos(A-B)+\frac{\sqrt{3}}{2})$, \\
故$S_{\max}=\dfrac{2+\sqrt{3}}{4}R^2(A=B=\frac{5\pi}{12})$
\end{enumerate}
\end{jieda}


\item 在$\triangle{ABC}$中, $a,b,c$分别为内角$A,B,C$的对边，且$2asinA=(2b+c)sinB+(2c+b)sinC$.
\begin{enumerate}
    \item 求A的大小
    \item 若$sinB + sinC = 1$，判断$\triangle{ABC}$的形状
\end{enumerate}
\begin{jieda}
    \begin{enumerate}
        \item 根据正弦定理有$2a^2=(2b+c)b+(2c+b)c$，即$a^2 = b^2+c^2+bc$，故$cosA = -\dfrac{1}{2}$，即$A = \dfrac{2\pi}{3}$
        \item $sinB + sinC = 1$即$sinB + sin(\dfrac{\pi}{3} - B) = 1$，解得$B = \dfrac{\pi}{6}$，故$B = C = \dfrac{\pi}{6}$，因此$\triangle{ABC}$为顶角$\dfrac{2\pi}{3}$的等腰三角形.
    \end{enumerate}
\end{jieda}

\item 在$\triangle{ABC}$中, 若$(a+b+c)(a-b+c)=3ac$, 且$\tan{A}+\tan{C}=3+\sqrt{3}$, $AB$边上的高为$4\sqrt{3}$, 求角$A$、$B$、$C$的大小与边$a$、$b$、$c$的长.
\begin{jieda}
$a^2+c^2-b^2=ac$, 故$\cos{B}=\frac{1}{2}$, 即$B=\frac{\pi}{3}$. \\
于是$A+C=\frac{2\pi}{3}$, 又$\tan{A}+\tan{C}=3+\sqrt{3}$, \\
故$\tan{A}=1$或$2+\sqrt{3}$. \\
\ps 可以$\tan{B}=-\tan(A+C)$, 亦可$\tan{A}+\tan{C}=\dfrac{\sin(A+C)}{\cos{A}\cos{C}}$. \\
$(A,B,C,a,b,c)=(\frac{\pi}{4},\frac{\pi}{3},\frac{5\pi}{12},8,4\sqrt{6},4+4\sqrt{3}), (\frac{5\pi}{12},\frac{\pi}{3},\frac{\pi}{4},8,12\sqrt{2}-4\sqrt{6},8\sqrt{3}-8)$. 
\end{jieda}

\end{enumerate}
C组：
\begin{enumerate}
    \item 在$\triangle{ABC}$中, $a,b,c$分别为内角$A,B,C$的对边，且满足$\sqrt{3}tanAtanB - tanA - tanB = \sqrt{3}$
    \begin{enumerate}
        \item 求C的大小
        \item 若$c = 2$，且$\triangle{ABC}$为锐角三角形，求$a^2+b^2$的取值范围
    \end{enumerate}
    \begin{jieda}
        \begin{enumerate}
            \item $\sqrt{3}tanAtanB - tanA - tanB = \sqrt{3}$即$-\sqrt{3}(1-tanAtanB) = -(tanA + tanB)$，即$\sqrt{3} = -tan(A+B) = tanC$，故$C = \dfrac{\pi}{3}$
            \item 根据正弦定理有$\dfrac{a}{sinA} = \dfrac{b}{sinB} = \dfrac{c}{sinC} = \dfrac{2}{\dfrac{\sqrt{3}}{2}} = \dfrac{4\sqrt{3}}{3}$，因此$a^2+b^2 = \dfrac{16}{3}(sin^2A + sin^2B) = \dfrac{16}{3}(sin^2A + sin^2(\dfrac{2\pi}{3}-A)) = \dfrac{16}{3}(1+\dfrac{1}{2}sin(2A - \dfrac{\pi}{6}))(A \in (\dfrac{\pi}{6}, \dfrac{\pi}{2}))$，因此$a^2+b^2 \in \left(\dfrac{20}{3}, 8 \right]$
        \end{enumerate}        
    \end{jieda}
    \item 已知$sin\alpha + sin\beta = sin\gamma, cos\alpha + cos\beta = -cos\gamma$，则$cos^2\alpha + cos^2\beta + cos^2\gamma = $\blkda{$\dfrac{3}{2}$}
    \begin{jieda}
        $\left \{ \begin{array}{l}
             sin^2\alpha + sin^2\beta + 2sin\alpha sin\beta = sin^2\gamma \\
             cos^2\alpha + cos^2\beta + 2cos\alpha cos\beta = cos^2\gamma
        \end{array}\right.$，两式相加有$2 + 2cos(\alpha - \beta) = 1$，即$cos(\alpha - \beta) = -\dfrac{1}{2}$ \\
        $\left \{ \begin{array}{l}
             sin\left(\dfrac{\alpha+\beta}{2}\right)cos\left(\dfrac{\alpha-\beta}{2}\right) = sin\gamma \\
             cos\left(\dfrac{\alpha+\beta}{2}\right)cos\left(\dfrac{\alpha-\beta}{2}\right) = -cos\gamma
        \end{array}\right.$两式相除有$tan\left(\dfrac{\alpha+\beta}{2}\right) = -tan\gamma$，即$cos(\alpha+\beta) = cos2\gamma$ \\
        $cos^2\alpha + cos^2\beta + cos^2\gamma = \dfrac{3}{2} + \dfrac{1}{2}\left(cos2\alpha+cos2\beta+cos2\gamma\right) = \dfrac{3}{2} + \dfrac{1}{2}(2cos(\alpha+\beta) \cdot cos(\alpha-\beta)+cos2\gamma) = \dfrac{3}{2}$
    \end{jieda}
    \item $cos\dfrac{\pi}{11} - cos\dfrac{2\pi}{11} + cos\dfrac{3\pi}{11}-cos\dfrac{4\pi}{11}+cos\dfrac{5\pi}{11} = $\blkda{$\dfrac{1}{2}$}
    \begin{jieda}
        $cos\dfrac{\pi}{11} - cos\dfrac{2\pi}{11} + cos\dfrac{3\pi}{11}-cos\dfrac{4\pi}{11}+cos\dfrac{5\pi}{11} \\ 
        = cos\dfrac{\pi}{11} + cos\dfrac{3\pi}{11} + cos\dfrac{5\pi}{11} + cos\dfrac{7\pi}{11} + cos\dfrac{9\pi}{11} \\
        = \dfrac{sin(\dfrac{\pi}{11}) \cdot (cos\dfrac{\pi}{11} + cos\dfrac{3\pi}{11} + cos\dfrac{5\pi}{11} + cos\dfrac{7\pi}{11} + cos\dfrac{9\pi}{11})}{sin(\dfrac{\pi}{11})} \\
        = \dfrac{(sin(\frac{2\pi}{11} + sin(\frac{0\pi}{11}) + sin(\frac{4\pi}{11} + sin(\frac{-2\pi}{11}) + sin(\frac{6\pi}{11}+sin(\frac{-4\pi}{11}) + sin(\frac{8\pi}{11}+sin(\frac{-6\pi}{11}) + sin(\frac{10\pi}{11}+sin(\frac{-8\pi}{11}))}{2sin\frac{\pi}{11}} \\
        = \dfrac{1}{2}$
    \end{jieda}
\end{enumerate}

\clearpage
\subsection{解三角形(3)}
\zsd 正余弦定理综合应用 \\
\begin{enumerate}
    \item 在$\triangle{ABC}$中，满足$\dfrac{a}{b} = \dfrac{cosA}{cosB}$的三角形是\blkda{等腰}三角形
    \item 在$\triangle{ABC}$中，若$sin^A + sin^2B < sin^2C$，则$\triangle{ABC}$的形状是\blkda{钝角}三角形
    \begin{jieda}
        $sin^2A + sin^2B < sin^2C$即$a^2+b^2 < c^2$，故$cosC < 0$.
    \end{jieda}
    \item 给出四个命题：
    \begin{dingautolist}{172}
        \item 若$sinA \cdot sinB \cdot sinC > 0$，则$\triangle{ABC}$一定是锐角三角形
        \item 若$cosA \cdot cosB \cdot cosC < 0$，则$\triangle{ABC}$一定是钝角三角形
        \item 若$cotA \cdot cotB > 1$，则$\triangle{ABC}$一定是钝角三角形
        \item 若$\dfrac{a}{b} = sinA$，则$\triangle{ABC}$一定是直角三角形
    \end{dingautolist}
    其中命题正确的有\blkda{\ding{173}\ding{174}\ding{175}}
    \begin{jieda}
        三角形内角正弦值均为证，故\ding{172}错，三角形内角余弦值中至多有一个为负，故\ding{173}正确，$cotA \cdot cotB > 1$即$tanA \cdot tanB < 1$，若$A, B$均为锐角，则$A < \dfrac{\pi}{4}, B < \dfrac{\pi}{4}$，因此$C > \dfrac{\pi}{2}$，故\ding{174}正确，$\dfrac{a}{b} = \dfrac{sinA}{sinB} = sinA$，即$sinB = 1$，故\ding{175}正确.
    \end{jieda}
    \item 在$\triangle{ABC}$中, 若其面积$S=\dfrac{a^2+b^2-c^2}{4\sqrt 3}$, 则$\angle{C}$=\blkda{$\dfrac{\pi}{6}$}. 
    \begin{jieda}
    $\iff \dfrac{1}{2}ab\sin{C}=\dfrac{2ab\cos{C}}{4\sqrt{3}}$. 
    \end{jieda}
    
    \item 在$\triangle{ABC}$中, $AB=2$, $AC=3$, $A=\dfrac{\pi}{3}$, 则$\triangle{ABC}$外接圆的半径为\blkda{$\dfrac{\sqrt{21}}{3}$}. 
    \begin{jieda}
    $2R=\dfrac{a}{\sin{A}}$, $a^2=b^2+c^2-2bc\cos{A}$. 
    \end{jieda}
        
    \item 在锐角$\triangle{ABC}$中, $\dfrac{b}{a}+\dfrac{a}{b}=6\cos{C}$, 则$\dfrac{\tan{C}}{\tan{A}}+\dfrac{\tan{C}}{\tan{B}}$=\blkda{4}. 
    \begin{jieda}
    $\iff a^2+b^2=6ab\cos{C}=3(a^2+b^2-c^2)\iff3c^2=2(a^2+b^2)$, 
    于是
    \begin{align*}
    \dfrac{\tan{C}}{\tan{A}}+\dfrac{\tan{C}}{\tan{B}}
    & =\dfrac{\sin{C}}{\cos{C}}\left(\dfrac{\cos{A}}{\sin{A}}+\dfrac{\cos{B}}{\sin{B}}\right)\\
    & =\dfrac{\sin{C}}{\cos{C}}\cdot\dfrac{\cos{A}\sin{B}+\cos{B}\sin{A}}{\sin{A}\sin{B}}\\
    & =\dfrac{\sin{C}}{\cos{C}}\cdot\dfrac{\sin(A+B)}{\sin{A}\sin{B}}\\
    & =\dfrac{c^2}{ab\cos{C}}\\
    & =\dfrac{2c^2}{a^2+b^2-c^2}\\
    & =4.
    \end{align*}
    \end{jieda}
    \item 在$\triangle{ABC}$中, $a^2-b^2=\sqrt{3}bc$, $\sin{C}=2\sqrt{3}\sin{B}$, 则$A=$\blkda{$\dfrac{\pi}{6}$}
    \begin{jieda}
    $\left \{\begin{array}{l}
         b^2 = a^2-\sqrt{3}bc  \\
         b^2 = a^2 + c^2 - 2bc \cdot cosA
    \end{array} \right.$
    因此$c^2 - 2bc \cdot cosA = -\sqrt{3}bc$，即$c - 2b\cdot cosA = -\sqrt{3}b$ \\
    $\sin{C}=2\sqrt{3}\sin{B}$即$c = 2\sqrt{3}b$，代入可得$cosA = \dfrac{\sqrt{3}}{2}$
    \end{jieda}
    
    \item 已知$\triangle{ABC}$的周长为$\sqrt{2}+1$, 且$\sin{A}+\sin{B}=\sqrt{2}\sin{C}$. 若$\triangle{ABC}$的面积为$\dfrac{1}{6}\sin{C}$, 则$C$=\blkda{$\dfrac{\pi}{3}$}. 
    \begin{jieda}
    根据正弦定理可以得到$a+b=\sqrt{2}c$，故根据题干有$\left \{ \begin{array}{l}
         a+b=\sqrt{2}c \\
         a+b+c = \sqrt{2}+1 \\
         \dfrac{1}{2}absinC = \dfrac{1}{6}sinC
    \end{array}\right.$
    因此$\left \{ \begin{array}{l}
         a+b =\sqrt{2} \\
         ab = \dfrac{1}{3}
    \end{array}\right.$，根据余弦定理有$cosC = \dfrac{a^2+b^2-c^2}{2ab} = -\dfrac{1}{2}$
    \end{jieda}

    \item 已知A、B两地间的距离为$10km$，B、C两地间的距离为$20km$，现测得$\angle ABC = \ang{120}$，则A、C两地间的距离为\blkda{$10\sqrt{7}$}

    \item $\triangle{ABC}$中，$\dfrac{b+a}{a} = \dfrac{sinB}{sinB - sinA}$，且$cos(A-B)+cosC = 1-cos2C$，则$\triangle{ABC}$是\blkda{直角}三角形

    \begin{figure}[h]
        \centering
        \includegraphics[width=0.15\linewidth]{bodyxb//src/jiesanjiao(3)_11.png}
        \caption*{第11题图}
    \end{figure}
    \item 如图所示，$l_1, l_2, l_3$是同一平面内三条平行直线，$l_1$与$l_2$间的距离是1，边长为4的正三角形的三顶点分别在$l_1, l_2, l_3$上，则$l_2$与$l_3$间的距离是\dotfill{(\kh{B})}
    \xx{$2\sqrt{3}$}{$\dfrac{3\sqrt{5}-1}{2}$}{$\dfrac{3\sqrt{15}}{4}$}{$2\sqrt{5}$}

    \begin{figure}[h]
        \centering
        \includegraphics[width=0.25\linewidth]{bodyxb//src/jiesanjiao(3)_12.png}
        \caption*{第12题图}
    \end{figure}

    \item 如图，从气球A上测得正前方的河流的两岸B，C的俯角分别为$\ang{75}$，$\ang{30}$，此时气球的高是$60m$，则河流的宽度BC等于\dotfill{(\kh{C})}
    \xx{$240(\sqrt{3}-1)$}{$180(\sqrt{2}-1)$}{$120(\sqrt{3}-1)$}{$30(\sqrt{3}+1)$}

    \begin{figure}[h]
        \centering
        \includegraphics[width=0.15\linewidth]{bodyxb//src/jiesanjiao(3)_13.png}
        \caption*{第13题图}
    \end{figure}
    
    \item 某同学骑电动车以24 km/h的速度沿正北方向的公路行驶，在点A处测得电视塔S在电动车的北偏东$\ang{30}$方向上，15 min后到点B处，测得电视塔S在电动车的北偏东$\ang{75}$方向上，则点B与电视塔的距离是\blkda{$3\sqrt{2}$}km. 

    \item 在不等边三角形$ABC$中，角$A, B, C$所对的边分别为$a, b, c$，其中$a$为最大边，如果$sin^2(B+C)<sin^2B+sin^2C$，则角$A$的取值范围为\dotfill{(\kh{D})}
    \xx{$\left(0, \dfrac{\pi}{2}\right)$}{$\left(\dfrac{\pi}{4}, \dfrac{\pi}{2}\right)$}{$\left(\dfrac{\pi}{6}, \dfrac{\pi}{3}\right)$}{$\left(\dfrac{\pi}{3}, \dfrac{\pi}{2}\right)$}

    \item 试判断下列三角形形状
    \begin{enumerate}
        \item 在$\triangle{ABC}$中, 已知$(a+b+c)(b+c-a) = 3bc$，$sinA = 2sinBcosC$
        \item 在$\triangle{ABC}$中, 已知$lg(sinA + sinC) = 2lg(sinB) - lg(sinC - sinA)$
    \end{enumerate}
    \begin{jieda}
        \begin{enumerate}
            \item 等边三角形
            \item 直角三角形
        \end{enumerate}
    \end{jieda}

    \item 某港口$O$要将一件重要物品用小艇送到一艘正在航行的轮船上．在小艇出发时，轮船位于港口$O$北偏西$\ang{30}$且与该港口相距20海里的$A$处，并正以30海里/小时的航行速度沿正东方向匀速行驶，假设该小艇沿直线方向以$v$海里/小时的航行速度匀速行驶，经过$t$小时与轮船相遇．
    \begin{enumerate}
        \item 若希望相遇时小艇的航行距离最小，则小艇航行速度的大小应为多少？
        \item 假设小艇的最高航行速度只能达到30海里/小时，试设计航行方案(即确定航行方向和航行速度的大小)，使得小艇能以最短时间与轮船相遇，并说明理由．
    \end{enumerate}
    \begin{jieda}
        \begin{enumerate}
            \item 设小艇与轮船相遇在B点，距离$O$港$S$海里，则在$\triangle{AOB}$中，$A = \dfrac{\pi}{6}$，因此$S = \sqrt{900t^2-600t+400} \geq 10\sqrt{3}$. 当$t = \dfrac{1}{3}$时取等号，此时$v = 30\sqrt{3}$
            \item 在$\triangle{AOB}$中应用余弦定理有$v^2t^2 = 400 + 900t^2 - 600t$，故$v^2 = 900 - \dfrac{600}{t} + \dfrac{400}{t^2} \leq 900$，因此$t \geq \dfrac{2}{3}$，即当$v = 30$时，$t$取最小值. \\
            故可设计方案如下：航行方向为北偏东$\ang{30}$，航行速度为30海里/小时，小艇能以最短时间与轮船相遇.
        \end{enumerate}
    \end{jieda}
    \item 在$\triangle{ABC}$中, 角$A, B, C$所对的边分别为$a, b, c$，且$sinA = \dfrac{\sqrt{5}}{5}$，$sinB = \dfrac{\sqrt{10}}{10}$
    \begin{enumerate}
        \item 求$A+B$的值
        \item 若$a-b = \sqrt{2} - 1$，求$a, b, c$的值.
    \end{enumerate}
    \begin{jieda}
        \begin{enumerate}
            \item $cosB = \dfrac{3\sqrt{10}}{10}$，$cosA = \dfrac{2\sqrt{5}}{5}$或$-\dfrac{2\sqrt{5}}{5}$ \\
            $cos(A+B) = \dfrac{\sqrt{2}}{2}$或$-\dfrac{7\sqrt{2}}{10}$，故$A+B = \dfrac{\pi}{4}$或$arccos(-\dfrac{7\sqrt{2}}{10})$
            \item 根据$\dfrac{a}{sinA} = \dfrac{b}{sinB}$可知$a:b = \sqrt{2}:1$，故$a = \sqrt{2}, b = 1$ \\
            $c^2 = a^2+b^2-2abcosC = 2+1+2\sqrt{2}\cdot cos(A+B) = 5$或$\dfrac{1}{5}$ \\
            因此$c = \sqrt{5}$或$\dfrac{\sqrt{5}}{5}$
        \end{enumerate}
    \end{jieda}
\end{enumerate}
